% !TeX spellcheck = en_US
\documentclass[10pt]{article}

\usepackage{geometry}
\usepackage[no-math]{fontspec}
\usepackage[fontspec-loader]{mathfont}
\usepackage{amsmath}
\usepackage[x11names, table]{xcolor}
\usepackage[colors = all*]{pggfplot}
\usepackage{plantico}
\usepackage{tabularx}
\usepackage{xltabular}
\usepackage{booktabs}
\usepackage{seqsplit}
\usepackage{caption}
\usepackage{newfloat}
\usepackage{xspace}
\usepackage[hidelinks]{hyperref}
\usepackage{polyglossia}
\setdefaultlanguage[variant = american]{english}
\usepackage[noabbrev]{cleveref}

%%% font definitions
\defaultfontfeatures{Ligatures = TeX}
\setsansfont{Arial}
\renewcommand{\familydefault}{\sfdefault}
\mathfont{Arial}
\newfontfamily{\courier}{Courier New}[Scale = 1.25]

%%% font sizes
\newcommand{\fignormal}{\scriptsize} % 7pt
\newcommand{\figsmall}{\fontsize{6}{7}\selectfont} % 6pt
\newcommand{\figtiny}{\tiny} % 5pt
\newcommand{\figlarge}{\footnotesize} % 8pt
\newcommand{\fighuge}{\small} % 9pt

%%% page layout
\geometry{letterpaper, textwidth = 180mm, textheight = 240mm, marginratio = 1:1}
\renewcommand{\textfraction}{0}
\makeatletter
	\setlength{\@fptop}{0pt}
\makeatother
    
%%% header and footer
\pagestyle{empty}

%%% no indentation
\setlength\parindent{0pt}

%%% table layout
\renewcommand{\arraystretch}{1.33}

\setlength{\aboverulesep}{0pt}
\setlength{\belowrulesep}{0pt}

\renewcommand{\tabularxcolumn}[1]{m{#1}}

%%% useful lengths
\newlength{\templength}

\setlength{\columnsep}{4mm}
\newlength{\twocolumnwidth}
\setlength{\twocolumnwidth}{0.5\textwidth}
\addtolength{\twocolumnwidth}{-.5\columnsep}
\newlength{\threecolumnwidth}
\setlength{\threecolumnwidth}{.333\textwidth}
\addtolength{\threecolumnwidth}{-.666\columnsep}
\newlength{\twothirdcolumnwidth}
\setlength{\twothirdcolumnwidth}{2\threecolumnwidth}
\addtolength{\twothirdcolumnwidth}{\columnsep}
\newlength{\fourcolumnwidth}
\setlength{\fourcolumnwidth}{.25\textwidth}
\addtolength{\fourcolumnwidth}{-.75\columnsep}
\newlength{\threequartercolumnwidth}
\setlength{\threequartercolumnwidth}{3\fourcolumnwidth}
\addtolength{\threequartercolumnwidth}{2\columnsep}

%%% tikz setup
\pgfmathsetseed{928}

\input{tikz_config}
%%% set raw data directory
\pggfdatadir{rawData/}


%%% set pggfplot options
\newlength{\groupplotsep}
\setlength{\groupplotsep}{0.75mm}

\pggfset{
	xlabel height = 7.85mm,
	ylabel width = 9.97mm,
	anchor = outer north west,
	thin,
	unbounded coords = jump,
	axis background/.style = {fill = white},
	label style = {node font = \fignormal, text depth = 0pt, align = center},
	title font = \figlarge,
	facet sep = \groupplotsep,
	tickwidth = \groupplotsep,
	tick label style = {node font = \figsmall, text depth = 0pt, /pgf/number format/fixed},
	tick align = outside,
	tick pos = lower,
	every x tick label/.append style = {inner xsep = 0pt},
	every y tick label/.append style = {inner ysep = 0pt},
	every tick/.append style = {black, thin},
	max space between ticks = 20,
	scaled ticks = false,
	legend style = {
		node font = \figsmall,
		font = \strut,
		fill = none,
		draw = none,
		inner ysep = 0pt,
		/tikz/every even column/.append style = {column sep = \groupplotsep},
	},
	legend cell align = right,
	legend plot pos = right,
	mark size = 1,
	boxplot/whisker extend = {\pgfkeysvalueof{/pgfplots/boxplot/box extend} * 0.5},
	boxplot/every median/.style = thick,
	sample size = {node font = \figtiny, text depth = 0pt},
	plot/diagonal/.default = {gray, dashed},
	plot/diagonal*/.default = {gray, dashed},
	pggf colorbar style = {
		yticklabel style = {node font = \figtiny},
		title style = {node font = \figsmall, inner sep = 0pt},
		tick align = outside,
		tick pos = lower,
		tickwidth = \groupplotsep
	},
	origin as zero,
	dendrogram style = {line cap = round, line join = round, semithick},
	dendrogram connector = {gray},
}

\pgfplotsset{
	boxplot/draw/average/.code = {},
	every non boxed x axis/.style = {},
	every non boxed y axis/.style = {},
	xytick/.style = {xtick = {#1}, ytick = {#1}},
	x decimals/.style = {x tick label style = {/pgf/number format/.cd, fixed, zerofill, precision = #1}},
	y decimals/.style = {y tick label style = {/pgf/number format/.cd, fixed, zerofill, precision = #1}},
	z decimals/.style = {z tick label style = {/pgf/number format/.cd, fixed, zerofill, precision = #1}},
	decimals/.style = {x decimals = #1, y decimals = #1, z decimals = #1},
}


%%% variant of /tikz/fill that accepts 'none' as a color (`fill = none` disables filling, `fill* = none` uses color 'none' for filling)
\makeatletter
	\tikzset{
		fill*/.code = {
      \tikz@addoption{\pgfsetfillcolor{#1}}%
      \def\tikz@fillcolor{#1}%
    	\tikz@addmode{\tikz@mode@filltrue}%
		}
	}
\makeatother


%%% set default plot options
\pggfset{
	plot defaults = {bar}{fill opacity = .5},
	plot defaults = {boxplot}{fill opacity = .5},
	plot defaults = {violin}{fill opacity = .5},
	plot defaults = {trendline}{Sienna1},
	plot defaults = {label}{node font = \figsmall},
	plot defaults = {stats}{node font = \figsmall},
	plot defaults = {signif}{node font = \figsmall, inner xsep = 0pt, inner ysep = .2em, level distance = \baselineskip, shift same sample = false},
	plot defaults = {zeroline}{thin},
}


%%% print ticklabels as range
\makeatletter
	\def\pggf@range#1to#2\relax{\pgfmathprintnumber{#1}--\pgfmathprintnumber{#2}}

	\pggf@set{
		xticklabels are range/.code = {%
			\def\pgfplots@user@ticklabel@list@x{%
				\pgfplotslistselectorempty\ticknum\of\pgfplots@xticklabels\to\tick%
				\expandafter\pggf@range\tick\relax%
			}%
		}
	}
	
	\newcommand{\pggf@bins}{\@ifnextchar[\pggf@bins@\pggf@bins@@}
	\def\pggf@bins@[#1,#2]{%
		\pgfmathprintnumber{#1}--\pgfmathprintnumber{#2}%
	}
	\def\pggf@bins@@(#1,#2]{%
		\pgfmathprintnumber{#1}--\pgfmathprintnumber{#2}%
	}
	
	\pggf@set{
		xticklabels are bins/.code = {%
			\def\pgfplots@user@ticklabel@list@x{%
				\pgfplotslistselectorempty\ticknum\of\pgfplots@xticklabels\to\tick%
				\expandafter\pggf@bins\tick\relax%
			}%
		}
	}
\makeatother

%%% useful commands
\renewcommand{\textprime}{\char"2032{}}
\newcommand{\textalpha}{\char"03B1{}}
\renewcommand{\textbeta}{\char"03B2{}}
\newcommand{\textDelta}{\char"0394{}}
\newcommand{\textlambda}{\char"03BB{}}
\renewcommand{\textle}{\char"2264{}}
\renewcommand{\textge}{\char"2265{}}

\newcommand{\enrichment}[1][]{%
	\if\relax#1\relax%
		\def\enrSuffix{}%
	\else%
		\def\enrSuffix{, #1}%
	\fi%
	$\log_2$(enhancer strength\enrSuffix)%
}
\newcommand{\dlratio}[1][]{%
	$\log_2$(NanoLuc/Luc)%
}

%% define shorthands for enhancers, species, insulators, ...
% use: \defname[<color model>]{<shorthand>}{<text form>}{<color>}
% sets both the text form (obtained with \usename{<shorthand>}) and the color
\newcommand{\defname}[4][HTML]{%
	\expandafter\def\csname#2\endcsname{#3\xspace}%
	\definecolor{#2}{#1}{#4}%
}

\newcommand{\usename}[1]{\csname#1\endcsname}

\newcommand{\defalias}[2]{
	\expandafter\def\csname#1\endcsname{\usename{#2}}%
	\colorlet{#1}{#2}
}

% define conditions
\defname[named]{light}{light}{OkabeIto5}
\defname[named]{dark}{dark}{OkabeIto1}
\defname[named]{warm}{warm}{OkabeIto7}
\defname[named]{cold}{cold}{OkabeIto3}

% define enhancers
\defname[named]{none}{\vphantom{A}none}{black}
\defname[named]{35S}{35S}{35Senhancer}
\defname[named]{AB80}{\textit{AB80}}{OrangeRed1}
\defname[named]{Cab-1}{\textit{Cab-1}}{Gold3}
\defname[named]{rbcS-E9}{\textit{rbcS-E9}}{DarkOrchid2}

% define species
\newif\ifspecieslong
\specieslongtrue

\defname[named]{At}{%
	\ifspecieslong%
		\vphantom{gh}\textit{Arabidopsis}%
	\else%
		At%
	\fi%
}{OkabeIto6}
\defalias{Arabidopsis}{At}

\defname[named]{Sl}{%
	\ifspecieslong%
		\vphantom{\textit{Ap}}tomato%
	\else%
		Sl%
	\fi%
}{OkabeIto2}
\defalias{tomato}{Sl}

\defname[named]{Zm}{%
	\ifspecieslong%
		\vphantom{\textit{Ap}}maize%
	\else%
		Zm%
	\fi%
}{OkabeIto8}
\defalias{maize}{Zm}

\defname[named]{Sb}{%
	\ifspecieslong%
		\vphantom{\textit{Ap}}sorghum%
	\else%
		Sb%
	\fi%
}{OkabeIto4}
\defalias{sorghum}{Sb}

\defname[named]{tobacco}{tobacco}{OkabeIto4}

% define TF names
\pgfplotstableread[col sep = tab]{rawData/TF_families.tsv}\TFtable

\pgfplotstableforeachcolumnelement{id}\of\TFtable\as\TFid{%
	\pgfplotstablegetelem{\pgfplotstablerow}{family}\of{\TFtable}
	\expandafter\edef\csname\TFid\endcsname{\pgfplotsretval}
}

% misc names
\expandafter\def\csname STARRseq\endcsname{Plant STARR-seq}
\expandafter\def\csname empty legend\endcsname{\relax}

\defname[named]{low GC}{low GC content}{gray}\colorlet{low GC}{gray!50!white}
\defname[named]{high GC}{high GC content}{gray}\colorlet{high GC}{gray!50!black}

%%% commands to generate random sequences (normal and one-hot encoded)
\newcommand{\randomNucs}[1]{%
	\foreach \x in {1, ..., #1}{%
		\pgfmathparse{int(random(4))}\ifcase\pgfmathresult\or A\or C\or G\or T\fi%
	}%
}

\def\zero{0}
\def\one{\textbf{1}}
\newcommand{\randomOhe}[1]{%
	\foreach \x in {1, ..., #1}{%
		\pgfmathparse{int(random(4))}%
		\ifcase\pgfmathresult\or%
			\one\llap{\raisebox{-.8\baselineskip}{\zero}}\llap{\raisebox{-1.6\baselineskip}{\zero}}\llap{\raisebox{-2.4\baselineskip}{\zero}}%
		\or%
			\zero\llap{\raisebox{-.8\baselineskip}{\one}}\llap{\raisebox{-1.6\baselineskip}{\zero}}\llap{\raisebox{-2.4\baselineskip}{\zero}}%
		\or%
			\zero\llap{\raisebox{-.8\baselineskip}{\zero}}\llap{\raisebox{-1.6\baselineskip}{\one}}\llap{\raisebox{-2.4\baselineskip}{\zero}}%
		\or%
			\zero\llap{\raisebox{-.8\baselineskip}{\zero}}\llap{\raisebox{-1.6\baselineskip}{\zero}}\llap{\raisebox{-2.4\baselineskip}{\one}}%
		\fi%
	}%
}

%%% declare figure types
\newif\ifnpc

% main figures
\newif\ifmain
\newcounter{fig}

\addto\captionsenglish{\renewcommand\figurename{Fig.}}

\newenvironment{fig}{%
	\begin{figure}[p]%
		\stepcounter{fig}%
		\pdfbookmark{\figurename\ \thefig}{figure\thefig}
		\tikzsetnextfilename{figure\thefig}%
		\fignormal%
		\centering%
}{%
	\end{figure}%
	\clearpage%
	\ifnpc%
		\makenextpagecaption%
		\global\npcfalse%
	\fi%
}

% extended data figures
\newif\ifext
\newcounter{efig}

\DeclareFloatingEnvironment[fileext = loef, name = Extended Data \figurename]{extfigure}

\newenvironment{efig}{%
	\begin{extfigure}[p]%
		\stepcounter{efig}%
		\setcounter{subfig}{0}%
		\pdfbookmark{\extfigurename\ \theefig}{ext_data_figure\theefig}
		\tikzsetnextfilename{ext_data_fig\theefig}%
		\fignormal%
		\centering%
}{%
	\end{extfigure}%
	\clearpage%
	\ifnpc%
		\makenextpagecaption%
		\global\npcfalse%
	\fi%
}

% supplementary figures
\newif\ifsupp
\newcounter{sfig}

\DeclareFloatingEnvironment[fileext = losf, name = Supplementary \figurename]{suppfigure}

\newenvironment{sfig}{%
	\begin{suppfigure}[p]%
		\stepcounter{sfig}%
		\setcounter{subfig}{0}%
		\pdfbookmark{\suppfigurename\ \thesfig}{supp_figure\thesfig}
		\tikzsetnextfilename{supp_fig\thesfig}%
		\fignormal%
		\centering%
}{%
	\end{suppfigure}%
	\clearpage%
	\ifnpc%
		\makenextpagecaption%
		\global\npcfalse%
	\fi%
}

% supplementary tables
\newif\ifstab
\newcounter{stab}

\addto\captionsenglish{\renewcommand\tablename{Supplementary Table}}

\newenvironment{stab}{%
	\begin{table}[p]%
		\stepcounter{stab}%
		\pdfbookmark{\tablename\ \thestab}{supp_table\thestab}%
		\fignormal%
}{%
	\end{table}%
	\clearpage%
}

\newenvironment{longstab}{%
	\begingroup%
	\stepcounter{stab}%
	\pdfbookmark{\tablename\ \thestab}{supp_table\thestab}%
	\fignormal%
}{%
	\endgroup%
	\clearpage%
}

%% command to put the caption on the next page
% use instead of a normal caption: \nextpagecaption{<caption text>}
\makeatletter
	\newcommand{\nextpagecaption}[1]{%
		\global\npctrue%
		\xdef\@npctype{\@currenvir}%
		\captionlistentry{#1}%
		\long\gdef\makenextpagecaption{%
				\csname\@npctype\endcsname%
					\ContinuedFloat%
					\caption{#1}%
				\csname end\@npctype\endcsname%
				\clearpage%
		}%
	}
\makeatother

%%% caption format
\DeclareCaptionJustification{nohyphen}{\hyphenpenalty = 10000}
\DeclareCaptionLabelSeparator{bar}{ | }

\captionsetup{
	labelsep = bar,
	justification = nohyphen,
	singlelinecheck = false,
	labelfont = bf,
	font = small,
	figureposition = below,
	tableposition = above
}
\captionsetup[figure]{skip = .5\baselineskip}

\newcommand{\titleend}{ }
\newcommand{\nextentry}{ }
\newcommand{\captiontitle}[2][]{\textbf{#2.}\titleend}

%%% subfigure labels
\newif\ifsubfigupper
\subfigupperfalse

\newcounter{subfig}[figure]

\tikzset{
	subfig label/.style = {anchor = north west, inner sep = 0pt, font = \normalsize\bfseries}
}

\newcommand{\subfiglabel}[2][]{
	\node[anchor = north west, inner sep = 0pt, font = \large\bfseries, #1] at (#2) {\strut\stepcounter{subfig}\ifsubfigupper\Alph{subfig}\else\alph{subfig}\fi};
}

\newcommand{\subfigrefsep}{,}
\newcommand{\subfigrefand}{,}
\newcommand{\subfigrefrange}{\textendash}

\newcommand{\subfigunformatted}[1]{\ifsubfigupper\uppercase{#1}\else\lowercase{#1}\fi}
\newcommand{\plainsubfigref}[1]{\textbf{\subfigunformatted{#1}}}
\newcommand{\subfig}[1]{\plainsubfigref{#1}\subfigrefsep}
\newcommand{\subfigtwo}[2]{\plainsubfigref{#1}\subfigrefand\plainsubfigref{#2}\subfigrefsep}
\newcommand{\subfigrange}[2]{\plainsubfigref{#1}\subfigrefrange\plainsubfigref{#2}\subfigrefsep}
\newcommand{\parensubfig}[2][]{(#1\plainsubfigref{#2})}
\newcommand{\parensubfigtwo}[3][]{(#1\plainsubfigref{#2}\subfigrefand\plainsubfigref{#3})}
\newcommand{\parensubfigrange}[3][]{(#1\plainsubfigref{#2}\subfigrefrange\plainsubfigref{#3})}

%%% cross-reference setup
\newcommand{\crefrangeconjunction}{\subfigrefrange}
\newcommand{\crefpairgroupconjunction}{\subfigrefand}
\newcommand{\crefmiddlegroupconjunction}{, }
\newcommand{\creflastgroupconjunction}{, }

\crefname{figure}{Fig.}{Figs.}
\crefname{extfigure}{Extended Data Fig.}{Extended Data Figs.}
\crefname{suppfigure}{Supplementary Fig.}{Supplementary Figs.}

%%% what to include
\maintrue
\exttrue
\supptrue
\stabtrue

%%%%%%%%%%%%%%%%%%%%%%%%%%%%%%%%%%%%
%%%   ||    preamble end    ||   %%%
%%%--\\//------------------\\//--%%%
%%%   \/   begin document   \/   %%%
%%%%%%%%%%%%%%%%%%%%%%%%%%%%%%%%%%%%

\begin{document}
	\sffamily
	\frenchspacing
	
	%% Main figures start
	\ifmain
		
		\begin{fig}
			\begin{tikzpicture}
	\pgfmathsetseed{928}
	
	%%% scheme of the enhancer screen %%%
	\coordinate (scheme) at (0, 0);

	%% DNA
	\coordinate (DNA) at ($(scheme) + (.85, -.45)$);
	
	\def\scalefactor{0.5}

	\tikzset{
		nucleosome core/.style = {cylinder, rotate = 150, shape aspect = 1, minimum width = \scalefactor * 1cm, minimum height = \scalefactor * 1cm, cylinder uses custom fill, cylinder end fill = gray!33, cylinder body fill = gray!66},
		DNA wrap/.style= {very thick, draw, color = black, line cap = round, line width = \scalefactor * \pgflinewidth}
	}
	
	\begin{scope}[scale = \scalefactor]

		\foreach \i [count = \ii, remember = \i as \li] in {0, 1, 2} { 
			\path node [nucleosome core] (nucleosome\i) at ($(DNA) + (\i, 0)$) {} ;
			\ifnum\ii>1
				\begin{pgfonlayer} {background}
					\draw[DNA wrap] ($(DNA) + (\li, 0)$) to[out = -90, in = 210] (nucleosome\i.65);
				\end{pgfonlayer}
			\fi
			\draw [DNA wrap] (nucleosome\i.65) edge [bend right = 15 / \scalefactor] (nucleosome\i.290);
			\draw [DNA wrap] (nucleosome\i.90) edge [bend right = 15 / \scalefactor] (nucleosome\i.265);
		}
		
		\draw[DNA wrap] ($(nucleosome0) + (-1, 0)$) coordinate (DNA start) to[out = 0, in = 210] (nucleosome0.65);
		\draw[DNA wrap, dotted] (DNA start) -- ++(-0.5, 0);
		
		\begin{pgfonlayer} {background}
			\draw[DNA wrap] ($(nucleosome2) + (.25, 0)$) to[out = -90, in = 180] ++(.5, -.5) coordinate (insulator start);
		\end{pgfonlayer}
		
		\foreach \i [count = \ii, remember = \i as \li] in {5, 6} {
			\path node [nucleosome core] (nucleosome\i) at ($(DNA) + (\i, 0)$) {} ;
			\ifnum\ii>1
				\begin{pgfonlayer} {background}
					\draw[DNA wrap] ($(DNA) + (\li, 0)$) to[out = -90, in = 210] (nucleosome\i.65);
				\end{pgfonlayer}
			\fi
			\draw [DNA wrap] (nucleosome\i.65) edge [bend right = 15 / \scalefactor] (nucleosome\i.290);
			\draw [DNA wrap] (nucleosome\i.90) edge [bend right = 15 / \scalefactor] (nucleosome\i.265);
		}
		
		\draw[DNA wrap] (insulator start) -- ($(nucleosome5) + (-1, -.5)$) coordinate (insulator end) to[out = 0, in = 200] (nucleosome5.65);
		\draw[line width = .15cm, DeepSkyBlue2, -Triangle Cap] (insulator start) -- (insulator end) node[pos = .5, anchor = north, align = center] {putative\\enhancer};
		
		\draw[fill = DeepSkyBlue3] ($(insulator start)!.25!(insulator end)$) ++(0, .5) circle (.45);
		\draw[fill = DeepSkyBlue1] ($(insulator start)!.67!(insulator end)$) ++(0, .4) circle (.35);
		
		\begin{pgfonlayer} {background}
			\draw[DNA wrap] ($(nucleosome6) + (.25, 0)$) to[out = -90, in = 180] ++(.5, -.5) coordinate (enhancer start);
		\end{pgfonlayer}
		
		\foreach \i [count = \ii, remember = \i as \li] in {9, 10} {
			\path node [nucleosome core] (nucleosome\i) at ($(DNA) + (\i, 0)$) {} ;
			\ifnum\ii>1
				\begin{pgfonlayer} {background}
					\draw[DNA wrap] ($(DNA) + (\li, 0)$) to[out = -90, in = 210] (nucleosome\i.65);
				\end{pgfonlayer}
			\fi
			\draw [DNA wrap] (nucleosome\i.65) edge [bend right = 15 / \scalefactor] (nucleosome\i.290);
			\draw [DNA wrap] (nucleosome\i.90) edge [bend right = 15 / \scalefactor] (nucleosome\i.265);
		}
		
		\draw[DNA wrap] (enhancer start) -- ($(nucleosome9) + (-1, -.5)$) coordinate (enhancer end) to[out = 0, in = 200] (nucleosome9.65);
		\draw[line width = .15cm, DodgerBlue2, -Triangle Cap] (enhancer start) -- (enhancer end) node[pos = .5, anchor = north, align = center] {putative\\enhancer};
		
		\draw[fill = DodgerBlue1] ($(enhancer start)!.7!(enhancer end)$) ++(0, .45) circle (.4);
		\draw[fill = DodgerBlue3] ($(enhancer start)!.2!(enhancer end)$) ++(0, .4) circle (.35);
		
		\begin{pgfonlayer} {background}
			\draw[DNA wrap] ($(nucleosome10) + (.25, 0)$) to[out = -90, in = 180] ++(.5, -.5) coordinate (promoter start);
		\end{pgfonlayer}
		
		\foreach \i [count = \ii, remember = \i as \li, evaluate = \i as \ni using int(\i)] in {13.5, 15, 16.5, 18, 19} {
			\path node [nucleosome core] (nucleosome\ni) at ($(DNA) + (\i, 0)$) {} ;
			\ifnum\ii>1
				\begin{pgfonlayer} {background}
					\draw[DNA wrap] ($(DNA) + (\li, 0)$) to[out = -90, in = 210] (nucleosome\ni.65);
				\end{pgfonlayer}
			\fi
			\draw [DNA wrap] (nucleosome\ni.65) edge [bend right = 15 / \scalefactor] (nucleosome\ni.290);
			\draw [DNA wrap] (nucleosome\ni.90) edge [bend right = 15 / \scalefactor] (nucleosome\ni.265);
		}	
		
		\draw[line width = .15cm, 35Spromoter] (promoter start) -- ++(1.25, 0) node[pos = .5, anchor = north] {core promoter} coordinate (promoter end);
		\draw[-{Stealth[round]}, thick] (promoter end) ++(.5\pgflinewidth, 0) |- ++(.6,.4);
		
		\draw[fill = DarkSeaGreen1] ($(promoter start)!.15!(promoter end)$) ++(0, .3) circle (.25);
		\draw[fill = DarkSeaGreen2] ($(promoter start)!.6!(promoter end)$) ++(0, .5) circle (.6 and .45);
		
		\draw[DNA wrap] (promoter end) -- ($(nucleosome13) + (-1.25, -.5)$) coordinate (gene end) to[out = 0, in = 200] (nucleosome13.65);
		
		\begin{pgfonlayer} {background}
			\draw[DNA wrap] ($(nucleosome19) + (.25, 0)$) to[out = -90, in = 180, out looseness = 3] ++(0.5, 0) coordinate (DNA end);
			\draw[DNA wrap, dotted] (DNA end) -- ++(0.5, 0);
		\end{pgfonlayer}
		
	\end{scope}

	%% accessibility
	\coordinate (accessibility) at ($(DNA start) + (0, -1.8)$);
	
	\path (accessibility -| nucleosome2) ++(.25, 0) coordinate (a1) (accessibility -| nucleosome5) ++(-.25, 0) coordinate (a2) (accessibility -| nucleosome6) ++(.25, 0) coordinate (a3) (accessibility -| nucleosome9) ++(-.25, 0) coordinate (a4) (accessibility -| nucleosome10) ++(.25, 0) coordinate (a5) (accessibility -| nucleosome13) ++(-.25, 0) coordinate (a6);
	\draw[fill = gray] (a1) cos ($(a1)!.25!(a2) + (0, .4)$) sin ($(a1)!.5!(a2) + (0, .8)$) cos ($(a1)!.75!(a2) + (0, .4)$) sin (a2) (a3) cos ($(a3)!.25!(a4) + (0, .375)$) sin ($(a3)!.5!(a4) + (0, .75)$) cos ($(a3)!.75!(a4) + (0, .375)$) sin (a4) (a5) cos ($(a5)!.25!(a6) + (0, .5)$) sin ($(a5)!.5!(a6) + (0, 1)$) cos ($(a5)!.75!(a6) + (0, .5)$) sin (a6);
	
	\draw (accessibility) -- (DNA end |- accessibility);
	\node[anchor = south west] at (scheme |- accessibility) {accessibility};
	
	%% species
	\setlength{\templength}{.7cm}
	
	\coordinate[shift = {(-4.3\templength, .8cm -.5\templength)}] (species) at (DNA end |- accessibility);
	
	\planticon[width = \templength]{arabidopsis} at ($(species) + (.5\templength, 0)$);
	\planticon[width = \templength]{tomatoFruit} at ($(species) + (1.6\templength, 0)$);
	\planticon[width = \templength]{maize} at ($(species) + (2.7\templength, 0)$);
	\planticon[width = \templength]{sorghum} at ($(species) + (3.8\templength, 0)$);
		
	%% fragments
	\coordinate[yshift = -.25cm] (fragments) at (accessibility);
	
	\draw[line width = .2cm, -Triangle Cap, DeepSkyBlue2] ($(fragments -| a1)!.5!(fragments -| a2)$) coordinate (frag1) +(-.3, 0) -- +(.3, 0);
	\draw[line width = .2cm, -Triangle Cap, DodgerBlue2] ($(fragments -| a3)!.5!(fragments -| a4)$) coordinate (frag2) +(-.3, 0) -- +(.3, 0);
	
	%% array
	\coordinate (array) at ($(scheme |- fragments) + (.1, -1.8)$);
	
	\draw[thin, fill = gray!50] (array) -- ++(3, 0) coordinate[pos = .5] (array mid) -- ++(0.6, 0.3) coordinate (array right) -- ++(-3, 0) -- cycle;
	\node[anchor = north, align = center, xshift = .3cm, inner xsep = 0pt] (array text) at (array mid) {array synthesis of \textasciitilde \textbf{175,000}\\accessible chromatin regions};
	
	\foreach \i/\col [evaluate = \i as \j using 7 - \i] in {1/DodgerBlue2, 2/DeepSkyBlue2, 3/MediumOrchid2, 4/RoyalBlue2, 5/Chocolate1, 6/Aquamarine2} {
		\foreach \k in {1, 2, 3, 4} {
			\draw[\col!50!white, rounded corners = 0.05cm] (array) ++(\i * 0.5 - 0.1 + rand * 0.06, 0.075 + rand * 0.03) -- ++(0, 0.1) -- ++(-0.05, 0.1) -- ++(0.1, 0.2) -- ++(-0.05, 0.1) -- ++(0, 0.1);
			\draw[\col, rounded corners = 0.05cm] (array) ++(\j * 0.5 + 0.2 + rand * 0.06, 0.225 + rand * 0.03)  -- ++(0, 0.1) -- ++(-0.05, 0.1) -- ++(0.1, 0.2) -- ++(-0.05, 0.1) -- ++(0, 0.1) coordinate (f\i);
		}
	}
	
	\draw[gray, ->, shorten both = .1cm] (frag1) ++(0, -.075) -- (f2);
	\draw[gray, ->, shorten both = .1cm] (frag2) ++(0, -.075) -- (f1);
	
	%% constructs
	\draw[very thick, -Latex] (array right) ++(.2, .2) -- ++(.6, 0) coordinate (arrow 1);
	
	\coordinate[xshift = .2cm] (constructs) at (arrow 1);
	
	\coordinate (construct 1) at ($(constructs) + (.1, .9)$);
	\coordinate (construct 2) at ($(constructs) + (0, .3)$);
	\coordinate (construct 3) at ($(constructs) + (-.1, -.3)$);
	\coordinate (construct 4) at ($(constructs) + (-.2, -.9)$);
	
	% construct 1
	\draw[line width = .2cm, -Triangle Cap, DeepSkyBlue2]  (construct 1) ++(.5, 0) -- ++(.6, 0) coordinate (enh east);
	\draw[line width = .2cm, 35Spromoter] (enh east) ++(.2, 0) -- ++(.3, 0) coordinate (pro east);
	\draw[-{Stealth[round]}, semithick] (pro east) ++(.5\pgflinewidth, 0) |- ++(.35, .3);
	\node[draw = black, thin, anchor = west, minimum width = 1cm, outer xsep = 0pt] (GFP) at ($(pro east) + (.3, 0)$) {GFP};
	\coordinate[xshift = .5cm] (construct 1 end) at (GFP.east);
	
	\begin{pgfonlayer} {background}
		\fill[Orchid1] (GFP.north west) ++(.15, 0) rectangle ($(GFP.south west) + (.05, 0)$);
		\draw[thick, dashed] (construct 1) -- ++(.4, 0) coordinate (c1) (construct 1 end) -- ++(-.4, 0) coordinate (c2);
		\draw[thick] (c1) -- (GFP.west) (GFP.east) -- (c2);
	\end{pgfonlayer}
	
	% construct 2
	\draw[line width = .2cm, Triangle Cap-, DeepSkyBlue2]  (construct 2) ++(.5, 0) -- ++(.6, 0) coordinate (enh east);
	\draw[line width = .2cm, 35Spromoter] (enh east) ++(.2, 0) -- ++(.3, 0) coordinate (pro east);
	\draw[-{Stealth[round]}, semithick] (pro east) ++(.5\pgflinewidth, 0) |- ++(.35, .3);
	\node[draw = black, thin, anchor = west, minimum width = 1cm, outer xsep = 0pt] (GFP) at ($(pro east) + (.3, 0)$) {GFP};
	\coordinate[xshift = .5cm] (construct 2 end) at (GFP.east);
	
	\begin{pgfonlayer} {background}
		\fill[MediumPurple2] (GFP.north west) ++(.15, 0) rectangle ($(GFP.south west) + (.05, 0)$);
		\draw[thick, dashed] (construct 2) -- ++(.4, 0) coordinate (c1) (construct 2 end) -- ++(-.4, 0) coordinate (c2);
		\draw[thick] (c1) -- (GFP.west) (GFP.east) -- (c2);
	\end{pgfonlayer}

	% construct 3
		\draw[line width = .2cm, -Triangle Cap, DodgerBlue2]  (construct 3) ++(.5, 0) -- ++(.6, 0) coordinate (enh east);
		\draw[line width = .2cm, 35Spromoter] (enh east) ++(.2, 0) -- ++(.3, 0) coordinate (pro east);
		\draw[-{Stealth[round]}, semithick] (pro east) ++(.5\pgflinewidth, 0) |- ++(.35, .3);
		\node[draw = black, thin, anchor = west, minimum width = 1cm, outer xsep = 0pt] (GFP) at ($(pro east) + (.3, 0)$) {GFP};
		\coordinate[xshift = .5cm] (construct 3 end) at (GFP.east);
		
		\begin{pgfonlayer} {background}
			\fill[HotPink1] (GFP.north west) ++(.15, 0) rectangle ($(GFP.south west) + (.05, 0)$);
			\draw[thick, dashed] (construct 3) -- ++(.4, 0) coordinate (c1) (construct 3 end) -- ++(-.4, 0) coordinate (c2);
			\draw[thick] (c1) -- (GFP.west) (GFP.east) -- (c2);
		\end{pgfonlayer}
	
	% construct 4
	\draw[line width = .2cm, Triangle Cap-, DodgerBlue2]  (construct 4) ++(.5, 0) -- ++(.6, 0) coordinate (enh east);
	\draw[line width = .2cm, 35Spromoter] (enh east) ++(.2, 0) -- ++(.3, 0) coordinate (pro east);
	\draw[-{Stealth[round]}, semithick] (pro east) ++(.5\pgflinewidth, 0) |- ++(.35, .3);
	\node[draw = black, thin, anchor = west, minimum width = 1cm, outer xsep = 0pt] (GFP) at ($(pro east) + (.3, 0)$) {GFP};
	\coordinate[xshift = .5cm] (construct 4 end) at (GFP.east);
	
	\begin{pgfonlayer} {background}
		\fill[SlateBlue3] (GFP.north west) ++(.15, 0) rectangle ($(GFP.south west) + (.05, 0)$);
		\draw[thick, dashed] (construct 4) -- ++(.4, 0) coordinate (c1) (construct 4 end) -- ++(-.4, 0) coordinate (c2);
		\draw[thick] (c1) -- (GFP.west) (GFP.east) -- (c2);
	\end{pgfonlayer}
	
	%% infiltration
	\draw[very thick, -Latex] (arrow 1 -| construct 2 end) ++(.2, 0) -- ++(.6, 0) coordinate (arrow 2);
	
	\coordinate[xshift = 1cm] (transformation) at (arrow 2);
		
	\coordinate[shift = {(.4, .4)}] (leaf center) at (transformation);
	
	\planticon*[scale = .65]{tobaccoLeaf} at (leaf center);
	
	\begin{scope}[xscale = -1]
		\planticon[width = .7cm]{pipe} at ($(transformation |- leaf center) + (.425, .33)$);
	\end{scope}
	
	\coordinate[shift = {(-.5, -.55)}] (protoplasts) at (transformation);
	
	\planticon[width = .6cm]{protoplast} at ($(protoplasts) + (90:.3)$);
	\planticon[width = .6cm]{protoplast} at ($(protoplasts) + (162:.3)$);
	\planticon[width = .6cm]{protoplast} at ($(protoplasts) + (234:.3)$);
	\planticon[width = .6cm]{protoplast} at ($(protoplasts) + (306:.3)$);
	\planticon[width = .6cm]{protoplast} at ($(protoplasts) + (18:.3)$);
	
	\planticon[width = .7cm, rotate = 15]{maize} at ($(protoplasts) + (.8, -.2)$);
	
	%% conditions
	\coordinate[xshift = .45cm + \columnsep] (conditions) at (scheme -| DNA end);
	
	\setlength{\templength}{.6cm}
	
	\node[anchor = north west, font = \figsmall, align = flush left, xshift = \templength] (light) at (conditions) {\textbf{light}\\2 days at 25\,\textdegree C\\16\,h light/8\,h dark};
	\planticon[height = \templength]{light} at ($(light.west) + (-.5\templength, 0)$);
	
	\node[anchor = north west, font = \figsmall, align = flush left, yshift = -.1\templength] (dark) at (light.south west) {\textbf{dark}\\2 days at 25\,\textdegree C\\constant dark};
	\planticon[height = \templength]{dark} at ($(dark.west) + (-.5\templength, 0)$);
	
	\node[anchor = north west, font = \figsmall, align = flush left, yshift = -.1\templength] (warm) at (dark.south west) {\textbf{warm}\\2 days at 25\,\textdegree C\\1 day at 35\,\textdegree C\\16\,h light/8\,h dark};
	\planticon[height = \templength]{warm} at ($(warm.west) + (-.5\templength, 0)$);
	
	\node[anchor = north west, font = \figsmall, align = flush left, yshift = -.1\templength] (cold) at (warm.south west) {\textbf{cold}\\2 days at 25\,\textdegree C\\1 day at 10\,\textdegree C\\16\,h light/8\,h dark};
	\planticon[height = \templength]{cold} at ($(cold.west) + (-.5\templength, 0)$);

	\node[anchor = north west, font = \figsmall, align = flush left, yshift = -.1\templength] (maize) at (cold.south west) {\textbf{maize}\\16\,h at 25\,\textdegree C\\constant dark};
	\planticon[height = \templength]{maize} at ($(maize.west) + (-.5\templength, 0)$);
	
	\node[anchor = south, rotate = -90] (tobacco) at ($(light.north east)!.5!(cold.south east)$) {\textbf{tobacco}};
	\begin{pgfonlayer}{background}
		\foreach \x in {light, dark, warm, cold} {
			\fill[\x!10] (tobacco.south |- \x.north) -- (\x.north west) to[out = 180, in = 90] ($(\x.west) + (-1.1\templength, 0)$) to[out = -90, in = 180] (\x.south west) -| cycle;
		}
		\fill[tobacco!20, draw = black, thin] (light.north -| tobacco.north) rectangle (cold.south -| tobacco.south);
	\end{pgfonlayer}
	
	
	\node[anchor = base, rotate = -90] at (maize.east -| tobacco.base) {\textbf{maize}};
	\begin{pgfonlayer}{background}
		\fill[maize!10] (tobacco.south |- maize.north) -- (maize.north west) to[out = 180, in = 90] ($(maize.west) + (-1.1\templength, 0)$) to[out = -90, in = 180] (maize.south west) -| cycle;
		\fill[maize!20, draw = black, thin] (maize.north -| tobacco.north) rectangle (maize.south -| tobacco.south);
	\end{pgfonlayer}
	

	%%% PCA of Plant STARR-seq experiments %%%
	\coordinate[xshift = \columnsep] (PCA) at (scheme -| tobacco.north);
	
	\distance{PCA}{array text.south -| \textwidth, 0}
	
	\tikzset{
		PC1/.store in = \PCone,
		PC2/.store in = \PCtwo
	}
	
	\begin{pggfplot}[%
		at = {(PCA -| \textwidth, 0)},
		anchor = north east,
		width = {\xdistance - 1.5\baselineskip},
		height = {-\ydistance - 1.5\baselineskip},
		tight,
		xlabel = {PC2 (\pgfmathprintnumber[precision = 0]{\PCtwo}\%)},
		ylabel = {PC1 (\pgfmathprintnumber[precision = 0]{\PCone}\%)},
		scatter styles = {light,dark,warm,cold,maize},
		ticks = none,
	]{tamsACR_PCA}
	
		\pggf_scatter(
			mark size = 2.5,
			opacity = 0.5,
			style from column = condition
		)
	
		\pggf_annotate(scatter legend*)
	
	\end{pggfplot}
	
	
	%%% enhancer strength by species (light only) %%%
	\coordinate[yshift = -\columnsep] (species) at (scheme |- array text.south);
	
	\planticon[height = .5cm]{light} at ($(species) + (.6, -.25)$);
	
	\begin{pggfplot}[% 
		at = (species),
		total width = \threecolumnwidth - 2.5\baselineskip,
		xlabel = species,
		ylabel = \enrichment,
		xticklabels are csnames,
		col title = light,
		title is csname,
		title color from name,
		enlarge x limits = 0,
		enlarge y limits = {lower = 0.1},
		ytick = {-12, -10, ..., 12},
		xticklabel style = {font = \specieslongfalse},
		only facet = 1,
	]{tamsACR_species}
	
		\pggf_annotate(
			zeroline,
			annotation = {
				\draw[gray, dashed] (\pggfxmin, -2) -- (\pggfxmax, -2) coordinate (repressive) (\pggfxmin, 1) -- (\pggfxmax, 1) coordinate (enhancing);
			}
		)
	
		\pggf_violin(
			style from column = {fill = xticklabels}
		)
	
	\end{pggfplot}
	
	\draw[decorate, decoration = {brace, amplitude = .15cm, raise = .075cm}] (pggfplot c1r1.north east) -- (enhancing -| pggfplot c1r1.east) node[pos = .5, anchor = south, rotate = -90, node font = \figsmall, align = center, yshift = .225cm] {enhancing\\(14\%)};
	\draw[decorate, decoration = {brace, amplitude = .15cm, raise = .075cm}] (enhancing -| pggfplot c1r1.east) -- (repressive -| pggfplot c1r1.east) node[pos = .4, anchor = south, rotate = -90, node font = \figsmall, align = center, yshift = .225cm] {inactive\\(76\%)};
	\draw[decorate, decoration = {brace, amplitude = .15cm, raise = .075cm}] (repressive -| pggfplot c1r1.east) -- (pggfplot c1r1.south east) node[pos = .55, anchor = south, rotate = -90, node font = \figsmall, align = center, yshift = .225cm] {repressive\\(11\%)};
	
	
	%%% correlation of enhancer strength in forward and reverse orientation (light only) %%%
	\coordinate (ori cor) at (species -| \textwidth - \twothirdcolumnwidth, 0);
	
	\planticon[height = .5cm]{light} at ($(ori cor) + (.6, -.25)$);
	
	\begin{pggfplot}[%
		at = {(ori cor)},
		total width = \threecolumnwidth,
		xlabel = {\textbf{forward}: \enrichment},
		ylabel = {\textbf{reverse}: \enrichment},
		facet 1/.append style = {pggf colorbar = {title = count, log2, horizontal}},
		col title = orientation (\usename{light}),
		title color = light,
		xytick = {-12, -10, ..., 12},
		only facet = 1,
		tight y
	]{tamsACR_cor_orientation}
	
		\pggf_annotate(diagonal)
	
		\pggf_hexbin(bins = 50)
	
		\pggf_stats(stats = {n,spearman,rsquare})
	
	\end{pggfplot}
	
	
	%%% orientation-dependence of enhancer strength %%%
	\coordinate (orientation) at (species -| \textwidth - \threecolumnwidth, 0);
	
	\begin{pggfplot}[% 
		at = (orientation),
		total width = \threecolumnwidth,
		xlabel = condition,
		ylabel = $\log_2$(orientation-dependence),
		xticklabels are csnames,
		title = orientation,
		enlarge x limits = 0,
		enlarge y limits = {lower = 0.1},
		ytick = {-12, -11, ..., 12},
		xticklabel style = {font = \vphantom{light}}
	]{tamsACR_orientation}
	
		\pggf_annotate(
			zeroline,
			annotation = {\draw[densely dashed] (\pggfxmin, {log2(1.5)}) -- (\pggfxmax, {log2(1.5)});},
			annotation = {\draw[densely dotted] (\pggfxmin, 1) -- (\pggfxmax, 1);}
		)
	
		\pggf_violin(
			style from column = {fill = xticklabels}
		)
	
	\end{pggfplot}
	
	
	%%% enhancer strength in stable Arabidopsis lines
	\coordinate[yshift = -\columnsep] (DL strength) at (scheme |- xlabel.south);
	
	\planticon[height = .5cm]{light} at ($(DL strength) + (.6, -.25)$);
	
	% define enhancers
	\defname[named]{Zm-23177}{Zm-23177}{PaulTolMuted1}
	\defname[named]{At-1}{At-1}{PaulTolMuted10}
	\defname[named]{At-13524}{At-13524}{PaulTolMuted2}
	\defname[named]{Sl-774}{Sl-774}{PaulTolMuted3}
	\defname[named]{At-2020}{At-2020}{PaulTolMuted4}
	\defname[named]{Sb-11289}{Sb-11289}{PaulTolMuted5}
	\defname[named]{At-11716}{At-11716}{PaulTolMuted6}
	\defname[named]{At-12871}{At-12871}{PaulTolMuted7}
	\defname[named]{At-9661}{At-9661}{PaulTolMuted8}
	\defname[named]{Sl-12881}{Sl-12881}{PaulTolMuted9}
	
	\begin{pggfplot}[%
		at = {(DL strength)},
		total width = \threecolumnwidth,
		ylabel = \dlratio,
		xticklabels are csnames,
		xticklabel style = {rotate = 45, anchor = north east, inner sep = 0.236em},
		col title = stable \textit{Arabidopsis} lines,
		title color = Arabidopsis,
		enlarge x limits = 0,
		enlarge y limits = {lower = 0.1},
		tight y
	]{DL_strength}
		
		\pggf_annotate(zeroline)
	
		\pggf_boxplot(
			style from column = {fill* = xticklabels},
			outlier options = {opacity = .5},
			outlier style from column = {xticklabels},
		)
	
	\end{pggfplot}
	
	
	%%% correlation between Plant STARR-seq and stable Arabidopsis lines
	\coordinate (DL cor) at (DL strength -| \textwidth - \twothirdcolumnwidth, 0);
	
	\planticon[height = .5cm]{light} at ($(DL cor) + (.6, -.25)$);
	
	\begin{pggfplot}[%
		at = {(DL cor)},
		total width = \fourcolumnwidth,
		xlabel = \textbf{Plant STARR-seq:} \enrichment,
		ylabel = \textbf{stable lines:} \dlratio,
		col title = \textit{Arabidopsis} vs. tobacco,
		col title style = {left color = Arabidopsis!20, right color = tobacco!20},
		legend style = {
			at = {(1, .5)},
			anchor = west,
		},
		scatter styles = {none,Zm-23177,At-1,At-13524,Sl-774,At-2020,Sb-11289,At-11716,At-12871,At-9661,Sl-12881},
		ytick = {-6, -3, ..., 18},
		legend cell align = left,
		legend plot pos = left,
	]{DL_cor}
	
		\pggf_trendline()
	
		\pggf_scatter(
			mark size = 2,
			style from column = enhancer,
			y error = CI
		)
	
		\pggf_annotate(scatter legend*)
	
		\pggf_stats(stats = {n,spearman,rsquare})
	
	\end{pggfplot}
	

	%%% correlation of enhancer strength and gene expression data (light only) %%%
	\coordinate (expression) at (DL strength -| \textwidth - \threecolumnwidth, 0);
	
	\planticon[height = .5cm]{light} at ($(expression) + (.6, -.25)$);
	
	\tikzset{
		leaves/.code = {
			\pgfmathsetlength{\templength}{-0.5\templength + 0.2 * rand * \templength}
			\pgfkeysalso{
				OkabeIto4,
				xshift = \templength
			}
		},
		other/.code = {
			\pgfmathsetlength{\templength}{0.5\templength + 0.2 * rand * \templength}
			\pgfkeysalso{
				black,
				xshift = \templength
			}
		},
		set box width/.code = {
			\pgfmathsetlength{\templength}{(\pgfkeysvalueof{/pgfplots/width} / 4) * 0.4}
		}
	}
	
	\def\tissue{\textbf{tissue:}\hspace*{-\groupplotsep}}
	\def\leaves{leaves}
	\def\other{other}

	\begin{pggfplot}[%
		at = {(expression)},
		total width = \threecolumnwidth,
		xlabel = species,
		ylabel = Pearson correlation coefficient,
		scatter styles = {tissue[empty legend],leaves,empty legend,other},
		enlarge x limits = 0,
		enlarge y limits = {lower = 0.1, upper = 0.1},
		xticklabels are csnames,
		col title = gene expression (\usename{light}),
		title color = light,
		xticklabel style = {font = \specieslongfalse},
		legend columns = 2,
		legend image post style = {xshift = -\templength},
		legend style = {name = legend, at = {(1, 1)}},
		only facet = 1
	]{tamsACR_expression}
	
		\pggf_annotate(zeroline)
	
		\pggf_boxplot(
			color = light,
			draw = black,
			forget plot,
			hide outliers
		)
	
		\pggf_scatter(
			style from column = tissue,
			mark size = 1.5,
			opacity = 0.75,
			set box width
		)
	
		\pggf_annotate(
			scatter legend*
		)
	
	\end{pggfplot}
	
	
	%%% subfigure labels %%%
	\subfiglabel{scheme}
	\subfiglabel[xshift = -1.5em]{conditions}
	\subfiglabel{PCA}
	\subfiglabel{species}
	\subfiglabel{ori cor}
	\subfiglabel{orientation}
	\subfiglabel{DL strength}
	\subfiglabel{DL cor}
	\subfiglabel{expression}
	
\end{tikzpicture}%
			\caption{%
				\captiontitle{Characterization of enhancer strength across species and conditions}%
				\subfigtwo{A}{B} Test sequences (170\,bp) were derived from accessible chromatin regions (ACRs) in the genomes of \usename{Arabidopsis} (At), \usename{tomato} (Sl), \usename{maize} (Zm), and \usename{sorghum} (Sb). The array-synthesized test sequences were cloned in the forward and reverse orientation upstream of a 35S minimal promoter driving the expression of a barcoded GFP reporter gene. The plasmid library was subjected to Plant STARR-seq in transiently transformed tobacco leaves and maize leaf protoplasts \parensubfig{A}. After transformation, the plants or protoplasts were subjected to different light and temperature conditions before RNA extraction \parensubfig{B}.\nextentry
				\subfig{C} Principal component analysis of enhancer strength determined in Plant STARR-seq replicate experiments.\nextentry
				\subfig{D} Enhancer strength of test sequences grouped by species of origin. The percentage of enhancing ($\log_2$(enhancer strength) > 1), inactive ($\log_2$(enhancer strength) between 1 and \textminus2), and repressive ($\log_2$(enhancer strength) < \textminus2) sequences is indicated.\nextentry
				\subfigtwo{E}{F} Correlation of \parensubfig{E} and fold-change in \parensubfig{F} enhancer strength between sequences tested in the forward and reverse orientation across the indicated Plant STARR-seq conditions. The dashed and dotted lines in \parensubfig {E} represent a 1.5-fold and 2-fold difference, respectively. \nextentry
				\subfigtwo{G}{H} The enhancer strength of ten ACR sequences selected from the Plant STARR-seq library and a negative control without an enhancer (none) was determined using a dual-luciferase assay in stable \usename{Arabidopsis} lines \parensubfig{G} and compared to the enhancer strength determined by Plant STARR-seq in tobacco leaves in the light \parensubfig{H}.\nextentry
				\subfig{I} Correlation between enhancer strength and gene expression in various tissues of the indicated species. ACR sequences were linked to the closest transcription start site. For each gene, only the ACR sequence with the highest enhancer strength was retained.\nextentry
				Violin plots in \parensubfigtwo{D}{F} represent the kernel density distribution. Box plots inside violin plots and in \parensubfigtwo{G}{I} represent the median (center line), upper and lower quartiles (box limits), 1.5$\times$ interquartile range (whiskers), and outliers (points) for all corresponding samples. Numbers at the bottom of each violin or box plot indicate the number of samples in each group. In the hexbin plot in \parensubfig{D}, color indicates the number of observations in each hexagon. In \parensubfig{H}, the solid line represents a linear regression line and error bars represent the 95\% confidence interval. Pearson's $R^2$, Spearman's $\rho$, and the number ($n$) of samples are indicated in \parensubfigtwo{E}{H}. 
			}%
			\label{fig:enhScreen}%
		\end{fig}
		
		\begin{fig}
			\begin{tikzpicture}
	
	%%% effect of transcription factor binding sites %%%
	\coordinate (TF effects) at (0, 0);
	
	\begin{pggfplot}[%
		at = {(TF effects)},
		yshift = -.5\baselineskip,
		total width = \textwidth,
		height = 2cm,
		xlabel = transcription factor family (motif ID),
		ylabel = condition,
		ylabel add to width = 9.164pt,
		tight y,
		enlargelimits = 0,
		grid = none,
		xticklabel style = {anchor = east, rotate = 90, inner ysep = 0pt, inner xsep = .3333em},
		xticklabels are csnames,
		facet 1/.append style = {
			pggf colorbar = {%
				horizontal,
				title = {$\log_2$(effect size)},
				title style = {yshift = .1em},
				height = .125 * \pgfkeysvalueof{/pgfplots/parent axis width},
				anchor = south west,
				at = {(parent axis.outer south west)},
				shift = {(-3\baselineskip, 2\baselineskip)},
				ytick = {-.3, 0, ..., 1},
				name = colorbar
			}
		},
		colormap 3 colors,
		y dendrogram height = .5cm,
		show y dendrogram
	]{TF_effects}
	
		\pggf_heatmap()
	
	\end{pggfplot}
	
	% mark validation TFs
	\pgfplotstableread[col sep = tab]{rawData/TF_position.tsv}\TFpostable
	
	\distance{pggfplot c1r1.west}{pggfplot c1r1.east}
	\pgfmathsetlength{\templength}{\xdistance/72}
	
	\pgfplotstableforeachcolumnelement{pos}\of\TFpostable\as\TFpos{%
		\node[anchor = south, inner xsep = 0pt] at ($(pggfplot c1r1.north west) + (-.5\templength, 0) + (\TFpos\templength, 0)$) {\textasteriskcentered};
	}
	
	
	%%% scheme of the TFBS shuffling experiment %%%
	\coordinate[yshift = -\columnsep] (shuffle scheme) at (TF effects |- xlabel.south);
	
	\node[anchor = north, font = \figsmall, text width = \fourcolumnwidth - 0.6666em, align = flush center, shift = {(.5\fourcolumnwidth, -.65)}] (seqs text) at (shuffle scheme) {50 genomic sequences with single strong binding site for a given transcription factor};

	\node[anchor = north, fill = DodgerBlue1, text = white, node font = \figsmall\bfseries, inner sep = .15em, shift = {(-.75, -.1)}] (TFBS 1) at (shuffle scheme -| seqs text) {TFBS};
	\node[anchor = south, fill = DodgerBlue2, text = white, node font = \figsmall\bfseries, inner sep = .15em, shift = {(.75, 0)}] (TFBS 3) at (seqs text.north) {TFBS};
	\node[fill = DodgerBlue3, text = white, node font = \figsmall\bfseries, inner sep = .15em] (TFBS 2) at ($(TFBS 1)!.5!(TFBS 3)$) {TFBS};
	
	\begin{pgfonlayer}{background}
		\draw[thick] (shuffle scheme |- TFBS 1) ++(.5, 0) -- ++(\fourcolumnwidth - 2cm, 0);
		\draw[thick] (shuffle scheme |- TFBS 2) ++(1, 0) -- ++(\fourcolumnwidth - 2cm, 0);
		\draw[thick] (shuffle scheme |- TFBS 3) ++(1.5, 0) -- ++(\fourcolumnwidth - 2cm, 0);
	\end{pgfonlayer}
	
	\draw[very thick, -Latex] (seqs text.south) -- ++(0, -.45) coordinate (arrow end);
	
	\node[anchor = north, font = \figsmall, text width = \fourcolumnwidth - 0.6666em, align = flush center, yshift = -.65cm] (shuffle text) at (arrow end) {replace original transcription factor binding site with randomly shuffled version};
	
	\begin{scope}
		\selectcolormodel{gray}
		\node[anchor = north, fill = DodgerBlue1, text = white, node font = \figsmall\bfseries, inner sep = .15em, shift = {(-.75, -.1)}] (TFBS 1) at (arrow end -| seqs text) {FSBT};
		\node[anchor = south, fill = DodgerBlue2, text = white, node font = \figsmall\bfseries, inner sep = .15em, shift = {(.75, 0)}] (TFBS 3) at (shuffle text.north) {TSFB};
		\node[fill = DodgerBlue3, text = white, node font = \figsmall\bfseries, inner sep = .15em] (TFBS 2) at ($(TFBS 1)!.5!(TFBS 3)$) {FBTS};
	\end{scope}
	
	\begin{pgfonlayer}{background}
		\draw[thick] (shuffle scheme |- TFBS 1) ++(.5, 0) -- ++(\fourcolumnwidth - 2cm, 0);
		\draw[thick] (shuffle scheme |- TFBS 2) ++(1, 0) -- ++(\fourcolumnwidth - 2cm, 0);
		\draw[thick] (shuffle scheme |- TFBS 3) ++(1.5, 0) -- ++(\fourcolumnwidth - 2cm, 0);
	\end{pgfonlayer}
	
	\draw[very thick, -Latex] (shuffle text.south) -- ++(0, -.45) coordinate (arrow end);
	
	\node[anchor = north, font = \figsmall, text width = \fourcolumnwidth - 0.6666em, align = flush center, yshift = -.65cm] (measure text) at (arrow end) {measure enhancer strength of sequences with WT or shuffled binding site};
	
	\node[node font = \figsmall\bfseries, align = flush center, yshift = .325cm] (Plant STARR-seq) at (measure text.north) {Plant\\STARR-seq};
	
	\coordinate[xshift = -.3cm] (leaf center) at (Plant STARR-seq.west);
	\planticon*[width = .6cm]{tobaccoLeaf} at (leaf center);
	\begin{scope}[xscale = -1]
		\planticon[width = .323cm]{pipe} at ($(leaf center) + (.35, .16)$);
	\end{scope}
	
	\coordinate[xshift = .3cm] (protoplasts) at (Plant STARR-seq.east);
	\planticon[width = .3cm]{protoplast} at ($(protoplasts) + (90:.15)$);
	\planticon[width = .3cm]{protoplast} at ($(protoplasts) + (162:.15)$);
	\planticon[width = .3cm]{protoplast} at ($(protoplasts) + (234:.15)$);
	\planticon[width = .3cm]{protoplast} at ($(protoplasts) + (306:.15)$);
	\planticon[width = .3cm]{protoplast} at ($(protoplasts) + (18:.15)$);
	\planticon[width = .3cm, rotate = 15]{maize} at ($(protoplasts) + (.375, -.1)$);
	
	
	%%% comparison of TFBS shuffle effect with main library %%%
	\coordinate (correlation) at (shuffle scheme -| \textwidth - \threequartercolumnwidth, 0);
	
	\begin{pggfplot}[%
		at = {(correlation)},
		total width = \threequartercolumnwidth,
		xlabel = {\textbf{transcription factor binding site shuffling:} $\log_2$(effect size)},
		ylabel = {\textbf{ACR library:} $\log_2$(effect size)},
		title is csname,
		title color from name,
		xytick = {-1.2, -.9, ..., 1.2}
	]{TFBS_shuffle_tamsACR}
	
		\pggf_annotate(diagonal)
	
		\pggf_trendline()
	
		\pggf_scatter(
			color from column name,
			mark size = 1.5
		)
	
		\pggf_stats(
			stats = {n,spearman,rsquare}
		)
	
	\end{pggfplot}
	
	
	%%% scheme of synthetic enhancer design %%%
	\coordinate[yshift = -\columnsep] (synTF scheme) at (TF effects |- xlabel.south);
	
	\pgfplotsset{letter colors = DNA}
	
	\pgfmathsetlength{\templength}{.45cm}
	
	\coordinate[yshift = -\columnsep - 2\templength] (TFs) at (synTF scheme);
	
	% nucleotide composition
	\coordinate[shift = {(.5\columnsep + \templength, 1.5\templength)}] (low GC) at (TFs);
	\coordinate[shift = {(.5\columnsep + \templength, -1.5\templength)}] (high GC) at (TFs);
	
	\pgfplotstableread[col sep = tab]{rawData/synEnh_nucFreq.tsv}\nucFreqTable
	\pgfplotstableforeachcolumn\nucFreqTable\as\GCbin{
		\pgfplotstablegetelem{0}{\GCbin}\of\nucFreqTable
		\foreach [expand list, evaluate = \y as \x using 180 - \y, remember = \x as \lastx (initially 180)] \base/\y in {\pgfplotsretval}{
			\fill[letter\base col!50] (\GCbin) -- ++(\lastx:\templength) arc[start angle = \lastx, end angle = \x, radius = \templength] coordinate[pos = .5] (c\base) -- cycle;
			\node[font = \figsmall] at ($(\GCbin)!.6!(c\base)$) {\base};
		}
	}
	
	% sequences
	\coordinate[xshift = .5\columnsep + \templength] (seqs) at (TFs -| low GC);
	\coordinate (low GC seq) at (low GC -| seqs);
	\coordinate (high GC seq) at (high GC -| seqs);
	
	\distance{seqs}{\threecolumnwidth - .5\columnsep, 0}
	
	\foreach \x/\y in {low GC/south, high GC/north}{
		\draw[line width = .06cm, \x] (\x\space seq) node[anchor = \y\space west, text = black, text depth = 0pt, xshift = -.5\columnsep] {25 sequences with \textbf{\x} content} -- ++(\xdistance, 0) coordinate[pos = .165] (\x\space 1) coordinate[pos = .5] (\x\space 2) coordinate[pos = .835] (\x\space 3) coordinate (\x\space seq end);
	}
	
	% TFs
	\foreach \x in {1, 2, 3}{
		\node[fill = DodgerBlue1, text = white, node font = \figsmall\bfseries, inner sep = .15em] (TFBS \x) at (TFs -| low GC \x) {TFBS};
		
		\draw[-Latex, shorten > = .08cm, DodgerBlue1!50] (TFBS \x.north) ++(0, .05) -- (low GC \x);
		\draw[-Latex, shorten > = .08cm, DodgerBlue1!50] (TFBS \x.south) ++(0, -.05) -- (high GC \x);
	}
	
	\node[anchor = base, node font = \figsmall] at ($(TFBS 1.base)!.5!(TFBS 2.base)$) {and/or};
	\node[anchor = base, node font = \figsmall] at ($(TFBS 2.base)!.5!(TFBS 3.base)$) {and/or};


	% tested TFBSs
	\coordinate[yshift = -.5\columnsep - \templength] (TFBSs) at (synTF scheme |- high GC);
	
	\node[anchor = north, font = \bfseries] (title) at ($(TFBSs)!.5!(TFBSs -| low GC seq end)$) {transcription factor binding sites (TFBS)\settowidth{\global\templength}{\figsmall\usename{TF_13}}};

	\pgfplotstableread[col sep = tab]{rawData/synEnh_TFs.tsv}\synEnhTFtable

	\pgfplotstableforeachcolumnelement{col1}\of\synEnhTFtable\as\x{
		\pgfmathsetmacro{\xi}{int(\pgfplotstablerow + 1)}
		\node[anchor = base west, node font = \figsmall, yshift = -\xi*1.15\baselineskip, text width = \templength, align = right] at (TFBSs |- title.base) {\usename{\x}};
	}
	\pgfplotstableforeachcolumnelement{col2}\of\synEnhTFtable\as\x{
		\pgfmathsetmacro{\xi}{int(\pgfplotstablerow + 1)}
		\node[anchor = base east, node font = \figsmall, yshift = -\xi*1.15\baselineskip] (TFBS name \xi) at (low GC seq end |- title.base) {\usename{\x}};
	}

	% connect TFBS list with TFBSs
	\draw[decorate, decoration = {brace, amplitude = .3cm}] ($(TFBS name 1.north east) + (-.1, 0)$) -- ($(TFBS name 6.south east) + (-.1, 0)$) coordinate[pos = .5, xshift = .3cm - \pgflinewidth] (curly);
	\draw[thick, -Latex, shorten > = .05cm, rounded corners = .15cm] (curly) -- ++(.2, 0) |- (TFBS 3.east);


	%%% strength of synthetic enhancers by number of inserted TFBSs %%%
	\coordinate (synEnh strength) at (synTF scheme -| \textwidth - \twothirdcolumnwidth, 0);
	
	% style to shift cld labels
	\tikzset{
		label shift/.code = {
			\ifnum1=#1\relax
				\pgfkeysalso{xshift = .2em}
			\fi
			\def\tmpA{maize}
			\ifx\tmpA\pggfcolname\relax
				\ifnum1=#1\relax
					\pgfkeysalso{yshift = .5ex}
				\fi
				\ifnum2=#1\relax
					\pgfkeysalso{yshift = 1.5ex}
				\fi
			\fi
		}
	}
	
	% group styles
	\pgfplotsset{
		low GC/.code = {
			\def\tmpA{dark}
			\ifx\tmpA\pggfcolname\relax
				\pgfkeysalso{fill = \pggfcolname!50!white}
			\else
				\pgfkeysalso{fill = \pggfcolname}
			\fi
		},
		high GC/.code = {
			\def\tmpA{dark}
			\ifx\tmpA\pggfcolname\relax
				\pgfkeysalso{fill = \pggfcolname}
			\else
				\pgfkeysalso{fill = \pggfcolname!50!black}
			\fi
		}
	}
	
	\begin{pggfplot}[%
		at = {(synEnh strength)},
		total width = \twothirdcolumnwidth,
		xlabel = number of inserted transcription factor binding sites,
		ylabel = \enrichment,
		enlarge x limits = 0,
		enlarge y limits = {lower = 0.25, upper = 0.15},
		title is csname,
		title color from name,
		legend pos = north west,
		ytick = {-12, -10, ..., 12},
		legend columns = 3,
		legend style = {anchor = north, at = {(.5, 1)}},
		legend image width = .75\baselineskip
	]{synEnh_strength}
		
		\pggf_annotate(
			manual legend = {
				{\hspace*{-\groupplotsep}\textbf{GC:}\hspace*{-\groupplotsep}} = {empty legend},
				low = {boxplot legend = y, fill = gray!25!white},
				high = {boxplot legend = y, fill = gray!75!black}
			},
			only facet = 1
		)
		
		\pggf_hline(densely dotted)
	
		\pggf_boxplot(
			style from column = {GC_bin},
			sample size = {anchor = west, rotate = 90},
			cld = {anchor = south, text depth = 0pt, label shift = #1}{max},
			forget plot
		)
	
	\end{pggfplot}
	
	
	%%% subfigure labels
	\subfiglabel[yshift = .5\baselineskip]{TF effects}
	\subfiglabel{shuffle scheme}
	\subfiglabel{correlation}
	\subfiglabel{synTF scheme}
	\subfiglabel{synEnh strength}

\end{tikzpicture}%
			\caption{%
				\captiontitle{Transcription factor binding sites are key determinants of enhancer strength}%
				\subfig{A} Putative transcription factor binding sites were identified by scanning the test sequences for 72 known binding motifs of the indicated transcription factor families. Effect size, calculated as the fold-change in GC-normalized (see Methods) enhancer strength between test sequences with or without one putative binding site, is indicated by color.\nextentry
				\subfigtwo{B}{C} To test the effect of a transcription factor binding site (TFBS), 50 test sequences containing a corresponding binding site were selected and mutated by replacing the binding site with a shuffled version. All sequences were subjected to Plant STARR-seq in the indicated conditions and species \parensubfig{B}. The effect size of 16 different binding sites (marked with asterisks in \plainsubfigref{A}) was determined as the fold-change between the wild-type and mutated sequences, and was compared to the effect size in the ACR library \parensubfig{C}. The solid and dashed lines represent a linear regression line and a $y = x$ line, respectively. Pearson's $R^2$, Spearman's $\rho$, and the number ($n$) of samples are indicated.
				\subfigtwo{D}{E} Synthetic enhancers were designed by inserting one to three transcription factor binding sites into 25 random sequences each with a nucleotide composition similar to an average dicot (low GC) or monocot (high GC) ACR \parensubfig{D}. The synthetic enhancers were subjected to Plant STARR-seq and their enhancer strengths are shown in box plots as defined in \cref{fig:enhScreen} \parensubfig{E}. Letters above the box plots indicate significance groups determined by post-hoc Tukey tests performed separately for each condition. Exact $p$ values are listed in Supplementary Data 4.% TODO: make sure reference is correct
			}%
			\label{fig:TFBS}%
		\end{fig}
		
		\begin{fig}
			\begin{tikzpicture}
	\pggfset{ylabel add to width = \baselineskip}

	%%% plantGREP model %%%
	\coordinate (plantGREP) at (0, 0);
	
	% sequence
	\node[anchor = west, node font = \figsmall\courier, shift = {(0, -1.25)}] (seq) at (plantGREP) {%
		\pgfmathsetseed{928}\randomNucs{10}\,\raisebox{.5ex}{\dots}\,\randomNucs{10}%
	};
	
	\draw[very thick, -Latex] (seq.east) -- ++(.7, 0) coordinate (ohe);
	
	% one-hot encoded sequence
	\node[anchor = west, node font = \figsmall\courier] (ohe) at (ohe) {
		\pgfmathsetseed{928}\randomOhe{10}\,\raisebox{-\baselineskip}{\dots}\,\randomOhe{10}%
	};
	
	\draw[very thick, -Latex] (ohe.east) -- ++(.7, 0) coordinate[xshift = .3333em] (model);
	
	% model
	\node[right = .05cm of model, align = center, draw, fill = OkabeIto3!25, rounded corners, font = \figsmall, yshift = -.5\baselineskip] (bi) {bidirectional\\convolutional\\layer};
	
	\node[right = .5cm of bi, align = center, draw, fill = OkabeIto4!25, rounded corners, font = \figsmall] (den) {\textbf{DenseNet}\vphantom{l}\\\figtiny 6 blocks with\\\figtiny 10 layers};
	
	\node[right = .5cm of den, align = center, draw, fill = OkabeIto7!25, rounded corners, font = \figsmall] (fc) {fully\\connected\\layer};
	
	\draw[very thick, -Latex] (bi) -- (den);
	\draw[very thick, -Latex] (den) -- (fc);
	
	\begin{pgfonlayer}{background}
		\draw[fill = OkabeIto5!25, rounded corners] (model) |- ($(fc.north east) + (.05, .05cm + \baselineskip)$) node[pos = .75, anchor = north, fill = none, draw = none, text = black, font = \bfseries] (model top) {plantGREP} |- ($(bi.south west) + (-.05, -.05)$) coordinate[pos = .25] (model out) coordinate[pos = .75] (model bottom) -- cycle;
	\end{pgfonlayer}
	
	\draw[very thick, -Latex] (model out) ++(.3333em, 0) -- ++(.7, 0) coordinate (predictions);
	
	% predictions
	\pgfmathsetmacro{\predLight}{2 + rand * 4}
	\pgfmathsetmacro{\predDark}{\predLight + rand * 0.5 * \predLight}
	\pgfmathsetmacro{\predWarm}{\predLight + rand * 0.1 * \predLight}
	\pgfmathsetmacro{\predCold}{\predLight + rand * 0.1 * \predLight}
	\pgfmathsetmacro{\predMaize}{1.5 + rand * 3}
	
	\node[anchor = west, node font = \figsmall] (predictions) at (predictions) {%
		\pgfmathsetseed{928}%
		\setlength{\tabcolsep}{.5em}%
		\begin{tabular}{rrrrr}%
			\toprule%
			\textbf{\usename{light}} & \textbf{\usename{dark}} & \textbf{\usename{warm}} & \textbf{\usename{cold}} & \textbf{\usename{maize}} \\%
			\midrule%
			\pgfmathprintnumber[fixed, zerofill, precision = 2]{\predLight} &
			\pgfmathprintnumber[fixed, zerofill, precision = 2]{\predDark} &
			\pgfmathprintnumber[fixed, zerofill, precision = 2]{\predWarm} &
			\pgfmathprintnumber[fixed, zerofill, precision = 2]{\predCold} &
			\pgfmathprintnumber[fixed, zerofill, precision = 2]{\predMaize} \\%
			\bottomrule%
		\end{tabular}%
	};
	
	% text
	\node[anchor = south, font = \strut, yshift = .3em] at (seq |- model top.north) {170\,bp sequence};
	\node[anchor = south, font = \strut, yshift = .3em] at (ohe |- model top.north) {one-hot encode};
	\node[anchor = south, font = \strut, yshift = .3em] at (model top.north) {convolutional neural network model};
	\node[anchor = south, font = \strut, yshift = .3em] (blub) at (predictions |- model top.north) {predict $\log_2$(enhancer strength)};
	
	
	%%% enhancer strength predictions from plantGREP %%%
	\coordinate[yshift = -\columnsep] (prediction) at (plantGREP |- model bottom);
	
	\begin{pggfplot}[%
		at = {(prediction)},
		total width = \textwidth,
		xlabel = {\textbf{measurement}: \enrichment},
		ylabel = {\textbf{prediction}:\\\enrichment},
		facet 1/.append style = {
			pggf colorbar = {%
				title = count,
				log2,
				horizontal,
				ytick = {0, 3, ..., 15}
			}
		},
		title is csname,
		title color from name,
		xytick = {-12, -10, ..., 12},
	]{plantGREP_predictions}
	
		\pggf_annotate(diagonal)
	
		\pggf_hexbin(bins = 50)
	
		\pggf_stats(stats = {n,spearman,rsquare})
	
	\end{pggfplot}
	
	
	%%% example DeepLIFT output %%%
	\coordinate[yshift = -\columnsep] (DeepLIFT) at (plantGREP |- xlabel.south);	
	
	\begin{pggfplot}[%
		at = {(DeepLIFT)},
		yshift = -.4cm,
		total width = \textwidth,
		height = .75cm,
		xlabel = position,
		ylabel = DeepLIFT\\contribution score,
		row title is csname,
		row title color from name,
		enlarge x limits = 0,
		xtick = {1, 10, 20, ..., 170},
		ytick = {-.2, 0, .2, .4},
		title font = \fignormal,
		title style = {inner xsep = 0pt},
	]{DeepLIFT_example}
	
		\pggf_annotate(
			annotation = {
				\coordinate (c1) at (3.5, \pggfymax);
				\coordinate (c2) at (4.5, \pggfymax);
				\coordinate (c3) at (13.5, \pggfymax);
				\coordinate (c4) at (101.5, \pggfymax);
				\coordinate (c5) at (126.5, \pggfymax);
			},
			only facet = 1
		)
	
		\pggf_annotate(
			zeroline,
			annotation = {
				\fill[gray, fill opacity = .1] (3.5, \pggfymin) rectangle (10.5, \pggfymax);
				\fill[gray, fill opacity = .1] (13.5, \pggfymin) rectangle (19.5, \pggfymax);
				\fill[gray, fill opacity = .1] (105.5, \pggfymin) rectangle (109.5, \pggfymax);
				\fill[gray, fill opacity = .1] (126.5, \pggfymin) rectangle (134.5, \pggfymax);
			}
		)
	
		\pggf_logo(letter colors = DNA)
	
	\end{pggfplot}
	\node[anchor = north] at (pggfplot c1r1.north) {sequence: At-12806 (fwd)};
	
	% matching TF motifs
	\distance{c1}{c2}
	
	\begin{pggfplot}[%
		anchor = south west,
		at = {(c1)},
		width = 9\xdistance,
		height = .4cm,
		axis lines = none,
		axis background/.style = {fill = none},
		tight
	]{TF_7_RC}
	
		\pggf_logo(letter colors = DNA)
	
	\end{pggfplot}
	\node[anchor = east, node font = \figtiny] at (pggfplot c1r1.west) {\usename{TF_7}};
	
	\begin{pggfplot}[%
		anchor = south west,
		at = {(c3)},
		width = 6\xdistance,
		height = .4cm,
		axis lines = none,
		axis background/.style = {fill = none},
		tight
	]{TF_13_RC}
	
		\pggf_logo(letter colors = DNA)
	
	\end{pggfplot}
	\node[anchor = west, node font = \figtiny] at (pggfplot c1r1.east) {\usename{TF_13}};
	
	\begin{pggfplot}[%
		anchor = south west,
		at = {(c4)},
		width = 8\xdistance,
		height = .4cm,
		axis lines = none,
		axis background/.style = {fill = none},
		tight
	]{TF_22}
	
		\pggf_logo(letter colors = DNA)
	
	\end{pggfplot}
	\node[anchor = west, node font = \figtiny] at (pggfplot c1r1.east) {partial \usename{TF_22}};
	
	\begin{pggfplot}[%
		anchor = south west,
		at = {(c5)},
		width = 8\xdistance,
		height = .4cm,
		axis lines = none,
		axis background/.style = {fill = none},
		tight
	]{TF_22}
	
		\pggf_logo(letter colors = DNA)
	
	\end{pggfplot}
	\node[anchor = west, node font = \figtiny] at (pggfplot c1r1.east) {\usename{TF_22}};
	
	
	%%% examples of TF-MoDISco patterns and matching TF motifs %%%
	\coordinate[yshift = -\columnsep] (TF matches) at (plantGREP |- xlabel.south);
	
	% modisco patterns
	\begin{pggfplot}[%
		at = {(TF matches)},
		total width = \textwidth,
		height = .75cm,
		xlabel = {},
		ylabel = average DeepLIFT\\contribution score,
		col titles from table = false,
		row title is csname,
		row title color from name,
		enlarge x limits = 0,
		xmajorticks = false,
		ytick = {-.1, 0, .1, .2},
		empty facet warning = false,
		title font = \fignormal,
		title style = {inner xsep = 0pt},
		tight y
	]{modisco_motifs_examples}
	
		\pggf_annotate(
			zeroline,
			skip facets = {6,14,30,36,40}
		)
		
		\pggf_annotate(
			annotation = {
				\node[anchor = north, align = center, node font = \figtiny, gray] at (current axis.north) {no matching\\pattern found};
			},
			only facets = {37, 39}
		)
	
		\pggf_logo(letter colors = DNA)
	
	\end{pggfplot}
	
	% matching TF motifs
	\tikzset{
		TF name/.style = {
			/pgfplots/extra description/.append code = {
				\node[anchor = north, text width = \pgfkeysvalueof{/pgfplots/width}, align = center, node font = \figsmall] at (.5, 0) {\usename{#1}};
			}
		}
	}
	
	\begin{pggfplot}[%
		anchor = north west,
		at = {(pggfplot c1r5.south west)},
		yshift = -3\groupplotsep,
		total width = \textwidth,
		height = .75cm,
		xlabel = {},
		ylabel = information\\content,
		col titles from table = false,
		row title = TFBS,
		row title color = gray,
		enlarge x limits = 0,
		enlarge y limits = {lower = 0},
		ymax = 2,
		ytick = {0, 1, 2},
		xmajorticks = false,
		yticklabel style = {font = \hphantom{0.}},
		title font = \fignormal,
		title style = {inner xsep = 0pt},
		tight y
	]{modisco_TFs_examples}
	
		\pggf_annotate(zeroline)
	
		\pggf_logo(letter colors = DNA)
	
	\end{pggfplot}
	
	
	%%% subfigure labels %%%
	\subfiglabel{plantGREP}
	\subfiglabel{prediction}
	\subfiglabel{DeepLIFT}
	\subfiglabel{TF matches}

\end{tikzpicture}%
			\caption{%
				\captiontitle{A deep learning model predicts enhancer strength and identifies underlying features}%
				\subfigtwo{A}{B} A convolutional neural network model, termed plantGREP, consisting of a bidirectional convolutional layer followed by a DenseNet architecture with six blocks of ten layers each and a final fully connected layer was developed to predict enhancer strength across five conditions using one-hot encoded DNA sequences as input \parensubfig{A}. The model was trained and validated on 90\% of the data. The remaining 10\% of the data (test set) were used to compare model predictions to the experimentally measured enhancer strength \parensubfig{B}. The color in the hexbin plots indicate the number of observations in each hexagon and the dashed line represents a $y = x$ line. Pearson's $R^2$, Spearman's $\rho$, and the number ($n$) of samples are indicated.\nextentry
				\subfig{C} DeepLIFT was used to calculate contribution scores for an example ACR sequence (At-12806). Known transcription factor binding motifs that match the observed pattern are shown above the plot.\nextentry
				\subfig{D} TF-MoDISco was used to identify commonly occurring patterns in the DeepLIFT scores across all sequences of the test set. Examples of patterns detected by TF-MoDISco (top) and matching transcription factor binding motifs (TFBS; bottom) are shown. Patterns were truncated to the size of respective TFBSs.
			}%
			\label{fig:plantGREP}%
		\end{fig}
		
		\begin{fig}
			\begin{tikzpicture}

	%%% in silico evolution scheme %%%
	\coordinate[yshift = -\columnsep] (evo scheme) at (0, 0);
	
	\node[anchor = north west, xshift = .5\columnsep] (spacer) at (evo scheme) {\hphantom{computational model}};
	
	\node[anchor = north, align = center, font = \bfseries, yshift = -.5\columnsep] (evo title) at (spacer.north) {\textit{in silico} evolution};
	
	\node[below = .1cm of evo title, align = center, inner ysep = .2em] (start) {WT sequence};
	
	\node[below = .5cm of start, align = center, inner ysep = .2em] (mutate) {digitally create all\\possible single-base\\substitution variants};
	
	\node[below = .5cm of mutate, align = center, inner ysep = .2em] (score) {score variants with\\\textbf{plantGREP}};
	
	\node[below = .5cm of score, align = center, inner ysep = .2em] (select) {select sequence\\with best score};
	
	\draw[thick, -Latex] (start) -- (mutate);
	
	\draw[thick, -Latex] (mutate) -- (score);
	
	\draw[thick, -Latex] (score) -- (select);
	
	\draw[thick, -Latex, rounded corners = .15cm] (select.east) -- ($(select -| spacer.east) + (\columnsep, 0)$) |- (mutate.east) node[anchor = south, rotate = -90, pos = .25, align = center] (n rounds) {$n$ rounds};
	
	
	%%% in silico evolution: number of rounds %%%
	\coordinate[xshift = \columnsep] (evo rounds) at (evo scheme -| n rounds.north);
	
	\distance{evo rounds}{\textwidth, 0}
	
	\begin{pggfplot}[%
		at = {(evo rounds)},
		total width = \xdistance,
		xlabel = number of \textit{in silico} evolution rounds,
		ylabel = \enrichment,
		enlarge x limits = 0,
		enlarge y limits = {lower = 0.1, upper = 0.1},
		title is csname,
		title color from name,
		xytick = {-12, -10, ..., 12},
	]{evolution_rounds}
	
		\pggf_annotate(zeroline)
	
		\pggf_violin(
			color from col name,
			draw = black,
			shade = left,
			shade color = white,
			cld = {anchor = south, text depth = 0pt}{max}
		)
	
	\end{pggfplot}
	
	
	%%% in silico evolution: species/condition specificity %%%
	\coordinate[yshift = -\columnsep] (evo specificity) at (evo scheme |- xlabel.south);
	
	\colorlet{start}{white}

	\tikzset{
		species-specificity/.style = {left color = tobacco!20, right color = maize!20, middle color = white},
		light-specificity/.style = {left color = light!20, right color = dark!20, middle color = white},
		temperature-specificity/.style = {left color = warm!20, right color = cold!20, middle color = white},
		TRUE/.style = {fill = pgffillcolor!50!white},
		FALSE/.style = {}
	}
	
	\expandafter\def\csname species-specificity\endcsname{species specificity}
	\expandafter\def\csname light-specificity\endcsname{light specificity}
	\expandafter\def\csname temperature-specificity\endcsname{temperature specificity}
	
	\begin{pggfplot}[%
		at = {(evo specificity)},
		total width = \textwidth,
		xlabel = number of \textit{in silico} evolution rounds,
		ylabel = $\log_2$(specificity),
		enlarge x limits = {abs = 1, rel = 0.005},
		enlarge y limits = {lower = 0.1, upper = 0.1},
		title is csname,
		title style from name,
		ytick = {-12, -10, ..., 12},
		xticklabels = {12, 6, 0, 6, 12},
		legend columns = 2,
		legend style = {anchor = north east, at = {(1, 1)}},
		legend image width = .75\baselineskip,
		tight y
	]{evolution_specificity}
	
		\pggf_annotate(
			zeroline,
			manual legend = {\textbf{objective:} = {empty legend}},
		)
		
		\pggf_annotate(
			annotation = {
				\draw[thick, -Latex] (-.25, .25) -- (-.25, 8.5) node[pos = .45, anchor = north, rotate = 90, align = center, font = \figsmall] {more active\\in \textbf{\usename{tobacco}}};
				\draw[thick, -Latex] (6.25, -.25) -- (6.25, -8.5) node[pos = .45, anchor = north, rotate = -90, align = center, font = \figsmall] {more active\\in \textbf{\usename{maize}}};
			},
			manual legend = {
				tobacco-specific = {violin legend = y, fill = tobacco, fill opacity = .5},
				\relax = {empty legend},
				maize-specific = {violin legend = y, fill = maize, fill opacity = .5}
			},
			only facet = 1
		)
		
		\pggf_annotate(
			annotation = {
				\draw[thick, -Latex] (-.25, .5) -- (-.25, 8.5) node[pos = .45, anchor = north, rotate = 90, align = center, font = \figsmall] {more active\\in \textbf{\usename{light}}};
				\draw[thick, -Latex] (6.25, -.5) -- (6.25, -8.5) node[pos = .45, anchor = north, rotate = -90, align = center, font = \figsmall] {more active\\in \textbf{\usename{dark}}};
			},
			manual legend = {
				light-specific = {violin legend = y, fill = light, fill opacity = .5},
				\relax = {empty legend},
				dark-specific = {violin legend = y, fill = dark, fill opacity = .5}
			},
			only facet = 2
		)
		
		\pggf_annotate(
			annotation = {
				\draw[thick, -Latex] (-.25, .5) -- (-.25, 8.5) node[pos = .45, anchor = north, rotate = 90, align = center, font = \figsmall] {more active\\in \textbf{\usename{warm}}};
				\draw[thick, -Latex] (6.25, -.5) -- (6.25, -8.5) node[pos = .45, anchor = north, rotate = -90, align = center, font = \figsmall] {more active\\in \textbf{\usename{cold}}};
			},
			manual legend = {
				warm-specific = {violin legend = y, fill = warm, fill opacity = .5},
				\relax = {empty legend},
				cold-specific = {violin legend = y, fill = cold, fill opacity = .5}
			},
			only facet = 3
		)
	
		\pggf_violin(
			style from column = {fill = color, fade},
			cld = {anchor = south, text depth = 0pt}{max},
			forget plot
		)
	
	\end{pggfplot}
	
	
	%%% subfigure labels %%%
	\subfiglabel{evo scheme}
	\subfiglabel{evo rounds}
	\subfiglabel{evo specificity}

\end{tikzpicture}%
			\caption{%
				\captiontitle{Building species- and condition-specific synthetic enhancers through \textit{in silico} evolution}%
				\subfigrange{A}{C} Improved enhancers were generated through \textit{in silico} evolution using plantGREP \parensubfig{A}. The starting sequences and sequences evolved for six or twelve rounds were subjected to Plant STARR-seq in the indicated conditions and species \parensubfigtwo{B}{C}. The objective was to evolve constitutive enhancers \parensubfig{B} or species- or condition-specific enhancers \parensubfig{C}. Species and condition specificity was calculated as the fold-change in enhancer strength between two species (tobacco [mean of light, dark, warm and cold] and maize) or conditions (light and dark for light specificity; warm and cold for temperature specificity).\nextentry
				Violin plots in \parensubfig{B} and \parensubfig{C} are as defined in \cref{fig:enhScreen} and letters above the violin plots indicate significance groups determined by post-hoc Tukey tests performed separately for each condition, species, or specificity. Exact $p$ values are listed in Supplementary Data 4.% TODO: make sure reference is correct
			}%
			\label{fig:synEnh}%
		\end{fig}
		
		\begin{fig}
			\begin{tikzpicture}
	
	\colorlet{genotype}{gray}
	\colorlet{phenotype}{OkabeIto4}
	

	%%% SlCLV3 promoter bashing %%%
	\coordinate (SlCLV3 bashing) at (0, 0);
	
	\def\WT{WT}
	\expandafter\def\csname CLV3p13\endcsname{\textit{slclv3\textsuperscript{pro-13}}}
	\expandafter\def\csname CLV3p15\endcsname{\textit{slclv3\textsuperscript{pro-15}}}
	\expandafter\def\csname CLV3p24\endcsname{\textit{slclv3\textsuperscript{pro-24}}}
	\expandafter\def\csname CLV3p26\endcsname{\textit{slclv3\textsuperscript{pro-26}}}
	
	\pgfplotsset{
		gRNA/.style = {arrows = -{Triangle[length = .9pt, width = .2cm]}, OkabeIto2},
		DNA/.style = {line width = .1cm, lightgray, -},
		deletion/.style = {line width = .05cm, OkabeIto7, densely dashed, -},
		xaxis/.is choice,
		xaxis/genotype/.style = {
			xtick = {-2.1, -1.5, -1, -.5},
			extra x ticks = {0},
			extra x tick labels = {ATG},
			xticklabel style = {name = xticklabel},
			xmin = -2.1,
			xmax = 0,
			x decimals = 1,
		},
		xaxis/phenotype/.style = {
			xmin = -3.5
		}		
	}
	
	\begin{pggfplot}[%
		at = {(SlCLV3 bashing)},
		total width = \twocolumnwidth,
		height = 2cm,
		xlabel = {\vphantom{position relative to}\\\vphantom{\textit{ZmCLE7} ATG (kb)}},
		ylabel = {\textit{SlCLV3}\\promoter allele},
		ylabel style = {xshift = -2cm/12},
		ylabel add to width = 4\baselineskip,
		label each x axis,
		enlargelimits = 0,
		ytick = {1, ..., 5},
		yticklabels are csnames,
		title color from name,
		axis background/.style = {fill = none},
		empty facet warning = false,
		tight y
	]{SlCLV3_alleles}
	
		\pggf_quiver(
			data = SlCLV3_alleles_DNA,
			DNA
		)
		
		\pggf_quiver(
			data = SlCLV3_alleles_deletion,
			deletion
		)
		
		\pggf_quiver(
			data = SlCLV3_alleles_gRNA,
			gRNA
		)
		
		\pggf_annotate(
			annotation = {
				\coordinate (sep) at (-1.0805, \pggfymax);
			},
			only facet = 1
		)
		
		\pggf_annotate(
			zeroline*,
			only facet = 2
		)
		
		\pggf_boxplot(
			data = SlCLV3_locules,
			fill = OkabeIto7,
			only facet = 2
		)
	
	\end{pggfplot}
	
	\planticon[height = .5cm]{tomatoFruit} at ($(pggfplot c2r1.north east) + (-.3, -.3)$);
	
	% xlabels
	\node[anchor = north, align = center] at (pggfplot c1r1 |- xlabel.north) {position relative to\\\textit{SlCLV3} ATG (kb)};
	\node[anchor = north, align = center] at (pggfplot c2r1 |- xlabel.north) {locule number per fruit};
	
	% save coordinates
	\coordinate (start) at (pggfplot c1r1.west |- xticklabel.south);
	\coordinate (end) at (pggfplot c1r1.east |- xticklabel.south);
	
	
	%%% SlCLV3 Plant STARR-seq %%%
	\coordinate[yshift = -\columnsep] (SlCLV3 STARR) at (SlCLV3 bashing |- xlabel.south);
	
	\tikzset{
		ROIstart/.store in = \ROIstart,
		ROIend/.store in = \ROIend
	}

	\begin{pggfplot}[%
		at = {(SlCLV3 STARR)},
		total width = \twocolumnwidth,
		xlabel = {position relative to \textit{SlCLV3} ATG (kb)},
		ylabel = \enrichment,
		xmin = -2.1,
		xmax = 0,
		enlarge x limits = 0,
		xtick = {-2.1, -1.8, ..., -.1},
		extra x ticks = {0},
		extra x tick labels = {ATG},
		x decimals = 1,
		row title = {Plant STARR-seq (\usename{dark})},
		row title color = dark,
		axis background/.style = {fill = none},
		legend style = {at = {(1, 1)}},
		every legend image post/.append style = {very thick, arrows = -{Latex[scale = .67]}},
		legend image width = transformdirectionx(0.17),
		tight y
	]{SlCLV3_PlantSTARRseq}
		
		\pggf_annotate(
			zeroline,
			annotation = {
				\coordinate (sep 2) at (-1.0805, \pggfymin);
				\fill[dark, fill opacity = .1] (\ROIstart, \pggfymin) coordinate (ROI Sl1) rectangle (\ROIend, \pggfymax) coordinate (ROI Sl2);
			},
			manual legend = {
				experiment = {black},
				prediction = {35Senhancer}
			}
		)
	
		\pggf_quiver(
			data = SlCLV3_PlantSTARRseq_prediction,
			35Senhancer,
			very thick,
			arrows = -{Latex[scale = .67]},
		)
	
		\pggf_quiver(
			data = SlCLV3_PlantSTARRseq_experiment,
			black,
			very thick,
			arrows = -{Latex[scale = .67]},
		)
	
	\end{pggfplot}
	
	\planticon[height = .5cm]{dark} at ($(pggfplot c1r1.north west) + (.3, -.3)$);
	
	% connection to previous plot
	\begin{pgfonlayer}{background}
		\draw[gray] (start) to[out = -90, in = 90] (pggfplot c1r1.north west) (end) to[out = -90, in = 90, looseness = .75] (pggfplot c1r1.north east);
		\draw[gray, dashed] (sep) -- (sep |- start) to[out = -90, in = 90] (sep 2 |- pggfplot c1r1.north) -- (sep 2);
	\end{pgfonlayer}
	
	
	
	%%% ZmCLE7 promoter bashing %%%
	\coordinate (ZmCLE7 bashing) at (SlCLV3 bashing -| \textwidth - \twocolumnwidth, 0);
	
	\expandafter\def\csname CLE7p1\endcsname{\textit{CLE7\textsuperscript{CR-pro1}}}
	\expandafter\def\csname CLE7p4\endcsname{\textit{CLE7\textsuperscript{CR-pro4}}}
	\expandafter\def\csname CLE7p6\endcsname{\textit{CLE7\textsuperscript{CR-pro6}}}
	
	\pgfplotsset{
		xaxis/genotype/.style = {
			xtick = {-2, -1.5, ..., -.1},
			extra x ticks = {0},
			extra x tick labels = {ATG},
			xticklabel style = {name = xticklabel},
			xmin = -2.313,
			xmax = 0,
			x decimals = 1,
		},
		xaxis/phenotype/.style = {
			xtick = {0, .2, ..., 1.5},
			x decimals = 1,
			xmin = -.15,
			xmax = 1.39
		}		
	}
	
	\begin{pggfplot}[%
		at = {(ZmCLE7 bashing)},
		total width = \twocolumnwidth,
		height = 2cm,
		xlabel = {\vphantom{position relative to}\\\vphantom{\textit{ZmCLE7} ATG (kb)}},
		ylabel = {\textit{ZmCLE7}\\promoter allele},
		ylabel style = {xshift = -.2cm},
		ylabel add to width = 4\baselineskip,
		label each x axis,
		enlargelimits = 0,
		ytick = {1, ..., 4},
		yticklabels are csnames,
		title color from name,
		axis background/.style = {fill = none},
		empty facet warning = false,
		tight y
	]{ZmCLE7_alleles}
	
		\pggf_quiver(
			data = ZmCLE7_alleles_DNA,
			DNA
		)
		
		\pggf_quiver(
			data = ZmCLE7_alleles_deletion,
			deletion
		)
		
		\pggf_quiver(
			data = ZmCLE7_alleles_gRNA,
			gRNA
		)
		
		\pggf_annotate(
			annotation = {
				\coordinate (sep) at (-0.6665, \pggfymax);
				\coordinate (sep 2) at (-1.5005, \pggfymax);
			},
			only facet = 1
		)
		
		\pggf_bar(
			data = ZmCLE7_expression,
			xbar,
			fill = OkabeIto5,
			draw = black,
			x error = {se},
			error bars/error bar style = {black, mark size = 2.5},
			only facet = 2
		)
		
		\pggf_annotate(
			zeroline*,
			only facet = 2
		)
		
		\pggf_scatter(
			data = ZmCLE7_expression,
			y filter/.expression = {(y == 3 && x == 0.498057387) ? y - .15 : ((y == 3 && x == 0.499394362) ? y + .15 : y)},
			only facet = 2
		)
	
	\end{pggfplot}
	
	\planticon[height = .5cm]{maize} at ($(pggfplot c2r1.north east) + (-.3, -.3)$);
	
	% xlabels
	\node[anchor = north, align = center] at (pggfplot c1r1 |- xlabel.north) {position relative to\\\textit{ZmCLE7} ATG (kb)};
	\node[anchor = north, align = center] at (pggfplot c2r1 |- xlabel.north) {\textit{ZmCLE7} expression};
	
	% save coordinates
	\coordinate (start) at (pggfplot c1r1.west |- xticklabel.south);
	\coordinate (start top) at (pggfplot c1r1.west);
	\coordinate (end) at (pggfplot c1r1.east |- xticklabel.south);
	
	
	%%% ZmCLE7 %%%
	\coordinate (ZmCLE7 plantGREP) at (SlCLV3 STARR -| ZmCLE7 bashing);
	
	\begin{pggfplot}[%
		at = {(ZmCLE7 plantGREP)},
		total width = \twocolumnwidth,
		xlabel = {position relative to \textit{ZmCLE7} ATG (kb)},
		ylabel = \enrichment,
		xmin = -2.313,
		xmax = 0,
		enlarge x limits = 0,
		xtick = {-2.4, -2.1, ..., -.1},
		extra x ticks = {0},
		extra x tick labels = {ATG},
		x decimals = 1,
		row title = {plantGREP (\usename{maize})},
		row title color = maize,
		axis background/.style = {fill = none},
		tight y
	]{ZmCLE7_prediction_points}
	
		\pggf_annotate(
			zeroline,
			annotation = {
				\coordinate (sep 3) at (-0.6665, \pggfymin);
				\coordinate (sep 4) at (-1.5005, \pggfymin);
				\fill[maize, fill opacity = .15] (\ROIstart, \pggfymin) coordinate (ROI Zm1) rectangle (\ROIend, \pggfymax) coordinate (ROI Zm2);
			}
		)
		
		\pggf_line(
			data = ZmCLE7_prediction_lines,
			35Senhancer,
			thick,
			on layer = axis foreground
		)
		
		\pggf_scatter(
			35Senhancer,
			opacity = .15
		)
	
	\end{pggfplot}
	
	\planticon[height = .5cm]{maize} at ($(pggfplot c1r1.north west) + (.3, -.3)$);
	
	% connection to previous plot
	\begin{pgfonlayer}{background}
		\draw[gray] (start top) -- (start) to[out = -90, in = 90] (pggfplot c1r1.north west) (end) to[out = -90, in = 90, looseness = .75] (pggfplot c1r1.north east);
		\draw[gray, dashed] (sep) -- (sep |- start) to[out = -90, in = 90, looseness = .75] (sep 3 |- pggfplot c1r1.north) -- (sep 3) (sep 2) -- (sep 2 |- start) to[out = -90, in = 90] (sep 4 |- pggfplot c1r1.north) -- (sep 4);
	\end{pgfonlayer}
	
	
	%%% DeepLIFT predictions for best SlCLV3 fragment %%%
	\coordinate[yshift = -\columnsep] (SlCLV3 DeepLIFT) at (SlCLV3 bashing |- xlabel.south);
	
	\begin{pggfplot}[%
		at = {(SlCLV3 DeepLIFT)},
		total width = \textwidth,
		height = 1.5cm,
		xlabel = {position relative to \textit{SlCLV3} ATG (bp)},
		ylabel = contribution\\score,
		ylabel add to width = 8pt,
		col title = DeepLIFT (\usename{dark}),
		col title color = dark,
		enlarge x limits = 0,
		ytick = {-0.2, -0.1, ..., 0.3},
		xtick = {-1550, -1540, ..., -1350}
	]{SlCLV3_DeepLIFT}
		
		\pggf_annotate(
			zeroline,
			annotation = {
				\fill[dark, fill opacity = .1] (-1503.5, \pggfymin) rectangle (-1498.5, \pggfymax) node[pos = .5, anchor = south, yshift = -.75cm, text = black, opacity = 1, text depth = 0pt, node font = \figsmall] {MYB};
				\fill[dark, fill opacity = .1] (-1492.5, \pggfymin) rectangle (-1480.5, \pggfymax) node[pos = .5, anchor = south, yshift = -.75cm, text = black, opacity = 1, text depth = 0pt, node font = \figsmall] {split bZIP?};
				\fill[dark, fill opacity = .1] (-1471.5, \pggfymin) rectangle (-1465.5, \pggfymax) node[pos = .5, anchor = south, yshift = -.75cm, text = black, opacity = 1, text depth = 0pt, node font = \figsmall] {MYB?};
				\fill[dark, fill opacity = .1] (-1432.5, \pggfymin) rectangle (-1427.5, \pggfymax) node[pos = .5, anchor = south, yshift = -.75cm, text = black, opacity = 1, text depth = 0pt, node font = \figsmall] {MYB};
				\fill[dark, fill opacity = .1] (-1412.5, \pggfymin) rectangle (-1405.5, \pggfymax) node[pos = .5, anchor = south, yshift = -.75cm, text = black, opacity = 1, text depth = 0pt, node font = \figsmall] {bZIP};
				\fill[dark, fill opacity = .1] (-1392.5, \pggfymin) rectangle (-1387.5, \pggfymax) node[pos = .5, anchor = south, yshift = -.75cm, text = black, opacity = 1, text depth = 0pt, node font = \figsmall] {MYB};
				\fill[dark, fill opacity = .1] (-1376.5, \pggfymin) rectangle (-1370.5, \pggfymax) node[pos = .5, anchor = south, yshift = -.75cm, text = black, opacity = 1, text depth = 0pt, node font = \figsmall] {MYB};
			}
		)
	
		\pggf_logo(letter colors = DNA)
	
	\end{pggfplot}
	
	\planticon[height = .5cm]{dark} at ($(pggfplot c1r1.north west) + (.3, -.3)$);
	
	\begin{pgfonlayer}{background}
		\draw[gray] (ROI Sl1) -- (pggfplot c1r1.above north west) (ROI Sl2 |- ROI Sl1) -- (pggfplot c1r1.above north east);
	\end{pgfonlayer}
	

	%%% DeepLIFT predictions for best ZmCLE7 region %%%
	\coordinate[yshift = -\columnsep] (ZmCLE7 DeepLIFT) at (SlCLV3 bashing |- xlabel.south);
	
	\begin{pggfplot}[%
		at = {(ZmCLE7 DeepLIFT)},
		total width = \textwidth,
		height = 1.5cm,
		xlabel = {position relative to \textit{ZmCLE7} ATG (bp)},
		ylabel = contribution\\score,
		ylabel add to width = 8pt,
		col title =  DeepLIFT (\usename{maize}),
		col title color = maize,
		enlarge x limits = 0,
		ytick = {-0.1, -0.05, ..., 0.15},
		xtick = {-550, -540, ..., -350}
	]{ZmCLE7_DeepLIFT}
	
		\pggf_annotate(
			zeroline,
			annotation = {
				\fill[maize, fill opacity = .15] (-487.5, \pggfymin) rectangle (-480.5, \pggfymax) node[pos = .5, anchor = south, yshift = -.75cm, text = black, opacity = 1, text depth = 0pt, node font = \figsmall] {bZIP};
				\fill[maize, fill opacity = .15] (-428.5, \pggfymin) rectangle (-424.5, \pggfymax) node[pos = .5, anchor = south, yshift = -.75cm, text = black, opacity = 1, text depth = 0pt, node font = \figsmall] {GATA};
				\fill[maize, fill opacity = .15] (-409.5, \pggfymin) rectangle (-405.5, \pggfymax) node[pos = .5, anchor = south, yshift = -.75cm, text = black, opacity = 1, text depth = 0pt, node font = \figsmall] {GATA};
				\fill[maize, fill opacity = .15] (-403.5, \pggfymin) rectangle (-394.5, \pggfymax) node[pos = .5, anchor = north, yshift = .75cm, text = black, opacity = 1, text depth = 0pt, node font = \figsmall] {2 $\times$ C2H2};
				\fill[maize, fill opacity = .15] (-364.5, \pggfymin) rectangle (-360.5, \pggfymax) node[pos = .5, anchor = south, yshift = -.75cm, text = black, opacity = 1, text depth = 0pt, node font = \figsmall] {GATA};
			}
		)
	
		\pggf_logo(letter colors = DNA)
	
	\end{pggfplot}
	
	\planticon[height = .5cm]{maize} at ($(pggfplot c1r1.north west) + (.3, -.3)$);
	
	\begin{pgfonlayer}{background}
		\draw[maize!50] (ROI Zm1) -- (pggfplot c1r1.above north west) (ROI Zm2 |- ROI Zm1) -- (pggfplot c1r1.above north east);
	\end{pgfonlayer}
	
	
	%%% subfigure labels %%%
	\subfiglabel{SlCLV3 bashing}
	\subfiglabel{SlCLV3 STARR}
	\subfiglabel{ZmCLE7 bashing}
	\subfiglabel{ZmCLE7 plantGREP}
	\subfiglabel{SlCLV3 DeepLIFT}
	\subfiglabel{ZmCLE7 DeepLIFT}

\end{tikzpicture}%
			\caption{%
				\captiontitle{Plant STARR-seq data and models can guide genome engineering of regulatory elements}%
				\subfig{A} Wang \textit{et al.} used CRISPR-Cas9 to introduce deletions in the upstream region of \textit{SlCLV3} and measured the number of locules per tomato fruit\textsuperscript{18}. Increased numbers of locules are indicative of decreased \textit{SlCLV3} expression levels.\nextentry% TODO: make sure references are correct
				\subfig{B} The enhancer strength in the dark of fragments from the \textit{SlCLV3} upstream region was determined experimentally by Plant STARR-seq (black arrows) and predicted using plantGREP (blue arrows). The orientation of the tested fragments is indicated by the direction of the arrows.\nextentry
				\subfig{C} Liu \textit{et al.} used CRISPR-Cas9 to introduce deletions in the upstream region of \textit{ZmCLE7} and measured its expression levels\textsuperscript{19}.\nextentry% TODO: make sure references are correct
				\subfig{D} The enhancer strength in maize of all possible 170-bp fragments from the \textit{ZmCLE7} upstream region was predicted with plantGREP. Points show individual predictions and the line represents a sliding average (window size 20\,bp).\nextentry
				\subfigtwo{E}{F} DeepLIFT contribution scores for regions with enhancer activity (shaded in \plainsubfigref{B} and \plainsubfigref{D}) of the \textit{SlCLV3} \parensubfig{E} and \textit{ZmCLE7} \parensubfig{F} upstream regions were calculated with plantGREP for the indicated conditions. Matches to known transcription factor binding sites are shaded and labeled.\nextentry
				Data in \parensubfig{A} and \parensubfig{C} was obtained from ref. \textsuperscript{18} and ref. \textsuperscript{19}, respectively.% TODO: make sure references are correct
			}%
			\label{fig:genomeScan}%
		\end{fig}
	
	\fi
	%% Main figures end

	%% Extended data figures start
	\ifext
	
		\begin{efig}
			\begin{tikzpicture}

	%%% enhancer strength by species (without light) %%%
	\coordinate (species) at (0, 0);
	
	\planticon[height = .5cm]{light} at ($(species) + (.6, -.25)$);
	
	\begin{pggfplot}[% 
		at = (species),
		total width = \textwidth - 2\baselineskip,
		xlabel = species,
		ylabel = \enrichment,
		xticklabels are csnames,
		col title = light,
		title is csname,
		title color from name,
		enlarge x limits = {lower = 0, upper = .15},
		enlarge y limits = {lower = 0.1},
		ytick = {-12, -10, ..., 12},
		xticklabel style = {font = \specieslongfalse},
		skip facet = 1,
	]{tamsACR_species}
	
		\pggf_annotate(
			zeroline,
			annotation = {
				\draw[gray, dashed] (\pggfxmin, -2) -- (\pggfxmax, -2) coordinate (repressive) (\pggfxmin, 1) -- (\pggfxmax, 1) coordinate (enhancing);
			}
		)
		
		\pggf_annotate(
			annotation = {
				\node[anchor = east, node font = \figsmall] at ($(enhancing)!.5!(\pggfxmax, \pggfymax)$) {19\%};
				\node[anchor = north east, node font = \figsmall] at (\pggfxmax, 0) {69\%};
				\node[anchor = east, node font = \figsmall] at ($(repressive)!.5!(\pggfxmax, \pggfymin)$) {12\%};
			},
			only facet = 1
		)
	
		\pggf_annotate(
			annotation = {
				\node[anchor = east, node font = \figsmall] at ($(enhancing)!.5!(\pggfxmax, \pggfymax)$) {22\%};
				\node[anchor = north east, node font = \figsmall] at (\pggfxmax, 0) {76\%};
				\node[anchor = east, node font = \figsmall] at ($(repressive)!.5!(\pggfxmax, \pggfymin)$) {2\%};
			},
			only facet = 2
		)
	
		\pggf_annotate(
			annotation = {
				\node[anchor = east, node font = \figsmall] at ($(enhancing)!.5!(\pggfxmax, \pggfymax)$) {11\%};
				\node[anchor = north east, node font = \figsmall] at (\pggfxmax, 0) {78\%};
				\node[anchor = east, node font = \figsmall] at ($(repressive)!.5!(\pggfxmax, \pggfymin)$) {11\%};
			},
			only facet = 3
		)
		
		\pggf_annotate(
			annotation = {
				\node[anchor = east, node font = \figsmall] at ($(enhancing)!.5!(\pggfxmax, \pggfymax)$) {5\%};
				\node[anchor = north east, node font = \figsmall] at (\pggfxmax, 0) {84\%};
				\node[anchor = east, node font = \figsmall] at ($(repressive)!.5!(\pggfxmax, \pggfymin)$) {11\%};
			},
			only facet = 4
		)
	
		\pggf_violin(
			style from column = {fill = xticklabels}
		)
	
	\end{pggfplot}

	\draw[decorate, decoration = {brace, amplitude = .15cm, raise = .075cm}] (pggfplot c4r1.north east) -- (enhancing -| pggfplot c4r1.east) node[pos = .5, anchor = south, rotate = -90, node font = \figsmall, align = center, yshift = .225cm] {enhancing};
	\draw[decorate, decoration = {brace, amplitude = .15cm, raise = .075cm}] (enhancing -| pggfplot c4r1.east) -- (repressive -| pggfplot c4r1.east) node[pos = .4, anchor = south, rotate = -90, node font = \figsmall, align = center, yshift = .225cm] {inactive};
	\draw[decorate, decoration = {brace, amplitude = .15cm, raise = .075cm}] (repressive -| pggfplot c4r1.east) -- (pggfplot c4r1.south east) node[pos = .55, anchor = south, rotate = -90, node font = \figsmall, align = center, yshift = .225cm] {repressive};
	

	%%% correlation of enhancer strength in forward and reverse orientation (without light) %%%
	\coordinate[yshift = -\columnsep] (ori cor) at (species |- xlabel.south);
	
	\begin{pggfplot}[%
		at = {(ori cor)},
		total width = \textwidth,
		xlabel = {\textbf{forward}: \enrichment},
		ylabel = {\textbf{reverse}: \enrichment},
		facet 1/.append style = {pggf colorbar = {title = count, log2, horizontal}},
		title is csname,
		title color from name,
		xytick = {-12, -10, ..., 12},
		skip facet = 1,
		tight y
	]{tamsACR_cor_orientation}
	
		\pggf_annotate(diagonal)
	
		\pggf_hexbin(bins = 50)
	
		\pggf_stats(stats = {n,spearman,rsquare})
	
	\end{pggfplot}
	
	
	%%% subfigure labels %%%
	\subfiglabel{species}
	\subfiglabel{ori cor}

\end{tikzpicture}%
			\caption{%
				\captiontitle{Regulatory activity spans large dynamic range and is orientation-independent}%
				\subfig{A} Violin plots (as defined in \cref{fig:enhScreen}) of the enhancer strength of test sequences grouped by species of origin. The percentage of enhancing ($\log_2$(enhancer strength) > 1), inactive ($\log_2$(enhancer strength) between 1 and \textminus2), and repressive ($\log_2$(enhancer strength) < \textminus2) sequences is indicated.\nextentry
				\subfig{B} Hexbin plots (as defined in \cref{fig:enhScreen}) for the correlation of enhancer strength between sequences tested in the forward and reverse orientation.
			}%
			\label{efig:species+ori}
		\end{efig}
		
		\begin{efig}
			\begin{tikzpicture}

	%%% enhancer strength by chromatin accessibility %%%
	\coordinate (accessibility) at (0, 0);
	
	\begin{pggfplot}[%
		at = (accessibility),
		total width = \textwidth,
		xlabel = chromatin accessibility,
		ylabel = \enrichment,
		title is csname,
		title color from name,
		enlarge x limits = 0,
		enlarge y limits = {lower = 0.1, upper = 0.1},
		ytick = {-12, -10, ..., 12}
	]{tamsACR_cutcount}
	
		\pggf_annotate(zeroline)
		
		\pggf_violin(
			color from col name,
			draw = black,
			shade = left,
			shade color = white,
			cld = {anchor = south, text depth = 0pt}{max}
		)
	
	\end{pggfplot}
	
	
	%%% enhancer strength by histone modifications %%%
	\coordinate[yshift = -\columnsep] (histones) at (accessibility |- xlabel.south);
	
	\planticon[height = .5cm]{arabidopsis} at ($(histones) + (.6, -.25)$);
	
	\expandafter\def\csname none\endcsname{\vphantom{$+$}$-$\\\vphantom{$+$}$-$}
	\expandafter\def\csname ac\endcsname{$+$\\\vphantom{$+$}$-$}
	\expandafter\def\csname me\endcsname{\vphantom{$+$}$-$\\$+$}
	\expandafter\def\csname ac+me\endcsname{$+$\\$+$}
	
	\begin{pggfplot}[%
		at = {(histones)},
		total width = \textwidth,
		xlabel = histone modification,
		ylabel = \enrichment,
		xticklabels are csnames,
		xticklabel style = {align = center, name = xticklabel},
		ytick = {-12, -10, ..., 12},
		enlarge x limits = 0,
		enlarge y limits = {lower = 0.1, upper = 0.1},
		title is csname,
		title color from name,
		tight y
	]{tamsACR_histones}
	
		\pggf_annotate(zeroline)
	
		\pggf_violin(
			cld = {anchor = south, text depth = 0pt}{max}
		)
	
	\end{pggfplot}
	
	\node[anchor = south east, align = right, node font = \figsmall, text depth = 0pt, inner xsep = 0pt, xshift = 3pt] at (pggfplot c1r1.west |- xticklabel.south) {\vphantom{$+$}\smash{acetylation}\\\vphantom{$+$}\smash{methylation}};
	
	
	%%% enhancer strength by genomic region %%%
	\coordinate[yshift = -\columnsep] (region) at (accessibility |- xlabel.south);
	
	\expandafter\def\csname intergenic\endcsname{intergenic\xspace}
	\expandafter\def\csname 5'UTR\endcsname{5\textprime~UTR\xspace}
	\expandafter\def\csname CDS\endcsname{CDS\xspace}
	\expandafter\def\csname intron\endcsname{intron\xspace}
	\expandafter\def\csname 3'UTR\endcsname{3\textprime~UTR\xspace}
	\expandafter\def\csname ncRNA\endcsname{ncRNA\xspace}
	
	\begin{pggfplot}[%
		at = (region),
		total width = \textwidth,
		xlabel = region,
		ylabel = \enrichment,
		xticklabels are csnames,
		xticklabel style = {rotate = 45, anchor = north east, inner sep = 0.236em},
		title is csname,
		title color from name,
		enlarge x limits = 0,
		enlarge y limits = {lower = 0.1, upper = 0.1},
		ytick = {-12, -10, ..., 12},
		sample size = {font = \fontsize{4}{5}\selectfont},
		tight y,
	]{tamsACR_region}
	
		\pggf_annotate(zeroline)
	
		\pggf_violin(
			draw = black,
			cld = {anchor = south, text depth = 0pt}{max}
		)
	
	\end{pggfplot}


	%%% enhancer strength by distance to closest TSS %%%
	\coordinate[yshift = -\columnsep] (TSS dist) at (accessibility |- xlabel.south);
	
	\expandafter\def\csname u5+\endcsname{$>5$\xspace}
	\expandafter\def\csname u5\endcsname{5\textendash1\xspace}
	\expandafter\def\csname u1\endcsname{1\textendash0\xspace}
	\expandafter\def\csname d1\endcsname{0\textendash1\xspace}
	\expandafter\def\csname d5\endcsname{1\textendash5\xspace}
	\expandafter\def\csname d5+\endcsname{$>5$\xspace}
	
	\begin{pggfplot}[% 
		at = (TSS dist),
		total width = \textwidth,
		xlabel = distance to closest transcription start site (kb),
		ylabel = \enrichment,
		xlabel style = {yshift = -.667\baselineskip},
		xticklabels are csnames,
		xticklabel style = {name = xticklabel},
		title is csname,
		title color from name,
		enlarge x limits = 0,
		enlarge y limits = {lower = 0.1, upper = 0.1},
		ytick = {-12, -10, ..., 12},
		sample size = {font = \fontsize{4}{5}\selectfont},
		tight y,
	]{tamsACR_TSS_dist}
	
		\pggf_annotate(
			zeroline,
			annotation = {
				\coordinate (c1) at (1, 0);
				\coordinate (c2) at (2, 0);
			}
		)
	
		\pggf_violin(
			color from col name,
			draw = black,
			shade = right,
			shade color = white,
			cld = {anchor = south, text depth = 0pt}{max}
		)
	
	\end{pggfplot}
	
	\distance{c1}{c2}
	
	\foreach \x in {1, ..., 5}{
		\draw[serif cm-serif cm] (pggfplot c\x r1.west |- xticklabel.south) ++(.05\xdistance, 0) -- ++(2.9\xdistance, 0) node[pos = .5, anchor = north, node font = \figsmall] {upstream};
		\draw[serif cm-serif cm] (pggfplot c\x r1.east |- xticklabel.south) ++(-.05\xdistance, 0) -- ++(-2.9\xdistance, 0) node[pos = .5, anchor = north, node font = \figsmall] {downstream};	
	}


	%%% subfigure labels %%%
	\subfiglabel{accessibility}
	\subfiglabel{histones}
	\subfiglabel{region}
	\subfiglabel{TSS dist}

\end{tikzpicture}%
			\caption{%
				\captiontitle{Genomic signatures of enhancers}%
				\subfigrange{A}{D} Enhancer strength of test sequences grouped by the relative accessibility of the corresponding region in the genome \parensubfig{A}, by the occurrence of acetylated and/or methylated histones in a 1-kb window surrounding the corresponding region in the genome \parensubfig{B}, by the genomic region from which they are derived \parensubfig{C}, or by their distance to the closest transcription start site \parensubfig{D}.\nextentry
				Violin plots are as defined in \cref{fig:enhScreen} and letters above the plots indicate significance groups determined by post-hoc Tukey tests performed separately for each condition and species. Exact $p$ values are listed in Supplementary Data 4.\nextentry% TODO: make sure reference is correct
				UTR, untranslated region; CDS, coding sequence; ncRNA, non-coding RNA.
			}%
			\label{efig:genomic_signature}
		\end{efig}
		
		\begin{efig}
			\begin{tikzpicture}

	%%% correlation of enhancer strength and gene expression data (without light) %%%
	\coordinate (gene expression) at (0, 0);
	
	\tikzset{
		leaves/.code = {
			\pgfmathsetlength{\templength}{-0.5\templength + 0.2 * rand * \templength}
			\pgfkeysalso{
				OkabeIto4,
				xshift = \templength
			}
		},
		other/.code = {
			\pgfmathsetlength{\templength}{0.5\templength + 0.2 * rand * \templength}
			\pgfkeysalso{
				black,
				xshift = \templength
			}
		},
		set box width/.code = {
			\pgfmathsetlength{\templength}{(\pgfkeysvalueof{/pgfplots/width} / 4) * 0.4}
		}
	}
	
	\def\tissue{\textbf{tissue:}\hspace*{-\groupplotsep}}
	\def\leaves{leaves}
	\def\other{other}

	\begin{pggfplot}[%
		at = {(gene expression)},
		total width = \threequartercolumnwidth,
		xlabel = species,
		ylabel = Pearson correlation coefficient,
		scatter styles = {tissue[empty legend],leaves,empty legend,other},
		enlarge x limits = 0,
		enlarge y limits = {lower = 0.1, upper = 0.1},
		xticklabels are csnames,
		title is csname,
		title color from name,
		xticklabel style = {font = \specieslongfalse},
		legend columns = 2,
		legend image post style = {xshift = -\templength},
		legend style = {name = legend, at = {(1, 1)}},
		skip facet = 1
	]{tamsACR_expression}
	
		\pggf_annotate(zeroline)
	
		\pggf_boxplot(
			color from column name,
			draw = black,
			forget plot,
			hide outliers
		)
	
		\pggf_scatter(
			style from column = tissue,
			mark size = 1.5,
			opacity = 0.75,
			set box width
		)
	
		\pggf_annotate(
			scatter legend*,
			only facet = 1
		)
	
	\end{pggfplot}
	
	
	%%% differentially expressed genes in light and dark %%%
	\coordinate (light specificity) at (gene expression -| \textwidth - \fourcolumnwidth, 0);
	
	\planticon[height = .5cm]{arabidopsis} at ($(light specificity) + (.6, -.25)$);
	
	\expandafter\def\csname dark_down\endcsname{downregulated}
	\expandafter\def\csname dark_up\endcsname{upregulated}
	
	\begin{pggfplot}[%
		at = {(light specificity)},
		total width = \fourcolumnwidth,
		xlabel = gene expression in dark,
		ylabel = linked ACR sequences (\%),
		ymin = 0,
		enlarge x limits = 0,
		enlarge y limits = {lower = 0, upper = .2},
		xticklabels are csnames,
		col title = light specificity,
		col title style = {left color = light!20, right color = dark!20, middle color = white},
		legend columns = 3,
		legend style = {anchor = north, at = {(.5, 1)}}
	]{tamsACR_dark_light_specificity}
	
		\pggf_bar(
			ybar,
			draw = black,
			style from column = {fill = group}
		)
	
	\end{pggfplot}
	
	
	%%% condition-/tissue-specificity of genes and corresponding ACRs %%%
	\coordinate[yshift = - \columnsep] (specificity) at (gene expression |- xlabel.south);
	
	\begin{pggfplot}[%
		at = {(specificity)},
		total width = \fourcolumnwidth,
		xlabel = tissue specificity,
		ylabel = condition specificity,
		enlarge y limits = {lower = 0.1, upper = 0.1},
		col title = specificity,
	]{tamsACR_specificity}
	
		\pggf_annotate(zeroline)
	
		\pggf_violin(
			fill = OkabeIto6,
			shade = left,
			shade color = white,
			cld = {anchor = south, text depth = 0pt}{max}
		)
	
	\end{pggfplot}
	
	
	%%% enriched GO terms %%%
	\coordinate (GO terms) at (specificity -| \textwidth - \threequartercolumnwidth, 0);
	
	\distance{col title.south}{col title.north}
	
	\begin{pggfplot}[%
		at = {(GO terms)},
		total width = \threequartercolumnwidth,
		xlabel = species,
		ylabel = GO term,
		ylabel add to width = 70pt,
		xticklabels are csnames,
		title is csname,
		title color from name,
		xticklabel style = {font = \specieslongfalse},
		y dir = reverse,
		enlargelimits = 0,
		colormap = {TF}{color(0) = (gray!20) color(1) = (OkabeIto4)},
		yticklabel style = {name = yticklabel},
		legend style = {anchor = east, at = {(0, 1)}, yshift = .5\ydistance - .5\pgflinewidth},
		legend columns = 3
	]{tamsACR_GO-terms}
	
		\pggf_heatmap(
			forget plot
		)
		
		\pggf_annotate(
			manual legend = {
				{\textbf{significant:}\hspace*{-\groupplotsep}} = {empty legend},
				yes = {area legend, fill = OkabeIto4},
				no = {area legend, fill = gray!20}
			},
			only facet = 1
		)
	
	\end{pggfplot}
	
	
	%%% subfigure labels %%%
	\subfiglabel{gene expression}
	\subfiglabel{light specificity}
	\subfiglabel{specificity}
	\subfiglabel{GO terms}
	
\end{tikzpicture}%
			\caption{%
				\captiontitle{Enhancer strength is correlated with gene expression \textit{in planta}}%
				\subfig{A} Correlation between enhancer strength and gene expression in various tissues of the indicated species. At, \usename{Arabidopsis}; Sl, \usename{tomato}; Zm, \usename{maize}; Sb, \usename{sorghum}. ACR sequences were linked to the closest transcription start site. For each gene, only the ACR sequence with the highest enhancer strength was retained.\nextentry
				\subfig{B} For genes down- or upregulated after shifting \usename{Arabidopsis} plants to the dark, the percentage of linked ACR sequences showing light-specific (1.5$\times$ stronger in the light than in the dark; light), dark-specific (1.5$\times$ stronger in the dark than in the light; dark), or constitutive (none) enhancer activity is indicated.\nextentry
				\subfig{C} Condition-specific enhancer activity of ACR sequences grouped by the relative tissue specificity of the linked genes. Condition and tissue specificity was determined using the tau index based on our Plant STARR-seq data and the gene expression data used in \parensubfig{A}, respectively. For each gene, only the ACR sequence with the highest condition specificity was retained. Violin plots are as defined in \cref{fig:enhScreen} and letters above the plots indicate significance groups determined by a post-hoc Tukey test. Exact $p$ values are listed in Supplementary Data 4.\nextentry% TODO: make sure reference is correct
				\subfig{D} GO term enrichment for genes linked to the top 5\% of test sequences with the highest enhancer strength in the indicated conditions and species. Significant (adjusted $p$ value \textle{} 0.05) enrichment is indicated in green; non-significant results are in gray. See Supplementary Data 2 for all GO terms and $p$ values.\nextentry
				Gene expression data was obtained from ref. \textsuperscript{33\textendash{}36}.% TODO: make sure references are correct
			}%
			\label{fig:gene_expression}%
		\end{efig}
		
		\begin{efig}
			\begin{tikzpicture}

	\pgfplotsset{set layers}
	
	\makeatletter
		\setlength\templength{\pggf@ylabelwidth}
	\makeatother

	%%% enhancer strength by GC content %%%
	\coordinate (GC) at (0, 0);
	
	\begin{pggfplot}[%
		anchor = north west,
		at = {($(GC) + (\templength, 0)$)},
		total width = \textwidth,
		xlabel = GC content (\%),
		ylabel = \enrichment,
		xticklabels are bins,
		ytick = {-12, -10, ..., 12},
		title is csname,
		title color from name,
		enlarge x limits = 0,
		enlarge y limits = {lower = 0.1, upper = 0.1},
	]{tamsACR_GC}
	
		\pggf_annotate(zeroline)
	
		\pggf_violin(
			color from col name,
			draw = black,
			shade = right,
			shade color = white,
			cld = {anchor = south, text depth = 0pt}{max}
		)
	
	\end{pggfplot}
	
	\coordinate (GC plot top) at (pggfplot c1r1.above north);
	
	
	%%% distribution of GC content in tamsACR sequences %%%
	\coordinate[yshift = -\columnsep] (GC density) at (GC |- xlabel.south);
	
	\begin{pggfplot}[%
		anchor = north east,
		at = {($(GC density) + (\twothirdcolumnwidth, 0)$)},
		total width = \twothirdcolumnwidth,                     
		xlabel = GC content (\%),
		ylabel = species,
		yticklabels are csnames,
		ylabel add to width = 22pt,
		scale ridges = 2,
		enlarge x limits = 0,
		enlarge y limits = {lower = .075},
		tight y
	]{GC_distribution}
	
		\pggf_ridgeline(
			close density,
			style from column = yticklabels,
			draw = black,
			fill = pgffillcolor!50!white
		)
	
	\end{pggfplot}
	
	
	%%% compare average GC content of ACRs to whole genome %%%
	\coordinate (ACR vs genome) at (GC density -| \textwidth - \threecolumnwidth, 0);
	
	% get ymin and ymax from last plot
	\pgfplotstableread[col sep = tab]{rawData/GC_distribution_axes.tsv}\axistable
	
	\pgfplotstablegetelem{0}{ymin}\of\axistable
	\edef\ymin{\expandonce{\pgfplotsretval}}
	
	\pgfplotstablegetelem{0}{ymax}\of\axistable
	\edef\ymax{\expandonce{\pgfplotsretval}}

	\begin{pggfplot}[%
		anchor = north east,
		at = {($(ACR vs genome) + (\threecolumnwidth, 0)$)},
		total width = \threecolumnwidth,
		xlabel = GC content (\%),
		ylabel = species,
		yticklabels are csnames,
		ylabel add to width = 22pt,
		ymin = \ymin,
		ymax = \ymax,
		scale ridges = 2,
		enlarge x limits = .095,
		enlarge y limits = {lower = .075},
		scatter styles = {%
			At:ACR[{At, mark = solido}],%
			At:genome[{At, mark = solidsquare}],%
			Sb:ACR[{Sb, mark = solido}],%
			Sb:genome[{Sb, mark = solidsquare}],%
			Sl:ACR[{Sl, mark = solido}],%
			Sl:genome[{Sl, mark = solidsquare}],%
			Zm:ACR[{Zm, mark = solido}],%
			Zm:genome[{Zm, mark = solidsquare}]%
		},
		legend image post style = {gray},
		legend columns = 2,
		legend style = {at = {(.5, .97)}, anchor = north},
		tight y
	]{GC_ACR_genome}
	
		\pggf_scatter(
			style from column = group,
			mark size = 3
		)
	
		\pggf_quiver(
			very thick,
			arrows = -{Latex[scale = .67]},
			on layer = axis foreground
		)
		
		\pggf_annotate(
			legend = {ACR, genome}
		)
	
	\end{pggfplot}
	
	
	%%% subfigure labels %%%
	\subfiglabel{GC |- GC plot top}
	\subfiglabel{GC density}
	\subfiglabel{ACR vs genome}

\end{tikzpicture}%
			\caption{%
				\captiontitle{GC content affects enhancer strength in \usename{tobacco}}%
				\subfig{A} Violin plots (as defined in \cref{fig:enhScreen}) of the enhancer strength of test sequences grouped by GC content. Letters above the plots indicate significance groups determined by post-hoc Tukey tests performed separately for each condition and species. Exact $p$ values are listed in Supplementary Data 4.\nextentry% TODO: make sure reference is correct
				\subfig{B} Distribution of GC content for the test sequences in the ACR library.\nextentry
				\subfig{C} Average GC content of the ACR sequences (circles) and the whole genome (squares) of the indicated species.
			}%
			\label{efig:GC_content}
		\end{efig}
		
		\begin{efig}
			\usetikzlibrary{decorations.markings}

\begin{tikzpicture}
	
	\defname[named]{strong_enh}{strong enhancers}{35Senhancer}
	\defname[named]{medium_enh}{medium enhancers}{strong_enh}\colorlet{medium_enh}{medium_enh!67!black}
	\defname[named]{weak_enh}{weak enhancers}{strong_enh}\colorlet{weak_enh}{weak_enh!33!black}
	\defname[named]{inactive}{inactive sequences}{gray}
	\defname[named]{repressive}{repressive elements}{Firebrick1}
	
	%%% scheme of category definition %%%
	\coordinate (category definition) at (0, 0);
	
	\tikzset{
		save coordinate/.style = {
			postaction = decorate,
			decoration = {
				markings,
				mark = at position 0.55 with {
					\coordinate (#1);
				}
			}
		},
		mean/.style = {
			densely dashed,
			save coordinate = mean
		},
		mean-1/.style = {
			densely dotted,
			save coordinate = mean-1
		},
		mean+1/.style = {
			densely dotted,
			save coordinate = mean+1
		},
		mean+2/.style = {
			densely dotted,
			save coordinate = mean+2
		},
		mean+3/.style = {
			densely dotted,
			save coordinate = mean+3
		},
		max/.style = {
			draw = none,
			save coordinate = max
		},
		min/.style = {
			draw = none,
			save coordinate = min
		}
	}
	
	\begin{pggfplot}[%
		at = {(category definition)},
		total width = \fourcolumnwidth,
		xlabel = {},
		ylabel = {\enrichment[GC-normalized]},
		ylabel style = {inner xsep = 0pt},
		ylabel add to width = .33\baselineskip,
		enlarge x limits = {rel = 0, abs = .75},
		enlarge y limits = {lower = 0.1},
		xtick = {1, 5},
		xticklabels = {non-ACR\\sequences, all test\\sequences},
		xticklabel style = {align = center},
		col title = light,
		title is csname,
		title color from name,
		axis background/.style = {fill = none},
		tight y
	]{category_def}
	
		\pggf_violin(
			fill = light
		)
	
		\pggf_hline(
			style from column = intercept_id
		)
		
		\pggf_annotate(
			annotation = {
				\node[anchor = east, fill = white, font = \figsmall, xshift = -.5em] at (mean) {mean};
				\draw[{|<[width = 3]}-{>|[width = 3]}, shorten both = -.5\pgflinewidth] (mean) -- (mean-1) node[pos = .5, anchor = west, font = \figsmall] {sd};
				\draw[{|<[width = 3]}-{>|[width = 3]}, shorten both = -.5\pgflinewidth] (mean) -- (mean+1) node[pos = .5, anchor = west, font = \figsmall] {sd};
				\draw[{|<[width = 3]}-{>|[width = 3]}, shorten both = -.5\pgflinewidth] (mean+1) -- (mean+2) node[pos = .5, anchor = west, font = \figsmall] {sd};
				\draw[{|<[width = 3]}-{>|[width = 3]}, shorten both = -.5\pgflinewidth] (mean+2) -- (mean+3) node[pos = .5, anchor = west, font = \figsmall] {sd};
			}
		)
	
	\end{pggfplot}
	
	\begin{pgfonlayer}{background}
		\shade[left color = white, right color = strong_enh!50] (mean+3) -| (max -| pggfplot c1r1.east) node[pos = .75, anchor = west, node font = \figsmall, xshift = .225cm] {strong enhancers} -- (pggfplot c1r1.north east) -| cycle;
		\shade[left color = white, right color = medium_enh!50] (mean+2) -| (mean+3 -| pggfplot c1r1.east) node[pos = .75, anchor = west, node font = \figsmall, xshift = .35cm] {medium enhancers} -| cycle;
		\shade[left color = white, right color = weak_enh!50] (mean+1) -| (mean+2 -| pggfplot c1r1.east) node[pos = .75, anchor = west, node font = \figsmall, xshift = .225cm] {weak enhancers} -| cycle;
		\shade[left color = white, right color = black!50] (mean-1) -| (mean+1 -| pggfplot c1r1.east) node[pos = .75, anchor = west, node font = \figsmall, xshift = .35cm] (inactive) {inactive sequences} -| cycle;
		\shade[left color = white, right color = repressive!50] (mean-1) -| (min -| pggfplot c1r1.east) node[pos = .75, anchor = west, node font = \figsmall, xshift = .225cm] {repressive elements} -- (pggfplot c1r1.south east) -| cycle;
	\end{pgfonlayer}
	
	\draw[decorate, decoration = {brace, amplitude = .15cm, raise = .075cm}] (max -| pggfplot c1r1.east) -- (mean+3 -| pggfplot c1r1.east);
	\draw[decorate, decoration = {brace, amplitude = .15cm, raise = .2cm}] (mean+3 -| pggfplot c1r1.east) -- (mean+2 -| pggfplot c1r1.east);
	\draw[decorate, decoration = {brace, amplitude = .15cm, raise = .075cm}] (mean+2 -| pggfplot c1r1.east) -- (mean+1 -| pggfplot c1r1.east);
	\draw[decorate, decoration = {brace, amplitude = .15cm, raise = .2cm}] (mean+1 -| pggfplot c1r1.east) -- (mean-1 -| pggfplot c1r1.east);
	\draw[decorate, decoration = {brace, amplitude = .15cm, raise = .075cm}] (mean-1 -| pggfplot c1r1.east) -- (min -| pggfplot c1r1.east);


	%%% number of sequences in each category %%%
	\coordinate[xshift = \columnsep] (category numbers) at (category definition -| inactive.east);
	
	\distance{category numbers}{\textwidth, 0}
	
	\begin{pggfplot}[
		at = {(category numbers)},
		total width = \xdistance,
		xlabel = number of test sequences (\times \pgfmathprintnumber{1000}),
		ylabel = {~},
		ylabel style = {inner sep = 0pt},
		ylabel add to width = 34pt,
		yticklabels are csnames,
		title is csname,
		title color from name,
		xmin = 0,
		enlarge x limits = {lower = 0},
		scaled x ticks = manual:{}{\pgfmathparse{##1/1000}},
		y dir = reverse,
		hide legend
	]{category_numbers}
	
		\pggf_bar(
			xbar,
			draw = black,
			style from column = {fill = yticklabels},
			point meta = x,
			nodes near coords,
			nodes near coords align = {horizontal},
			nodes near coords style = {
				text = black,
				opacity = 1,
				font = \figsmall,
				text depth = 0pt,
				/pgf/number format/fixed
			},
			coordinate style/.condition = {x > 150000}{anchor = east}
		)
	
	\end{pggfplot}


	%%% correlation between TF effects and rel. enrichment of TFBSs in sequences grouped by their strength %%%
	\coordinate[yshift = -\columnsep] (correlation) at (category definition |- xlabel.south);
	
	\begin{pggfplot}[%
		at = {(correlation)},
		total width = \textwidth,
		xlabel = {$\log_2$(effect size)},
		ylabel = {$\log_2$(average TFBS counts, normalized to inactive sequences)},
		ylabel add to width = .33\baselineskip,
		xtick = {-1.2, -.9, ..., 1.2},
		title is csname,
		title color from name,
	]{categories_TFs}
	
		\pggf_scatter(color from col name)
	
		\pggf_trendline(/pgfplots/color from row name)
	
		\pggf_stats(stats = {slope, goodness of fit})
	
	\end{pggfplot}


	%%% subfigure labels %%%
	\subfiglabel{category definition}
	\subfiglabel{category numbers}
	\subfiglabel{correlation}

\end{tikzpicture}%
			\caption{%
				\captiontitle{Activating transcription factor binding sites are enriched in active enhancers}%
				\subfigtwo{A}{B} Test sequences in the ACR library were categorized based on the mean (mean) and standard deviation (sd) of the GC-normalized enhancer strength of non-ACR sequences. The categorization scheme is shown for the light condition as an example \parensubfig{A}. The number of sequences in each category is indicated for all conditions and assay systems \parensubfig{B}.\nextentry
				\subfig{C} The average count of putative transcription factor binding sites (TFBS) in the test sequences of each category was normalized to the average count in inactive sequences and compared to the effect size of the binding site as determined in \cref{fig:TFBS}\subfigunformatted{A}.
			}%
			\label{efig:categories_TFs}%
		\end{efig}
		
		\begin{efig}
			\begin{tikzpicture}

	%%% scheme of the TFBS insertion experiment %%%
	\coordinate (scheme) at (0, 0);
	
	\node[anchor = north, font = \figsmall, text width = \fourcolumnwidth - 0.6666em, align = flush center, shift = {(.5\fourcolumnwidth, -.65)}] (seqs text) at (scheme) {50 random sequences without strong transcription factor binding sites};

	\node[anchor = north, transparent, node font = \figsmall\bfseries, inner sep = .15em, shift = {(-1, -.1)}] (TFBS 1) at (scheme -| seqs text) {TFBS};
	\node[anchor = south, transparent, node font = \figsmall\bfseries, inner sep = .15em, shift = {(1, 0)}] (TFBS 3) at (seqs text.north) {TFBS};
	\node[transparent, node font = \figsmall\bfseries, inner sep = .15em] (TFBS 2) at ($(TFBS 1)!.5!(TFBS 3)$) {TFBS};
	
	\begin{pgfonlayer}{background}
		\draw[thick] (scheme |- TFBS 1) ++(.5, 0) -- ++(\fourcolumnwidth - 2cm, 0);
		\draw[thick] (scheme |- TFBS 2) ++(1, 0) -- ++(\fourcolumnwidth - 2cm, 0);
		\draw[thick] (scheme |- TFBS 3) ++(1.5, 0) -- ++(\fourcolumnwidth - 2cm, 0);
	\end{pgfonlayer}
	
	\draw[very thick, -Latex] (seqs text.south) -- ++(0, -.45) coordinate (arrow end);
	
	\node[anchor = north, font = \figsmall, text width = \fourcolumnwidth - 0.6666em, align = flush center, yshift = -.65cm] (shuffle text) at (arrow end) {insert optimal transcription factor binding site in one of three positions};
	
	\node[anchor = north, fill = DodgerBlue1, text = white, node font = \figsmall\bfseries, inner sep = .15em, shift = {(-1, -.1)}] (TFBS 1) at (arrow end -| seqs text) {TFBS};
	\node[anchor = south, fill = DodgerBlue1, text = white, node font = \figsmall\bfseries, inner sep = .15em, shift = {(1, 0)}] (TFBS 3) at (shuffle text.north) {TFBS};
	\node[fill = DodgerBlue1, text = white, node font = \figsmall\bfseries, inner sep = .15em] (TFBS 2) at ($(TFBS 1)!.5!(TFBS 3)$) {TFBS};
	
	\begin{pgfonlayer}{background}
		\draw[thick] (scheme |- TFBS 1) ++(.5, 0) -- ++(\fourcolumnwidth - 2cm, 0);
		\draw[thick] (scheme |- TFBS 2) ++(1, 0) -- ++(\fourcolumnwidth - 2cm, 0);
		\draw[thick] (scheme |- TFBS 3) ++(1.5, 0) -- ++(\fourcolumnwidth - 2cm, 0);
	\end{pgfonlayer}
	
	\draw[very thick, -Latex] (shuffle text.south) -- ++(0, -.45) coordinate (arrow end);
	
	\node[anchor = north, font = \figsmall, text width = \fourcolumnwidth - 0.6666em, align = flush center, yshift = -.65cm] (measure text) at (arrow end) {measure enhancer strength of sequences with or without a binding site};
	
	\node[node font = \figsmall\bfseries, align = flush center, yshift = .325cm] (Plant STARR-seq) at (measure text.north) {Plant\\STARR-seq};
	
	\coordinate[xshift = -.3cm] (leaf center) at (Plant STARR-seq.west);
	\planticon*[width = .6cm]{tobaccoLeaf} at (leaf center);
	\begin{scope}[xscale = -1]
		\planticon[width = .323cm]{pipe} at ($(leaf center) + (.35, .16)$);
	\end{scope}
	
	\coordinate[xshift = .3cm] (protoplasts) at (Plant STARR-seq.east);
	\planticon[width = .3cm]{protoplast} at ($(protoplasts) + (90:.15)$);
	\planticon[width = .3cm]{protoplast} at ($(protoplasts) + (162:.15)$);
	\planticon[width = .3cm]{protoplast} at ($(protoplasts) + (234:.15)$);
	\planticon[width = .3cm]{protoplast} at ($(protoplasts) + (306:.15)$);
	\planticon[width = .3cm]{protoplast} at ($(protoplasts) + (18:.15)$);
	\planticon[width = .3cm, rotate = 15]{maize} at ($(protoplasts) + (.375, -.1)$);


	%%% comparison of TFBS insertion effect with ACR library %%%
	\coordinate (correlation) at (scheme -| \textwidth - \threequartercolumnwidth, 0);
	
	\begin{pggfplot}[%
		at = {(correlation)},
		total width = \threequartercolumnwidth,
		xlabel = {\textbf{transcription factor binding site insertion:} $\log_2$(effect size)},
		ylabel = {\textbf{ACR library:} $\log_2$(effect size)},
		title is csname,
		title color from name,
		xytick = {-2.1, -1.8, ..., 2.1}
	]{TFBS_insertion_tamsACR}
	
		\pggf_annotate(diagonal)
	
		\pggf_trendline()
	
		\pggf_scatter(
			color from column name,
			mark size = 1.5
		)
	
		\pggf_stats(
			stats = {n,spearman,rsquare}
		)
	
	\end{pggfplot}


	%%% effect of TFBS insertion %%%
	\coordinate[yshift = -\columnsep] (insertion effect) at (scheme |- xlabel.south);
	
	\tikzset{
		issignif/.code = {
			\pgfmathparse{#1 <= 0.05}
			\ifnum1=\pgfmathresult\relax
				\pgfkeysalso{fill = OkabeIto4}
			\else
				\pgfkeysalso{fill = gray}
			\fi
		}
	}

	\begin{pggfplot}[%
		at = {(insertion effect)},
		total width = \textwidth,
		xlabel = transcription factor family (motif ID),
		ylabel = {\enrichment[rel. to no TFBS]},
		ylabel add to width = .5\baselineskip,
		xticklabel style = {anchor = east, rotate = 90, inner ysep = 0pt, inner xsep = .3333em},
		xticklabels are csnames,
		title is csname,
		title color from name,
		enlarge x limits = 0,
		enlarge y limits = {lower = 0.1},
		ytick = {-10, -9, ..., 10},
		legend style = {anchor = north east, at = {(1, 1)}},
		legend image width = .75\baselineskip,
		every legend image post/.append style = {fill opacity = 0.5},
		tight y
	]{TFBS_insertion_effect}
	
		\pggf_annotate(zeroline)
		
		\pggf_annotate(
			manual legend = {
				significant = {boxplot legend = y, fill = OkabeIto4},
				not significant = {boxplot legend = y, fill = gray}
			},
			only facet = 1
		)
	
		\pggf_boxplot(
			style from column = {issignif = p_value},
			sample size = {font = \fontsize{4}{5}\selectfont},
			forget plot
		)
	
	\end{pggfplot}
	
	
	%%% subfigure labels %%%
	\subfiglabel{scheme}
	\subfiglabel[yshift = .5\baselineskip]{insertion effect}
	\subfiglabel{correlation}

\end{tikzpicture}%
			\caption{%
				\captiontitle{Inserted transcription factor binding sites can increase enhancer strength}%
				\subfigrange{A}{C} A transcription factor binding site (TFBS) was inserted into 50 random sequences and the sequences were subjected to Plant STARR-seq \parensubfig{A} to determine their enhancer strength relative to the corresponding random sequence without a binding site (no TFBS) \parensubfig{B}. The effect size of each binding site was determined as the fold-change between the sequences with and without it, and was compared to the effect size in the  ACR library (see \cref{fig:TFBS}\subfigunformatted{B}) \parensubfig{C}. The solid and dashed lines represent a linear regression line and a $y = x$ line, respectively. Pearson's $R^2$, Spearman's $\rho$, and the number ($n$) of samples are indicated.\nextentry
				Box plots in \parensubfig{B} are as defined in \cref{fig:enhScreen}. Significant differences from a null distribution were determined using one-sample Wilcoxon rank-sum tests. Box plots for significant groups (Bonferroni-adjusted $p$ value \textle{} 0.05) are shown in green. Non-significant groups are represented by gray box plots. Exact $p$ values are listed in Supplementary Data 4.% TODO: make sure reference is correct
			}%
			\label{efig:TFBS_insertion}
		\end{efig}
		
		\begin{efig}
			\begin{tikzpicture}
	
	%%% prediction accuracy of synthetic enhancer strength with different inputs %%%
	%% explain input data
	\coordinate (scheme) at (0, 0);
	
	\foreach \x in {1, 2, 3}{
		\node[anchor = west, node font = \figsmall, yshift = -\x * .75cm] (seq\x) at (scheme) {seq\x};
	}
	
	\distance{seq3.east}{\threecolumnwidth, 0}
	\pgfmathsetlength{\templength}{(\xdistance - 2\columnsep) / 4}
	
	\pgfmathsetseed{928}
	
	\colorlet{seq1}{gray!75!white}
	\colorlet{seq2}{gray}
	\colorlet{seq3}{gray!75!black}
	
	% sequence- and position-specific effect
	\node[anchor = north west, node font = \figsmall\bfseries, text width = \templength, align = center, inner xsep = 0pt, shift = {(.5\columnsep, -.5\columnsep)}] (seq+pos) at (scheme -| seq3.east) {seq+pos};
	
	\foreach \x in {1, 2, 3}{
		\draw[ultra thick, seq\x] (seq\x -| seq+pos.west) -- ++(\templength, 0) coordinate[pos = .165] (p1) coordinate[pos = .5] (p2) coordinate[pos = .835] (p3);
		\foreach \y in {1, 2, 3}{
			\draw[->] (p\y) +(0, .05) -- ++(0, -.35) node[anchor = north, node font = \figtiny] {%
				\pgfmathparse{rnd}%
				\expandafter\xdef\csname tmp\x\y\endcsname{\pgfmathresult}%
				\pgfmathprintnumber[fixed, zerofill, precision = 1]{\csname tmp\x\y\endcsname}%
			};
		}
	}
	
	% sequence-specific effect
	\node[anchor = base west, node font = \figsmall\bfseries, text width = \templength, align = center, inner xsep = 0pt, xshift = .5\columnsep] (seq only) at (seq+pos.base east) {seq only};
	
	\foreach \x in {1, 2, 3}{
		\draw[ultra thick, seq\x] (seq\x -| seq only.west) -- ++(\templength, 0) coordinate[pos = .165] (p1) coordinate[pos = .5] (p2) coordinate[pos = .835] (p3);
		\coordinate[yshift = -.35cm] (e\x) at (p2);
		\foreach \y in {1, 2, 3}{
			\draw[->] (p\y) +(0, .05) -- ++(0, -.05) to[out = -90, in = 90] ($(e\x) + (0, .05)$) -- (e\x);
		}
		\node[anchor = north, node font = \figtiny] at (e\x) {%
			\pgfmathparse{(\csname tmp\x1\endcsname + \csname tmp\x2\endcsname + \csname tmp\x3\endcsname) / 3}%
			\pgfmathprintnumber[fixed, zerofill, precision = 1]{\pgfmathresult}%
		};
	}
	
	% position-specific effect
	\node[anchor = base west, node font = \figsmall\bfseries, text width = \templength, align = center, inner xsep = 0pt, xshift = .5\columnsep] (pos only) at (seq only.base east) {pos only};
	
	\foreach \x in {1, 2, 3}{
		\draw[ultra thick, seq\x] (seq\x -| pos only.west) -- ++(\templength, 0) coordinate[pos = .165] (p1) coordinate[pos = .5] (p2) coordinate[pos = .835] (p3);
	}
	\foreach \y in {1, 2, 3}{
		\draw[->] (seq1 -| p\y) +(0, .05) -- (e3 -| p\y) node[anchor = north, node font = \figtiny] {%
			\pgfmathparse{(\csname tmp1\y\endcsname + \csname tmp2\y\endcsname + \csname tmp3\y\endcsname) / 3}%
			\pgfmathprintnumber[fixed, zerofill, precision = 1]{\pgfmathresult}%
		};
	}
	
	% mean effect
	\node[anchor = base west, node font = \figsmall\bfseries, text width = \templength, align = center, inner xsep = 0pt, xshift = .5\columnsep] (mean) at (pos only.base east) {mean};
	
	\foreach \x in {1, 2, 3}{
		\draw[ultra thick, seq\x] (seq\x -| mean.west) -- ++(\templength, 0) coordinate[pos = .165] (p1) coordinate[pos = .5] (p2) coordinate[pos = .835] (p3);
	}
	\foreach \y in {1, 2, 3}{
		\draw[->] (seq1 -| p\y) +(0, .05) -- (seq3 -| p\y) -- ++(0, -.05) to[out = -90, in = 90] ($(p2 |- e3) + (0, .05)$) -- (p2 |- e3);
	}
	\node[anchor = north, node font = \figtiny] (effect) at (p2 |- e3) {%
		\pgfmathparse{%
			(%
				\csname tmp11\endcsname + \csname tmp12\endcsname + \csname tmp13\endcsname +%
				\csname tmp21\endcsname + \csname tmp22\endcsname + \csname tmp23\endcsname +%
				\csname tmp31\endcsname + \csname tmp32\endcsname + \csname tmp33\endcsname
			) / 9%
		}%
		\pgfmathprintnumber[fixed, zerofill, precision = 1]{\pgfmathresult}%
	};
	
	% ACR lib
	\node[anchor = north east, node font = \figsmall\bfseries, text width = \templength, align = right, yshift = -\columnsep] (ACR lib) at (seq+pos.east |- effect.south) {ACR lib:};
	\node[anchor = base west, node font = \figsmall, text width = 3\templength, align = left] at (ACR lib.base east) {transcription factor effects from ACR library};
	
	% random
	\node[anchor = north east, node font = \figsmall\bfseries, text width = \templength, align = right, yshift = -3\columnsep] (random) at (seq+pos.east |- effect.south) {random:};
	\node[anchor = base west, node font = \figsmall] at (random.base east) {random permutation of \textbf{seq+pos}};
	
	
	%% plot prediction accuracy
	\coordinate (accuracy) at (scheme -| \textwidth - \twothirdcolumnwidth, 0);
	
	\expandafter\def\csname id+pos\endcsname{\vphantom{ly}seq+pos}
	\expandafter\def\csname no_pos\endcsname{\vphantom{ly}seq only}
	\expandafter\def\csname no_id\endcsname{\vphantom{ly}pos only}
	\expandafter\def\csname mean\endcsname{\vphantom{ly}mean}
	\expandafter\def\csname tamsACR\endcsname{\vphantom{ly}ACR lib}
	\expandafter\def\csname random\endcsname{\vphantom{ly}random}
	
	\begin{pggfplot}[%
		at = {(accuracy)},
		total width = \twothirdcolumnwidth,
		xlabel = method to obtain transcription factor effects,
		ylabel = prediction accuracy ($R^2$),
		ymin = 0,
		enlarge x limits = 0,
		enlarge y limits = {lower = 0},
		title is csname,
		title color from name,
		xticklabels are csnames,
		xticklabel style = {rotate = 45, anchor = north east, inner sep = 0.236em, text depth = -2pt},
		tight y
	]{synEnh_pred_accuracy}
	
		\pggf_bar(
			ybar,
			color from column name,
			draw = black
		)
	
	\end{pggfplot}
	
	
	%%% predict synthetic enhancer strength %%%
	\coordinate[yshift = -\columnsep] (correlation) at (scheme |- xlabel.south);
	
	\begin{pggfplot}[%
		at = {(correlation)},
		total width = \textwidth,
		xlabel = {\textbf{measurement}: \enrichment},
		ylabel = {\textbf{prediction (seq+pos)}:\\\enrichment},
		ylabel add to width = 3pt,
		facet 1/.append style = {
			pggf colorbar = {%
				title = count,
				log2,
				horizontal,
				ytick = {0, 2, ..., 15}
			}
		},
		title is csname,
		title color from name,
		xytick = {-12, -10, ..., 12},
	]{synEnh_prediction}
	
		\pggf_annotate(diagonal)
	
		\pggf_hexbin(bins = 50)
	
		\pggf_stats(stats = {n,spearman,rsquare})
	
	\end{pggfplot}
	
	
	%%% subfigure labels %%%
	\subfiglabel{scheme}
	\subfiglabel{accuracy}
	\subfiglabel{correlation}

\end{tikzpicture}%
			\caption{%
				\captiontitle{The strength of synthetic enhancers can be predicted}%
				\subfigrange{A}{C} The strength of synthetic enhancers with two or three transcription factor binding sites was predicted based on the strength of the random sequence without any binding sites and the effects observed for inserting individual binding sites. The transcription factor binding site effects were calculated individually for each random sequence and each position (seq+pos), as the mean across all positions for each sequence (seq only), as the mean across all sequences for each position (pos only), or as the mean across all sequences and positions (mean). Additionally, transcription factor binding site effects were obtained from the ACR library experiment (ACR lib; see \cref{fig:TFBS}\subfigunformatted{A}) or from a random permutation of the seq+pos method \parensubfig{A}. For each method of obtaining transcription factor effects, the prediction accuracy was determined as the correlation between the predicted and experimentally determined synthetic enhancer strengths \parensubfig{B}. Hexbin plots (as defined in \cref{fig:enhScreen}) of this correlation are shown for the seq+pos method \parensubfig{C}.
			}%
			\label{efig:synEnh_pred_accuracy}%
		\end{efig}
		
		\begin{efig}
			\begin{tikzpicture}

	\pgfplotsset{letter colors = DNA}
	
	%%% scheme of synthetic enhancer design %%%
	\coordinate (scheme) at (0, 0);
	
	\pgfmathsetlength{\templength}{.45cm}
	
	\coordinate[yshift = -\columnsep - 2\templength] (TFs) at (scheme);
	
	% nucleotide composition
	\coordinate[shift = {(.5\columnsep + \templength, 1.5\templength)}] (low GC) at (TFs);
	\coordinate[shift = {(.5\columnsep + \templength, -1.5\templength)}] (high GC) at (TFs);
	
	\pgfplotstableread[col sep = tab]{rawData/synEnh_nucFreq.tsv}\nucFreqTable
	\pgfplotstableforeachcolumn\nucFreqTable\as\GCbin{
		\pgfplotstablegetelem{0}{\GCbin}\of\nucFreqTable
		\foreach [expand list, evaluate = \y as \x using 180 - \y, remember = \x as \lastx (initially 180)] \base/\y in {\pgfplotsretval}{
			\fill[letter\base col!50] (\GCbin) -- ++(\lastx:\templength) arc[start angle = \lastx, end angle = \x, radius = \templength] coordinate[pos = .5] (c\base) -- cycle;
			\node[font = \figsmall] at ($(\GCbin)!.6!(c\base)$) {\base};
		}
	}
	
	% sequences
	\coordinate[xshift = .5\columnsep + \templength] (seqs) at (TFs -| low GC);
	\coordinate (low GC seq) at (low GC -| seqs);
	\coordinate (high GC seq) at (high GC -| seqs);
	
	\distance{seqs}{\threecolumnwidth - .5\columnsep, 0}
	
	\foreach \x/\y in {low GC/south, high GC/north}{
		\draw[line width = .06cm, \x] (\x\space seq) node[anchor = \y\space west, text = black, text depth = 0pt, xshift = -.5\columnsep] {25 sequences with \textbf{\x} content} -- ++(\xdistance, 0) coordinate[pos = .165] (\x\space 1) coordinate[pos = .5] (\x\space 2) coordinate[pos = .835] (\x\space 3) coordinate (\x\space seq end);
	}
	
	% TFs
	\foreach \x in {1, 2, 3}{
		\node[fill = DodgerBlue1, text = white, node font = \figsmall\bfseries, inner sep = .15em] (TFBS \x) at (TFs -| low GC \x) {TFBS};
		
		\draw[-Latex, shorten > = .08cm, DodgerBlue1!50] (TFBS \x.north) ++(0, .05) -- (low GC \x);
		\draw[-Latex, shorten > = .08cm, DodgerBlue1!50] (TFBS \x.south) ++(0, -.05) -- (high GC \x);
	}
	
	\node[anchor = base, node font = \figsmall] at ($(TFBS 1.base)!.5!(TFBS 2.base)$) {and};
	\node[anchor = base, node font = \figsmall] at ($(TFBS 2.base)!.5!(TFBS 3.base)$) {and};

	% tested TFBSs
	\coordinate[yshift = -.5\columnsep - \templength] (TFBSs) at (scheme |- high GC);
	
	\node[anchor = north, font = \bfseries] (title) at ($(TFBSs)!.5!(TFBSs -| low GC seq end)$) {transcription factor binding sites (TFBS)\settowidth{\global\templength}{\figsmall tobacco-specific:}};

	\pgfplotstableread[col sep = tab]{rawData/synEnh_TF_specificities.tsv}\synEnhTFtable

	\pgfplotstableforeachcolumnelement{specificity}\of\synEnhTFtable\as\x{
		\pgfmathsetmacro{\xi}{int(\pgfplotstablerow + 1)}
		\node[anchor = base west, node font = \figsmall, yshift = -\xi*1.15\baselineskip, text width = \templength, align = right] at (title.base west) {\x:};
	}
	\pgfplotstableforeachcolumnelement{TF}\of\synEnhTFtable\as\x{
		\pgfmathsetmacro{\xi}{int(\pgfplotstablerow + 1)}
		\node[anchor = base east, node font = \figsmall, yshift = -\xi*1.15\baselineskip] (TFBS name \xi) at (title.base east) {\usename{\x}};
	}

	% connect TFBS list with TFBSs
	\draw[decorate, decoration = {brace, amplitude = .3cm}] ($(TFBS name 1.north east) + (-.1, 0)$) -- ($(TFBS name 6.south east) + (-.1, 0)$) coordinate[pos = .5, xshift = .3cm - \pgflinewidth] (curly);
	\draw[thick, -Latex, shorten > = .05cm, rounded corners = .15cm] (curly) -- (curly -| \threecolumnwidth - .5\pgflinewidth, 0) |- (TFBS 3.east);
	
	
	%%% species-/condition-specificity of synthetic enhancers %%%
	\coordinate (specificity) at (scheme -| \textwidth - \twothirdcolumnwidth, 0);

	\tikzset{
		species-specificity/.style = {left color = tobacco!20, right color = maize!20, middle color = white},
		light-specificity/.style = {left color = light!20, right color = dark!20, middle color = white},
		temperature-specificity/.style = {left color = warm!20, right color = cold!20, middle color = white}
	}
	
	\expandafter\def\csname species-specificity\endcsname{species specificity}
	\expandafter\def\csname light-specificity\endcsname{light specificity}
	\expandafter\def\csname temperature-specificity\endcsname{temperature specificity}
	
	\def\noTFBS{\vphantom{tblih}none}
	\expandafter\def\csname tobacco-specific\endcsname{\vphantom{tblih}tobacco-\\specific}
	\expandafter\def\csname maize-specific\endcsname{\vphantom{tblih}maize-\\specific}
	\expandafter\def\csname light-specific\endcsname{\vphantom{tblih}light-\\specific}
	\expandafter\def\csname dark-specific\endcsname{\vphantom{tblih}dark-\\specific}
	\expandafter\def\csname warm-specific\endcsname{\vphantom{tblih}warm-\\specific}
	\expandafter\def\csname cold-specific\endcsname{\vphantom{tblih}cold-\\specific}
	
	\begin{pggfplot}[%
		at = {(specificity)},
		total width = \twothirdcolumnwidth,
		xlabel = inserted transcription factor binding site,
		ylabel = $\log_2$(specificity),
		enlarge x limits = {abs = 1, rel = 0.005},
		enlarge y limits = {lower = 0.1, upper = 0.1},
		title is csname,
		title style from name,
		ytick = {-12, -10, ..., 12},
		xticklabel style = {align = center},
		xticklabels are csnames,
		tight y
	]{synEnh_specificity}
	
		\pggf_annotate(zeroline)
		
		\pggf_annotate(
			annotation = {
				\draw[thick, -Latex] (-.25, .25) -- (-.25, 4.5) node[pos = .45, anchor = north, rotate = 90, align = center, font = \figsmall] {more active\\in \textbf{\usename{tobacco}}};
				\draw[thick, -Latex] (4.25, -.25) -- (4.25, -3.5) node[pos = .45, anchor = north, rotate = -90, align = center, font = \figsmall] {more active\\in \textbf{\usename{maize}}};
			},
			only facet = 1
		)
		
		\pggf_annotate(
			annotation = {
				\draw[thick, -Latex] (-.25, .5) -- (-.25, 4.5) node[pos = .45, anchor = north, rotate = 90, align = center, font = \figsmall] {more active\\in \textbf{\usename{light}}};
				\draw[thick, -Latex] (4.25, -.5) -- (4.25, -3.5) node[pos = .45, anchor = north, rotate = -90, align = center, font = \figsmall] {more active\\in \textbf{\usename{dark}}};
			},
			only facet = 2
		)
		
		\pggf_annotate(
			annotation = {
				\draw[thick, -Latex] (-.25, .5) -- (-.25, 4.5) node[pos = .45, anchor = north, rotate = 90, align = center, font = \figsmall] {more active\\in \textbf{\usename{warm}}};
				\draw[thick, -Latex] (4.25, -.5) -- (4.25, -3.5) node[pos = .45, anchor = north, rotate = -90, align = center, font = \figsmall] {more active\\in \textbf{\usename{cold}}};
			},
			only facet = 3
		)
	
		\pggf_boxplot(
			style from column = {fill = target_species},
			cld = {anchor = south, text depth = 0pt}{max}
		)
	
	\end{pggfplot}
	

	%%% strength of species /condition-specific synthetic enhancers %%%
	\coordinate[yshift = -\columnsep] (strength) at (scheme |- xlabel.south);
	
	\colorlet{high}{OkabeIto4}
	\colorlet{any}{gray}
	\colorlet{low}{OkabeIto2}
	
	\expandafter\def\csname tobacco-specific\endcsname{\vphantom{tblih}tobacco-specific}
	\expandafter\def\csname maize-specific\endcsname{\vphantom{tblih}maize-specific}
	\expandafter\def\csname light-specific\endcsname{\vphantom{tblih}light-specific}
	\expandafter\def\csname dark-specific\endcsname{\vphantom{tblih}dark-specific}
	\expandafter\def\csname warm-specific\endcsname{\vphantom{tblih}warm-specific}
	\expandafter\def\csname cold-specific\endcsname{\vphantom{tblih}cold-specific}
	
	\begin{pggfplot}[%
		at = {(strength)},
		total width = \textwidth,
		xlabel = inserted transcription factor binding site,
		ylabel = \enrichment,
		enlarge x limits = 0,
		enlarge y limits = {lower = 0.1, upper = 0.15},
		title is csname,
		title color from name,
		xticklabel style = {rotate = 45, anchor = north east, inner sep = 0.236em},
		xticklabels are csnames,
		legend style = {anchor = north east, at = {(1, 1)}},
		legend columns = 2,
		legend image width = .75\baselineskip,
		every legend image post/.append style = {fill opacity = 0.5},
		tight y
	]{synEnh_specific_strength}
	
		\pggf_annotate(
			manual legend = {
				{\textbf{expected activity:}\hspace*{-\groupplotsep}} = {empty legend},
				high = {boxplot legend = y, fill = high},
				\relax = {empty legend},
				any = {boxplot legend = y, fill = any},
				\relax = {empty legend},
				low = {boxplot legend = y, fill = low}
			},
			only facet = 1
		)
	
		\pggf_hline(densely dotted)
	
		\pggf_boxplot(
			style from column = {fill = target},
			cld = {anchor = south, text depth = 0pt}{max},
			forget plot
		)
	
	\end{pggfplot}
	
	
	%%% subfigure labels %%%
	\subfiglabel{scheme}
	\subfiglabel{strength}
	\subfiglabel{specificity}

\end{tikzpicture}%
			\caption{%
				\captiontitle{Rationally designed synthetic enhancers show little condition specificity}%
				\subfigrange{A}{C} Synthetic enhancers were designed by inserting three binding sites for transcription factors with species- or condition-specific effects into 25 random sequences each with a nucleotide composition similar to an average dicot (low GC) or monocot (high GC) ACR \parensubfig{A}. The strength \parensubfig{B} and species or condition specificity \parensubfig{C} of these synthetic enhancers was determined by Plant STARR-seq in the indicated conditions. Species and condition specificity was calculated as the fold-change in enhancer strength between two species (tobacco [mean of light, dark, warm and cold] and maize) or conditions (light and dark for light specificity; warm and cold for temperature specificity). Box plots are as defined in \cref{fig:enhScreen} and letters above the box plots in indicate significance groups determined by post-hoc Tukey tests performed separately for each condition, species, or specificity. Exact $p$ values are listed in Supplementary Data 4. In \parensubfig{B}, box plots are colored according to the expected activity in the indicated condition.% TODO: make sure reference is correct
			}%
			\label{efig:synEnh_specificity}%
		\end{efig}
		
		\begin{efig}
			\begin{tikzpicture}

	%%% linear model for TF hits %%%
	\coordinate (lm hits) at (0, 0);
	
	% sequence
	\node[anchor = west, node font = \figsmall\courier, shift = {(0, -1.25)}] (seq) at (lm hits) {%
		\pgfmathsetseed{928}\randomNucs{10}\,\raisebox{.5ex}{\dots}\,\randomNucs{10}%
	};
	
	\draw[very thick, -Latex] (seq.east) -- ++(.5, 0) coordinate (ohe);
	
	% TF hits
	\node[anchor = west, node font = \figsmall] (hits) at (ohe) {
		\setlength{\tabcolsep}{.5em}%
		\begin{tabular}{rrrcrrr}%
			\toprule%
			\textbf{TF-1} & \textbf{TF-2} & \textbf{TF-3} & \raisebox{.5ex}{\textbf{\dots}} & \textbf{TF-70} & \textbf{TF-71} & \textbf{TF-72} \\%
			\midrule%
			\pgfmathparse{int(random(4) - 1)}\pgfmathresult &
			\pgfmathparse{int(random(4) - 1)}\pgfmathresult &
			\pgfmathparse{int(random(4) - 1)}\pgfmathresult &
			\raisebox{.5ex}{\dots} &
			\pgfmathparse{int(random(4) - 1)}\pgfmathresult &
			\pgfmathparse{int(random(4) - 1)}\pgfmathresult &
			\pgfmathparse{int(random(4) - 1)}\pgfmathresult \\%
			\bottomrule%
		\end{tabular}%
	};
	
	\draw[very thick, -Latex] (hits.east) -- ++(.5, 0) coordinate[xshift = .3333em] (model);
	
	% model
	\node[right = .05cm of model, align = flush right, font = \figsmall, yshift = -.5\baselineskip] (lm) {%
		$\log_2(\text{enhancer strength}) =$\hspace*{7em}\\
		$\beta_0 + \beta_1\text{TF-1} + \beta_2\text{TF-2} + \beta_3\text{TF-3} + \dots$\\
		$+ \beta_{70}\text{TF-70} + \beta_{71}\text{TF-71} + \beta_{72}\text{TF-72}$%
	};
	
	\begin{pgfonlayer}{background}
		\draw[fill = OkabeIto5!25, rounded corners] (model) |- ($(lm.north east) + (.05, .05cm + \baselineskip)$) node[pos = .75, anchor = north, fill = none, draw = none, text = black, font = \bfseries] (model top) {linear regression} |- ($(lm.south west) + (-.05, -.05)$) coordinate[pos = .25] (model out) coordinate[pos = .75] (model bottom) -- cycle;
	\end{pgfonlayer}
	
	\draw[very thick, -Latex] (model out) ++(.3333em, 0) -- ++(.5, 0) coordinate (predictions);
	
	% predictions
	\pgfmathsetmacro{\predLight}{2 + rand * 4}
	\pgfmathsetmacro{\predDark}{\predLight + rand * 0.5 * \predLight}
	\pgfmathsetmacro{\predWarm}{\predLight + rand * 0.1 * \predLight}
	\pgfmathsetmacro{\predCold}{\predLight + rand * 0.1 * \predLight}
	\pgfmathsetmacro{\predMaize}{1.5 + rand * 3}
	
	\node[anchor = west, node font = \figsmall] (predictions) at (predictions) {%
		\pgfmathsetseed{928}%
		\setlength{\tabcolsep}{.5em}%
		\begin{tabular}{rrrrr}%
			\toprule%
			\textbf{\usename{light}} & \textbf{\usename{dark}} & \textbf{\usename{warm}} & \textbf{\usename{cold}} & \textbf{\usename{maize}} \\%
			\midrule%
			\pgfmathprintnumber[fixed, zerofill, precision = 2]{\predLight} &
			\pgfmathprintnumber[fixed, zerofill, precision = 2]{\predDark} &
			\pgfmathprintnumber[fixed, zerofill, precision = 2]{\predWarm} &
			\pgfmathprintnumber[fixed, zerofill, precision = 2]{\predCold} &
			\pgfmathprintnumber[fixed, zerofill, precision = 2]{\predMaize} \\%
			\bottomrule%
		\end{tabular}%
	};
	
	% text
	\node[anchor = south, font = \strut, yshift = .3em] at (seq |- model top.north) {170\,bp sequence};
	\node[anchor = south, font = \strut, yshift = .3em] at (hits |- model top.north) {count occurence of transcription factor binding sites};
	\node[anchor = south, font = \strut, yshift = .3em] at (model top.north) {linear model};
	\node[anchor = south, font = \strut, yshift = .3em] (blub) at (predictions |- model top.north) {predict $\log_2$(enhancer strength)};
	

	%%% enhancer strength predictions from linear model based on TF hits %%%
	\coordinate[yshift = -\columnsep] (cor hits) at (lm hits |- model bottom);
	
	\begin{pggfplot}[%
		at = {(cor hits)},
		total width = \textwidth,
		xlabel = {\textbf{measurement}: \enrichment},
		ylabel = {\textbf{prediction}:\\\enrichment},
		facet 1/.append style = {
			pggf colorbar = {%
				title = count,
				log2,
				horizontal,
				ytick = {0, 3, ..., 15}
			}
		},
		title is csname,
		title color from name,
		xytick = {-12, -10, ..., 12},
		ylabel add to width = .6\baselineskip,
	]{lm_TF_hits_predictions}
	
		\pggf_annotate(diagonal)
	
		\pggf_hexbin(bins = 50)
	
		\pggf_stats(stats = {n,spearman,rsquare})
	
	\end{pggfplot}
	
	
	%%% linear model for TF scores %%%
	\coordinate[yshift = -\columnsep] (lm scores) at (lm hits |- xlabel.south);
	
	% sequence
	\node[anchor = west, node font = \figsmall\courier, shift = {(0, -1.25)}] (seq) at (lm scores) {%
		\pgfmathsetseed{928}\randomNucs{10}\,\raisebox{.5ex}{\dots}\,\randomNucs{10}%
	};
	
	\draw[very thick, -Latex] (seq.east) -- ++(.5, 0) coordinate (ohe);
	
	% TF hits
	\node[anchor = west, node font = \figsmall] (hits) at (ohe) {
		\pgfmathsetseed{928}%
		\setlength{\tabcolsep}{.5em}%
		\begin{tabular}{rrrcrrr}%
			\toprule%
			\textbf{TF-1} & \textbf{TF-2} & \textbf{TF-3} & \raisebox{.5ex}{\textbf{\dots}} & \textbf{TF-70} & \textbf{TF-71} & \textbf{TF-72} \\%
			\midrule%
			\pgfmathparse{rnd * random(3)}\pgfmathprintnumber[fixed, zerofill, precision = 2]{\pgfmathresult} &
			\pgfmathparse{rnd * random(3)}\pgfmathprintnumber[fixed, zerofill, precision = 2]{\pgfmathresult} &
			\pgfmathparse{rnd * random(3)}\pgfmathprintnumber[fixed, zerofill, precision = 2]{\pgfmathresult} &
			\raisebox{.5ex}{\dots} &
			\pgfmathparse{rnd * random(3)}\pgfmathprintnumber[fixed, zerofill, precision = 2]{\pgfmathresult} &
			\pgfmathparse{rnd * random(3)}\pgfmathprintnumber[fixed, zerofill, precision = 2]{\pgfmathresult} &
			\pgfmathparse{rnd * random(3)}\pgfmathprintnumber[fixed, zerofill, precision = 2]{\pgfmathresult} \\%
			\bottomrule%
		\end{tabular}%
	};
	
	\draw[very thick, -Latex] (hits.east) -- ++(.5, 0) coordinate[xshift = .3333em] (model);
	
	% model
	\node[right = .05cm of model, align = flush right, font = \figsmall, yshift = -.5\baselineskip] (lm) {%
		$\log_2(\text{enhancer strength}) =$\hspace*{7em}\\
		$\beta_0 + \beta_1\text{TF-1} + \beta_2\text{TF-2} + \beta_3\text{TF-3} + \dots$\\
		$+ \beta_{70}\text{TF-70} + \beta_{71}\text{TF-71} + \beta_{72}\text{TF-72}$%
	};
	
	\begin{pgfonlayer}{background}
		\draw[fill = OkabeIto5!25, rounded corners] (model) |- ($(lm.north east) + (.05, .05cm + \baselineskip)$) node[pos = .75, anchor = north, fill = none, draw = none, text = black, font = \bfseries] (model top) {linear regression} |- ($(lm.south west) + (-.05, -.05)$) coordinate[pos = .25] (model out) coordinate[pos = .75] (model bottom) -- cycle;
	\end{pgfonlayer}
	
	\draw[very thick, -Latex] (model out) ++(.3333em, 0) -- ++(.5, 0) coordinate (predictions);
	
	% predictions
	\pgfmathsetmacro{\predLight}{2 + rand * 4}
	\pgfmathsetmacro{\predDark}{\predLight + rand * 0.5 * \predLight}
	\pgfmathsetmacro{\predWarm}{\predLight + rand * 0.1 * \predLight}
	\pgfmathsetmacro{\predCold}{\predLight + rand * 0.1 * \predLight}
	\pgfmathsetmacro{\predMaize}{1.5 + rand * 3}
	
	\node[anchor = west, node font = \figsmall] (predictions) at (predictions) {%
		\pgfmathsetseed{928}%
		\setlength{\tabcolsep}{.5em}%
		\begin{tabular}{rrrrr}%
			\toprule%
			\textbf{\usename{light}} & \textbf{\usename{dark}} & \textbf{\usename{warm}} & \textbf{\usename{cold}} & \textbf{\usename{maize}} \\%
			\midrule%
			\pgfmathprintnumber[fixed, zerofill, precision = 2]{\predLight} &
			\pgfmathprintnumber[fixed, zerofill, precision = 2]{\predDark} &
			\pgfmathprintnumber[fixed, zerofill, precision = 2]{\predWarm} &
			\pgfmathprintnumber[fixed, zerofill, precision = 2]{\predCold} &
			\pgfmathprintnumber[fixed, zerofill, precision = 2]{\predMaize} \\%
			\bottomrule%
		\end{tabular}%
	};
	
	% text
	\node[anchor = south, font = \strut, yshift = .3em] at (seq |- model top.north) {170\,bp sequence};
	\node[anchor = south, font = \strut, yshift = .3em] at (hits |- model top.north) {sum of transcription factor motif scores};
	\node[anchor = south, font = \strut, yshift = .3em] at (model top.north) {linear model};
	\node[anchor = south, font = \strut, yshift = .3em] (blub) at (predictions |- model top.north) {predict $\log_2$(enhancer strength)};
	
	
	%%% enhancer strength predictions from linear model based on TF scores %%%
	\coordinate[yshift = -\columnsep] (cor scores) at (cor hits |- model bottom);
	
	\begin{pggfplot}[%
		at = {(cor scores)},
		total width = \textwidth,
		xlabel = {\textbf{measurement}: \enrichment},
		ylabel = {\textbf{prediction}:\\\enrichment},
		facet 1/.append style = {
			pggf colorbar = {%
				title = count,
				log2,
				horizontal,
				ytick = {0, 3, ..., 15}
			}
		},
		title is csname,
		title color from name,
		xytick = {-12, -10, ..., 12},
		ylabel add to width = .6\baselineskip,
	]{lm_TF_scores_predictions}
	
		\pggf_annotate(diagonal)
	
		\pggf_hexbin(bins = 50)
	
		\pggf_stats(stats = {n,spearman,rsquare})
	
	\end{pggfplot}
	
	
	%%% compare linear model coefficients %%%
	\coordinate[yshift = -\columnsep] (coefficients) at (cor hits |- xlabel.south);
	
	\makeatletter
	
	\tikzset{
		label TF/.code 2 args = {
			\edef\tmpA{nan}
			\edef\tmpB{#1}
			\ifx\tmpA\tmpB\relax\else
				\pgfkeysalso{pin = {[node font = \figsmall, black, inner sep = .1em]#1:\pgfmathprintnumber[precision = 0]{#2}}}
			\fi
		}
	}
	
	\makeatother
	
	\begin{pggfplot}[%
		at = {(coefficients)},
		total width = \textwidth,
		xlabel = \textbf{TFBS counts:} model coefficient,
		ylabel = \textbf{TFBS scores:} model coefficient,
		title is csname,
		title color from name,
		xytick = {-1.2, -.9, ..., 1.2},
		ylabel add to width = .6\baselineskip,
	]{lm_TF_coefficients}
	
		\pggf_annotate(diagonal)
	
		\pggf_scatter(
			color from column name,
			nodes near coords = {},
			nodes near coords align = {anchor = center},
			nodes near coords style = {inner sep = 1.5pt, text = black, label TF = {\pin}{\TFid}},
			pin distance = .5em,
			visualization depends on = \thisrow{TF} \as \TFid,
			visualization depends on = \thisrow{pin} \as \pin
		)
	
		\pggf_stats(stats = {n,spearman,rsquare})
	
	\end{pggfplot}
	
	
	%%% subfigure labels %%%
	\subfiglabel{lm hits}
	\subfiglabel{cor hits}
	\subfiglabel{lm scores}
	\subfiglabel{cor scores}
	\subfiglabel{coefficients}

\end{tikzpicture}%
			\caption{%
				\captiontitle{Linear models based on transcription factor binding sites predict enhancer strength with low accuracy}%
				\subfigrange{A}{D} Linear models based on the counts of putative transcription factor binding sites \parensubfigtwo{A}{B} or on the cumulative scores of binding motif matches within each sequence \parensubfigtwo{C}{D} were trained to predict enhancer strength. The predictions of the models were compared with the experimentally determined enhancer strengths in a held-out test set comprising 10\% of the test sequences \parensubfigtwo{B}{D}.\nextentry
				\subfig{E} Correlation between the coefficients assigned by the linear models in \parensubfigrange{A}{D} to the 72 different transcription factor binding motifs. Selected motifs associated with increased or decreased enhancer strength are labeled.\nextentry
				Hexbin plots in \parensubfig{B} and \parensubfig{D} are as defined in \cref{fig:enhScreen}. In \parensubfig{B}, \parensubfig{C}, and \parensubfig{E}, the dashed lines represents a $y = x$ line. Pearson's $R^2$, Spearman's $\rho$, and the number ($n$) of samples are indicated.
				\label{efig:lm_TFs}%
			}
		\end{efig}
		
		\begin{efig}
			\begin{tikzpicture}

	%%% linear model for TF hits %%%
	\coordinate (lm kmers) at (0, 0);
	
	% sequence
	\node[anchor = west, node font = \figsmall\courier, shift = {(0, -1.25)}] (seq) at (lm kmers) {%
		\pgfmathsetseed{928}\randomNucs{10}\,\raisebox{.5ex}{\dots}\,\randomNucs{10}%
	};
	
	\draw[very thick, -Latex] (seq.east) -- ++(.5, 0) coordinate (ohe);
	
	% kmers
	\node[anchor = west, node font = \figsmall] (kmers) at (ohe) {
		\setlength{\tabcolsep}{.5em}%
		\begin{tabular}{ccccc}%
			\toprule%
			\textbf{kmer-1} & \textbf{kmer-2} & \raisebox{.5ex}{\textbf{\dots}} & \textbf{kmer-4095} & \textbf{kmer-4096} \\
			\figtiny AAAAAA & \figtiny AAAAAC & \raisebox{.5ex}{\dots} & \figtiny TTTTTG & \figtiny TTTTTG \\%
			\midrule%
			\pgfmathparse{int(random(4) - 1)}\pgfmathresult &
			\pgfmathparse{int(random(4) - 1)}\pgfmathresult &
			\raisebox{.5ex}{\dots} &
			\pgfmathparse{int(random(4) - 1)}\pgfmathresult &
			\pgfmathparse{int(random(4) - 1)}\pgfmathresult \\%
			\bottomrule%
		\end{tabular}%
	};
	
	\draw[very thick, -Latex] (kmers.east) -- ++(.5, 0) coordinate[xshift = .3333em] (model);
	
	% model
	\node[right = .05cm of model, align = flush right, font = \figsmall, yshift = -.5\baselineskip] (lm) {%
		$\log_2(\text{enhancer strength}) =$\hspace*{7em}\\
		$\beta_0 + \beta_1\text{kmer-1} + \beta_2\text{kmer-2} + \dots$\\
		$+ \beta_{4095}\text{kmer-4095} + \beta_{4096}\text{kmer-4096}$%
	};
	
	\begin{pgfonlayer}{background}
		\draw[fill = OkabeIto5!25, rounded corners] (model) |- ($(lm.north east) + (.05, .05cm + \baselineskip)$) node[pos = .75, anchor = north, fill = none, draw = none, text = black, font = \bfseries] (model top) {linear regression} |- ($(lm.south west) + (-.05, -.05)$) coordinate[pos = .25] (model out) coordinate[pos = .75] (model bottom) -- cycle;
	\end{pgfonlayer}
	
	\draw[very thick, -Latex] (model out) ++(.3333em, 0) -- ++(.5, 0) coordinate (predictions);
	
	% predictions
	\pgfmathsetmacro{\predLight}{2 + rand * 4}
	\pgfmathsetmacro{\predDark}{\predLight + rand * 0.5 * \predLight}
	\pgfmathsetmacro{\predWarm}{\predLight + rand * 0.1 * \predLight}
	\pgfmathsetmacro{\predCold}{\predLight + rand * 0.1 * \predLight}
	\pgfmathsetmacro{\predMaize}{1.5 + rand * 3}
	
	\node[anchor = west, node font = \figsmall, inner xsep = 0pt, xshift = .3333em] (predictions) at (predictions) {%
		\pgfmathsetseed{928}%
		\setlength{\tabcolsep}{.5em}%
		\begin{tabular}{rrrrr}%
			\toprule%
			\textbf{\usename{light}} & \textbf{\usename{dark}} & \textbf{\usename{warm}} & \textbf{\usename{cold}} & \textbf{\usename{maize}} \\%
			\midrule%
			\pgfmathprintnumber[fixed, zerofill, precision = 2]{\predLight} &
			\pgfmathprintnumber[fixed, zerofill, precision = 2]{\predDark} &
			\pgfmathprintnumber[fixed, zerofill, precision = 2]{\predWarm} &
			\pgfmathprintnumber[fixed, zerofill, precision = 2]{\predCold} &
			\pgfmathprintnumber[fixed, zerofill, precision = 2]{\predMaize} \\%
			\bottomrule%
		\end{tabular}%
	};
	
	% text
	\node[anchor = south, font = \strut, yshift = .3em] at (seq |- model top.north) {170\,bp sequence};
	\node[anchor = south, font = \strut, yshift = .3em] at (kmers |- model top.north) {count occurence of all possible 6-mers};
	\node[anchor = south, font = \strut, yshift = .3em] at (model top.north) {linear model};
	\node[anchor = south, font = \strut, yshift = .3em] (blub) at (predictions |- model top.north) {predict $\log_2$(enhancer strength)};
	
	
	%%% enhancer strength predictions from linear model based on kmer counts %%%
	\coordinate[yshift = -\columnsep] (cor kmers) at (lm kmers |- model bottom);
	
	\begin{pggfplot}[%
		at = {(cor kmers)},
		total width = \textwidth,
		xlabel = {\textbf{measurement}: \enrichment},
		ylabel = {\textbf{prediction}:\\\enrichment},
		facet 1/.append style = {
			pggf colorbar = {%
				title = count,
				log2,
				horizontal,
				ytick = {0, 3, ..., 15}
			}
		},
		title is csname,
		title color from name,
		xytick = {-12, -10, ..., 12},
		ylabel add to width = .6\baselineskip,
		tight y
	]{lm_kmers_predictions}
	
		\pggf_annotate(diagonal)
	
		\pggf_hexbin(bins = 50)
	
		\pggf_stats(stats = {n,spearman,rsquare})
	
	\end{pggfplot}	
	
	
	%%% subfigure labels %%%
	\subfiglabel{lm kmers}
	\subfiglabel{cor kmers}
	
\end{tikzpicture}%
			\caption{%
				\captiontitle{A linear model based on 6-mer counts predicts enhancer strength with limited accuracy}%
				\subfigtwo{A}{B} A linear model based on the counts of all possible 6-mers in a given sequence was trained to predict enhancer strength \parensubfig{A}. The model predictions were compared with the experimentally determined enhancer strengths in a held-out test set comprising 10\% of the test sequences \parensubfig{B}. Hexbin plots are as defined in \cref{fig:enhScreen}.
			}%
			\label{efig:lm_kmers}%
		\end{efig}
		
	\fi
	%%% Extended data figures end
	
	%%% Supplementary figures start
	\ifsupp
	
		\begin{sfig}
			\input{SuppFig1}%
			\caption{%
				\captiontitle{Plant STARR-seq is reproducible}%
				\subfigrange{A}{E} Hexbin plots (color represents the count of points in each hexagon) of the correlation between enhancer strength in independent Plant STARR-seq replicate experiments in the indicated conditions. The dashed lines represents a $y = x$ line. Pearson's $R^2$, Spearman's $\rho$, and the number ($n$) of samples are indicated.
			}%
			\label{sfig:tamsACR_repCor}%
		\end{sfig}
		
		\begin{sfig}
			\begin{tikzpicture}

	%%% correlation of enhancer strength across conditions %%%
	\coordinate (cor) at (0, 0);
	
	\tikzset{
		condX/.store in = \condX,
		condY/.store in = \condY
	}
	
	\begin{pggfplot}[%
		at = {(cor)},
		total width = \textwidth,
		total height = \textwidth,
		xlabel = {\textbf{\expandafter\csname\condX\endcsname}: \enrichment},
		ylabel = {\textbf{\expandafter\csname\condY\endcsname}: \enrichment},
		facet 1/.append style = {pggf colorbar = {title = count, log2, horizontal}},
		title is csname,
		title color from name,
		xytick = {-12, -10, ..., 12},
		empty facet warning = false,
		label each axis,
	]{tamsACR_cor_conds}
	
		\pggf_annotate(
			diagonal,
			skip facets = {2, 3, 4, 7, 8, 12}
		)
	
		\pggf_hexbin(bins = 50)
	
		\pggf_stats(stats = {n,spearman,rsquare})
	
	\end{pggfplot}

\end{tikzpicture}%
			\caption{%
				\captiontitle{Plant STARR-seq detects species- and condition-specific differences in enhancer strength}%
				Hexbin plots (as defined in \cref{sfig:tamsACR_repCor}) of the correlation between enhancer strength in different conditions.
			}%
			\label{sfig:tamsACR_condCor}%
		\end{sfig}
		
		\begin{sfig}
			\begin{tikzpicture}

	%%% dual-luciferase construct
	\coordinate[yshift = -\columnsep] (DL construct) at (0, 0);
	
	\coordinate[shift = {(.1, -1)}] (construct) at (DL construct);
	
	\draw[dashed, thick] (construct) -- ++(.5, 0) coordinate (c1);
	
	\node[draw = black, thin, anchor = west, text depth = 0pt] (Luc) at ($(c1) + (.1, 0)$) {Luc};
	
	\draw[line width = .2cm, DarkOliveGreen3] (Luc.east) ++(.4, 0) coordinate (c2) -- ++(.5, 0) coordinate[pos = .5] (pro1) coordinate (c3);
	\draw[-{Stealth[round]}, thick] (c2) ++(-.5\pgflinewidth, 0) |- ++(-.4, .3);
	
	\node[draw = black, thin, anchor = west, text depth = 0pt] (BlpR) at ($(c3) + (.3, 0)$) {BlpR};
	
	\draw[line width = .2cm, PaleGreen3] (BlpR.east) ++(.4, 0) coordinate (c4) -- ++(.5, 0) coordinate (c5);
	\draw[-{Stealth[round]}, thick] (c4) ++(-.5\pgflinewidth, 0) |- ++(-.4, .3);
	
	\draw[line width = .25cm, -Triangle Cap, 35Senhancer]  (c5) ++(.4, 0) -- ++(.6, 0) coordinate[pos = .5] (enh) coordinate (c6);
	
	\draw[line width = .2cm, 35Spromoter] (c6) ++(.1, 0) -- ++(.3, 0) coordinate[pos = .5] (pro2) coordinate (c7);
	\draw[-{Stealth[round]}, thick] (c7) ++(.5\pgflinewidth, 0) |- ++(.4, .3);
	
	\node[draw = black, thin, anchor = west, text depth = 0pt] (NLuc) at ($(c7) + (.4, 0)$) {NanoLuc};
	
	\draw[dashed, thick] (NLuc.east) ++(.1, 0) coordinate (c8) -- ++(.5, 0) coordinate (construct end);
	
	\begin{pgfonlayer} {background}
		\draw[thick] (c1) -- (Luc.west) (Luc.east) -- (BlpR.west) (BlpR.east) -- (NLuc.west) (NLuc.east) -- (c8);
	\end{pgfonlayer}
	
	\draw[Latex-] (enh) ++(0, .175) -- ++(0, .4) node[anchor = south, text depth = 0pt] {enhancer};
	\draw[Latex-] (pro1) ++(0, -.15) -- ++(0, -.4) node[anchor = north, text depth = 0pt] {\textit{AtUBQ10} promoter};
	\draw[Latex-] (pro2) ++(0, -.15) -- ++(0, -.4) node[anchor = north, text depth = 0pt] (l35Spr) {35S minimal promoter};

	\draw[decorate, decoration = {brace, amplitude = .4cm, aspect = .825}] (l35Spr.south -| construct end) -- (l35Spr.south -| construct) coordinate[pos = .825, yshift = -.4cm] (curly);
	
	\draw[ultra thick, -Latex] (curly) ++(0, .75\pgflinewidth) -- ++(0, -.3) arc[start angle = 180, end angle = 270, radius = .8cm] -- ++(.3, 0) coordinate[shift = {(1, .1)}] (Arabidopsis);
	\planticon[scale = .9]{arabidopsis} at (Arabidopsis);
	\node [anchor = north east, align = center, inner xsep = 0pt, shift = {(.1, -.1)}] at (curly) {integrate\\into\\\Arabidopsis\\genome};
	
	\draw[ultra thick, -Latex] (Arabidopsis) ++(1, -.1) -- ++(.8, 0) node[anchor = west, align = center] {measure\\luciferase (Luc)\\and nanoluciferase\\(NanoLuc) activity};

\end{tikzpicture}%
			\caption{%
				\captiontitle{A dual-luciferase assay to measure enhancer strength in stable \usename{Arabidopsis} lines}%
				Transgenic Arabidopsis lines were generated with T-DNAs harboring a constitutively expressed luciferase (Luc) gene and a nanoluciferase (NanoLuc) gene under control of a 35S minimal promoter coupled to a candidate enhancer. Luciferase and nanoluciferase activity was measured in the T2 generation.
			}%
			\label{sfig:DL_assay}%
		\end{sfig}
		
		\begin{sfig}
			\begin{tikzpicture}

	%%% enhancer strength of non-ACR controls %%%
	\coordinate (nonACR controls) at (0, 0);
	
	\pgfplotsset{
		non-ACR/.style = {fill = pgffillcolor!50!black},
		ACR/.style = {}
	}
	
	\begin{pggfplot}[% 
		at = (nonACR controls),
		yshift = -1.5\baselineskip,
		total width = \textwidth,
		xlabel = species,
		ylabel = \enrichment,
		xticklabels are csnames,
		title is csname,
		title color from name,
		enlarge x limits = 0,
		enlarge y limits = {lower = 0.1, upper = 0.15},
		ytick = {-12, -10, ..., 12},
		xticklabel style = {font = \specieslongfalse},
		legend style = {anchor = south east, at = {(1, 1)}, yshift = 2\baselineskip},
		legend columns = 3,
		legend image width = .75\baselineskip,
		every legend image post/.append style = {fill opacity = 0.5},
	]{tamsACR_non-ACR_ctrls}
	
		\pggf_annotate(zeroline)
		
		\pggf_annotate(
			manual legend = {
				{\textbf{sequence origin:}\hspace*{-\groupplotsep}} = {empty legend},
				ACR = {half violin left legend = y, fill = gray!25!white},
				non-ACR = {half violin right legend = y, fill = gray!75!black}
			},
			only facet = 5
		)
	
		\pggf_violin(
			half violin,
			style from column = {fill = xticklabels, ACR},
			sample size = {rotate = 90},
			sample size left = {anchor = south west},
			sample size right = {anchor = north west},
			forget plot
		)
		
		\pggf_signif()
	
	\end{pggfplot}

\end{tikzpicture}%
			\caption{%
				\captiontitle{Accessible chromatin regions are enriched for active enhancers}%
				Enhancer strength of sequences derived from accessible chromatin regions (ACR, left) and non-ACR sequences (right) randomly sampled from inaccessible regions in the genomes of the indicated species. At, \usename{Arabidopsis}; Sl, \usename{tomato}; Zm, \usename{maize}; Sb, \usename{sorghum}. Violin plots represent the kernel density distribution and the box plots inside represent the median (center line), upper and lower quartiles (box limits), and 1.5$\times$ interquartile range (whiskers) for all corresponding sequences. Significant differences between ACR and non-ACR sequences were determined by two-sided Wilcoxon rank-sum tests. The obtained $p$ values were Bonferroni-adjusted separately for each condition and species and are indicated: \printsigniflevels. Exact $p$ values are listed in Supplementary Data 4.% TODO: make sure reference is correct
			}%
			\label{sfig:accessibility}%
		\end{sfig}
		
		\begin{sfig}
			\begin{tikzpicture}

	%%% examples of TF effects %%%
	\coordinate (TF effects) at (0, 0);
	
	% set styles for half violin plots
	\pgfplotsset{
		absent/.code = {
			\def\tmpA{dark}
			\ifx\tmpA\pggfrowname\relax
				\pgfkeysalso{fill = \pggfrowname}
			\else
				\pgfkeysalso{fill = \pggfrowname!50!black}
			\fi
		},
		present/.code = {
			\def\tmpA{dark}
			\ifx\tmpA\pggfrowname\relax
				\pgfkeysalso{fill = \pggfrowname!50!white}
			\else
				\pgfkeysalso{fill = \pggfrowname}
			\fi
		}
	}
	
	\pgfplotscreateplotcyclelist{none}{black}
	
	\begin{pggfplot}[% 
		at = {(TF effects)},
		total width = \textwidth,
		xlabel = transcription factor family (motif ID),
		ylabel = {\enrichment[GC-normalized]},
		xticklabel style = {rotate = 45, anchor = north east, inner sep = 0.236em},
		xticklabels are csnames,
		title is csname,
		title color from name,
		enlarge x limits = 0,
		enlarge y limits = {lower = 0.1, upper = 0.15},
		ytick = {-12, -10, ..., 12},
		cycle list name = none,
		legend style = {anchor = south east, at = {(1, 1)}},
		legend columns = 3,
		legend image width = .75\baselineskip,
		every legend image post/.append style = {fill opacity = 0.5},
		tight y
	]{TF_effect_examples}
	
		\pggf_annotate(zeroline)
		
		\pggf_annotate(
			manual legend = {
				{\textbf{transcription factor binding site:}\hspace*{-\groupplotsep}} = {empty legend},
				absent = {half violin left legend = y, fill = gray!75!black},
				present = {half violin right legend = y, fill = gray!25!white}
			},
			only facet = 1
		)
	
		\pggf_violin(
			half violin,
			style from column = {hits},
			draw = black,
			sample size = {rotate = 90},
			sample size left = {anchor = south west},
			sample size right = {anchor = north west},
			forget plot
		)
		
		\pggf_signif()
	
	\end{pggfplot}

\end{tikzpicture}%
			\caption{%
				\captiontitle{The presence of transcription factor binding sites affects enhancer strength}%
				Putative transcription factor binding sites were identified by scanning the test sequences with known binding motifs for transcription factors of the indicated families. The enhancer strength of test sequences without or with one putative binding site is shown in violin plots as defined in \cref{sfig:accessibility}. Significant differences between sequences with and without a binding site were determined by two-sided Wilcoxon rank-sum tests. The obtained $p$ values were Bonferroni-adjusted separately for each condition and species and are indicated: \printsigniflevels. Exact $p$ values are listed in Supplementary Data 4.% TODO: make sure reference is correct
			}%
			\label{sfig:TF_effects}%
		\end{sfig}
		
		\begin{sfig}
			\begin{tikzpicture}

	%%% correlation of enhancer strength across conditions %%%
	\coordinate (cor) at (0, 0);
	
	\tikzset{
		condX/.store in = \condX,
		condY/.store in = \condY
	}
	
	\begin{pggfplot}[%
		at = {(cor)},
		total width = \textwidth,
		total height = \textwidth,
		xlabel = {\textbf{\expandafter\csname\condX\endcsname}: $\log_2$(effect size)},
		ylabel = {\textbf{\expandafter\csname\condY\endcsname}: $\log_2$(effect size)},
		facet 1/.append style = {pggf colorbar = {title = count, log2, horizontal}},
		title is csname,
		title color from name,
		empty facet warning = false,
		label each axis,
	]{TF_effect_cor_conds}
	
		\pggf_annotate(
			diagonal,
			skip facets = {2, 3, 4, 7, 8, 12}
		)
	
		\pggf_scatter()
	
		\pggf_stats(stats = {n,spearman,rsquare})
	
	\end{pggfplot}

\end{tikzpicture}%
			\caption{%
				\captiontitle{Transcription factor binding sites show species- and condition-specific effects on enhancer strength}%
				Correlation between the effects of transcription factor binding sites (as determined in \cref{fig:TFBS}) in different conditions. Each dot represents a different transcription factor binding motif. Pearson's $R^2$, Spearman's $\rho$, and the number ($n$) of samples are indicated.
			}%
			\label{sfig:cor_TF_effects}%
		\end{sfig}
		
		\begin{sfig}
			\begin{tikzpicture}

	\pggfset{
		ylabel style = {xshift = -1em},
	}

	%%% Calculate plot width %%%
	\makeatletter
		\figlarge
		\pgfmathsetlength{\templength}{(\textwidth - \columnsep - 2\pggf@ylabelwidth - 2\groupplotsep - 2 * (.85 * \baselineskip + 2 * \pgfkeysvalueof{/pgf/inner ysep})) / 4}
		\fignormal
	\makeatother

	%%% Replicate correlation light %%%
	\coordinate (light) at (0, 0);
	
	\planticon[height = .5cm]{light} at ($(light) + (.25, -.75)$);
	
	\begin{pggfplot}[%
		at = {(light)},
		width = \templength,
		xlabel = \enrichment,
		ylabel = \enrichment,
		row title = {light},
		row title is csname,
		row title color from name,
		xytick = {-12, -10, ..., 12},
		facet 1/.append style = {pggf colorbar = {title = count, log2, horizontal}}
	]{eVal_cor_reps_light}
	
		\pggf_annotate(diagonal)
	
		\pggf_hexbin(bins = 50)
	
		\pggf_stats(stats = {n,spearman,rsquare})
	
	\end{pggfplot}
	
	
	%%% Replicate correlation dark %%%
	\coordinate[xshift = \columnsep] (dark) at (pggfplot c3r1.outer north east);
	
	\planticon[height = .5cm]{dark} at ($(dark) + (.25, -.75)$);
	
	\begin{pggfplot}[%
		at = {(dark)},
		width = \templength,
		xlabel = \enrichment,
		ylabel = \enrichment,
		col title = 1 vs. 2,
		row title = {dark},
		row title is csname,
		row title color from name,
		xytick = {-12, -10, ..., 12},
		facet 1/.append style = {pggf colorbar = {title = count, log2, horizontal}}
	]{eVal_cor_reps_dark}
	
		\pggf_annotate(diagonal)
	
		\pggf_hexbin(bins = 50)
	
		\pggf_stats(stats = {n,spearman,rsquare})
	
	\end{pggfplot}
	
	
	%%% Replicate correlation warm %%%
	\coordinate[yshift = -\columnsep] (warm) at (light |- xlabel.south);
	
	\planticon[height = .5cm]{warm} at ($(warm) + (.25, -.75)$);
	
	\begin{pggfplot}[%
		at = {(warm)},
		width = \templength,
		xlabel = \enrichment,
		ylabel = \enrichment,
		row title = {warm},
		row title is csname,
		row title color from name,
		xytick = {-12, -10, ..., 12},
		facet 1/.append style = {pggf colorbar = {title = count, log2, horizontal}}
	]{eVal_cor_reps_warm}
	
		\pggf_annotate(diagonal)
	
		\pggf_hexbin(bins = 50)
	
		\pggf_stats(stats = {n,spearman,rsquare})
	
	\end{pggfplot}
	
	
	%%% Replicate correlation cold %%%
	\coordinate[xshift = \columnsep] (cold) at (pggfplot c3r1.outer north east);
	
	\planticon[height = .5cm]{cold} at ($(cold) + (.25, -.75)$);
	
	\begin{pggfplot}[%
		at = {(cold)},
		width = \templength,
		xlabel = \enrichment,
		ylabel = \enrichment,
		col title = 1 vs. 2,
		row title = {cold},
		row title is csname,
		row title color from name,
		xytick = {-12, -10, ..., 12},
		facet 1/.append style = {pggf colorbar = {title = count, log2, horizontal}}
	]{eVal_cor_reps_cold}
	
		\pggf_annotate(diagonal)
	
		\pggf_hexbin(bins = 50)
	
		\pggf_stats(stats = {n,spearman,rsquare})
	
	\end{pggfplot}
	
	
	%%% Replicate correlation maize %%%
	\coordinate[yshift = -\columnsep] (maize) at (light |- xlabel.south);
	
	\planticon[height = .5cm]{maize} at ($(maize) + (.25, -.75)$);
	
	\begin{pggfplot}[%
		at = {(maize)},
		width = \templength,
		xlabel = \enrichment,
		ylabel = \enrichment,
		col title = 1 vs. 2,
		row title = {maize},
		row title is csname,
		row title color from name,
		xytick = {-12, -10, ..., 12},
		facet 1/.append style = {
			pggf colorbar = {%
				title = count,
				log2,
				horizontal,
				ytick = {0, 3, ..., 15}
			}
		},
		tight y
	]{eVal_cor_reps_maize}
	
		\pggf_annotate(diagonal)
	
		\pggf_hexbin(bins = 50)
	
		\pggf_stats(stats = {n,spearman,rsquare})
	
	\end{pggfplot}
	
	
	%%% PCA %%%
	\coordinate[xshift = \columnsep] (PCA) at (pggfplot c1r1.outer north east);
	
	\tikzset{
		PC1/.store in = \PCone,
		PC2/.store in = \PCtwo
	}
	
	\begin{pggfplot}[%
		at = {(PCA)},
		width = \templength,
		tight,
		xlabel = {PC2 (\pgfmathprintnumber[precision = 0]{\PCtwo}\%)},
		ylabel = {PC1 (\pgfmathprintnumber[precision = 0]{\PCone}\%)},
		scatter styles = {light,dark,warm,cold,maize},
		ticks = none,
		title = PCA,
	]{eVal_PCA}
	
		\pggf_scatter(
			mark size = 2.5,
			opacity = 0.5,
			style from column = condition
		)
	
		\pggf_annotate(scatter legend*)
	
	\end{pggfplot}
	
	
	%%% subfigure labels %%%
	\subfiglabel{light}
	\subfiglabel{dark}
	\subfiglabel{warm}
	\subfiglabel{cold}
	\subfiglabel{maize}
	\subfiglabel{PCA}

\end{tikzpicture}%
			\caption{%
				\captiontitle{The validation library yielded reproducible results}%
				\subfigrange{A}{E} Hexbin plots (as defined in \cref{sfig:tamsACR_repCor}) of the correlation between enhancer strength in independent Plant STARR-seq replicate experiments with the validation library in the indicated conditions.\nextentry
				\subfig{F} Principal component analysis of enhancer strength determined in Plant STARR-seq replicate experiments with the validation library.
			}%
			\label{sfig:eVal_repCor}%
		\end{sfig}

		\begin{sfig}
			\begin{tikzpicture}

	%%% correlation of enhancer strength across tamsACR and eVal libraries %%%
	\coordinate (cor) at (0, 0);
	
	\begin{pggfplot}[%
		at = {(cor)},
		total width = \textwidth,
		xlabel = {\textbf{ACR library}: \enrichment},
		ylabel = {\textbf{validation library}:\\\enrichment},
		facet 1/.append style = {pggf colorbar = {title = count, log2, horizontal}},
		title is csname,
		title color from name,
		xytick = {-12, -10, ..., 12},
		ylabel add to width = .6\baselineskip,
		tight y
	]{tamsACR_eVal_cor}
	
		\pggf_annotate(diagonal)
	
		\pggf_hexbin(bins = 50)
	
		\pggf_stats(stats = {n,spearman,rsquare})
	
	\end{pggfplot}

\end{tikzpicture}%
			\caption{%
				\captiontitle{The validation library results reproduce the results obtained with the ACR library}%
				Hexbin plots (as defined in \cref{sfig:tamsACR_repCor}) of the correlation between the strength of enhancers included in the ACR library and the validation library.
			}%
			\label{sfig:cor_tamsACR_eVal}%
		\end{sfig}
		
		\begin{sfig}
			\begin{tikzpicture}

	%%% effect of TFBS shuffling %%%
	\coordinate (shuffle effect) at (0, 0);

	\tikzset{
		issignif/.code = {
			\pgfmathparse{#1 <= 0.05}
			\ifnum1=\pgfmathresult\relax
				\pgfkeysalso{fill = OkabeIto4}
			\else
				\pgfkeysalso{fill = gray}
			\fi
		}
	}

	\begin{pggfplot}[%
		at = {(shuffle effect)},
		total width = \textwidth,
		xlabel = transcription factor family (motif ID),
		ylabel = {\enrichment[rel. to WT]},
		xticklabel style = {anchor = east, rotate = 90, inner ysep = 0pt, inner xsep = .3333em},
		xticklabels are csnames,
		title is csname,
		title color from name,
		enlarge x limits = 0,
		enlarge y limits = {lower = 0.1},
		ytick = {-10, -9, ..., 10},
		legend style = {anchor = south east, at = {(1, 0)}, yshift = \baselineskip},
		legend image width = .75\baselineskip,
		every legend image post/.append style = {fill opacity = 0.5},
		tight y
	]{TFBS_shuffle_effect}
	
		\pggf_annotate(zeroline)
		
		\pggf_annotate(
			manual legend = {
				significant = {boxplot legend = y, fill = OkabeIto4},
				not significant = {boxplot legend = y, fill = gray}
			},
			only facet = 1
		)
	
		\pggf_boxplot(
			style from column = {issignif = p_value},
			sample size = {font = \fontsize{4}{5}\selectfont},
			forget plot
		)
	
	\end{pggfplot}

\end{tikzpicture}%
			\caption{%
				\captiontitle{Removing transcription factor binding sites affects enhancer strength}%
				Relative enhancer strength of sequences with mutations destroying a single strong ($p$ value \textle{} 0.0001) binding site for transcription factors of the indicated family. For each sequence, the enhancer strength was normalized to the strength of the corresponding wild-type sequence.\nextentry
				Box plots represent the median (center line), upper and lower quartiles (box limits), 1.5$\times$ interquartile range (whiskers), and outliers (points) for all corresponding sequences. Numbers at the bottom of the plots indicate the number of samples in each group. Significant differences from a null distribution were determined using one-sample Wilcoxon rank-sum tests. Box plots for significant groups (Bonferroni-adjusted $p$ value \textle{} 0.05) are shown in green. Non-significant groups are represented by gray box plots. Exact $p$ values are listed in Supplementary Data 4.% TODO: make sure reference is correct
			}%
			\label{sfig:TFBS_shuffle_strength}%
		\end{sfig}
		
		\begin{sfig}
			\begin{tikzpicture}

	%%% Variability of TF effect sizes in synthetic enhancers %%%
	%% scheme %%
	\coordinate (scheme) at (0, 0);
	
	
	
	\foreach \x in {1, 2, 3}{
		\node[anchor = west, node font = \figsmall, yshift = -.75cm - \x * .75cm] (seq\x) at (scheme) {seq\x};
	}
	
	\distance{seq3.east}{\threecolumnwidth, 0}
	\pgfmathsetlength{\templength}{(\xdistance - 2\columnsep) / 3}
	
	\pgfmathsetseed{928}
	
	\colorlet{seq1}{gray!75!white}
	\colorlet{seq2}{gray}
	\colorlet{seq3}{gray!75!black}
	
	% TFBS insertion effect
	\node[anchor = north west, node font = \figsmall\bfseries, text width = \templength, align = center, inner xsep = 0pt, shift = {(.5\columnsep, -.5\columnsep)}, align = center] (seq+pos) at (scheme -| seq3.east) {~\\TFBS\\insertion\\effect size};
	
	\foreach \x in {1, 2, 3}{
		\draw[ultra thick, seq\x] (seq\x -| seq+pos.west) -- ++(\templength, 0) coordinate[pos = .165] (p1) coordinate[pos = .5] (p2) coordinate[pos = .835] (p3);
		\foreach \y in {1, 2, 3}{
			\draw[->] (p\y) +(0, .05) -- ++(0, -.35) node[anchor = north, node font = \figtiny] {%
				\pgfmathparse{rnd}%
				\expandafter\xdef\csname tmp\x\y\endcsname{\pgfmathresult}%
				\pgfmathprintnumber[fixed, zerofill, precision = 1]{\csname tmp\x\y\endcsname}%
			};
		}
	}
	
	% standard deviation across positions
	\node[anchor = base west, node font = \figsmall\bfseries, text width = \templength, align = center, inner xsep = 0pt, xshift = .5\columnsep, align = center] (seq only) at (seq+pos.base east) {standard\\deviation\\across\\positions};
	
	\foreach \x in {1, 2, 3}{
		\draw[ultra thick, seq\x] (seq\x -| seq only.west) -- ++(\templength, 0) coordinate[pos = .165] (p1) coordinate[pos = .5] (p2) coordinate[pos = .835] (p3);
		\coordinate[yshift = -.35cm] (e\x) at (p2);
		\foreach \y in {1, 2, 3}{
			\draw[->] (p\y) +(0, .05) -- ++(0, -.05) to[out = -90, in = 90] ($(e\x) + (0, .05)$) -- (e\x);
		}
		\node[anchor = north, node font = \figtiny] at (e\x) {%
			\pgfmathsetmacro{\tmpmean}{(\csname tmp\x1\endcsname + \csname tmp\x2\endcsname + \csname tmp\x3\endcsname) / 3}%
			\pgfmathparse{sqrt(((\csname tmp\x1\endcsname - \tmpmean)^2 + (\csname tmp\x2\endcsname - \tmpmean)^2 + (\csname tmp\x3\endcsname - \tmpmean)^2) / 2)}
			\pgfmathprintnumber[fixed, zerofill, precision = 1]{\pgfmathresult}%
		};
	}
	
	% standard deviation across sequences
	\node[anchor = base west, node font = \figsmall\bfseries, text width = \templength, align = center, inner xsep = 0pt, xshift = .5\columnsep, align = center] (pos only) at (seq only.base east) {standard\\deviation\\across\\sequences};
	
	\foreach \x in {1, 2, 3}{
		\draw[ultra thick, seq\x] (seq\x -| pos only.west) -- ++(\templength, 0) coordinate[pos = .165] (p1) coordinate[pos = .5] (p2) coordinate[pos = .835] (p3);
	}
	\foreach \y in {1, 2, 3}{
		\draw[->] (seq1 -| p\y) +(0, .05) -- (e3 -| p\y) node[anchor = north, node font = \figtiny] {%
			\pgfmathsetmacro{\tmpmean}{(\csname tmp1\y\endcsname + \csname tmp2\y\endcsname + \csname tmp3\y\endcsname) / 3}%
			\pgfmathparse{sqrt(((\csname tmp1\y\endcsname - \tmpmean)^2 + (\csname tmp2\y\endcsname - \tmpmean)^2 + (\csname tmp3\y\endcsname - \tmpmean)^2) / 2)}
			\pgfmathprintnumber[fixed, zerofill, precision = 1]{\pgfmathresult}%
		};
	}
	
	
	%% plot %%
	\coordinate (TF sd) at (scheme -| \textwidth - \twothirdcolumnwidth, 0);
	
	\expandafter\def\csname acrossPos\endcsname{across\\positions}
	\expandafter\def\csname acrossSeq\endcsname{across\\sequences}
	
	\begin{pggfplot}[%
		at = {(TF sd)},
		total width = \twothirdcolumnwidth,
		xlabel = {},
		ylabel = {standard deviation of effect sizes},
		title is csname,
		title color from name,
		enlarge x limits = 0,
		enlarge y limits = {lower = 0.1, upper = 0.2},
		y decimals = 1,
		xticklabels are csnames,
		xticklabel style = {align = center},
		tight y
	]{synEnh_TF_sd}
	
		\pggf_annotate(zeroline)
	
		\pggf_violin()
		
		\pggf_signif()
	
	\end{pggfplot}
	
	
	%%% subfigure labels %%%
	\subfiglabel{scheme}
	\subfiglabel{TF sd}

\end{tikzpicture}%
			\caption{%
				\captiontitle{Sequence context has a stronger impact on enhancer strength than position of a transcription factor binding site}%
				A transcription factor binding site (TFBS) was inserted into 50 random sequences in one of three different positions, and the effect size of each binding site was determined as the fold-change in enhancer strength between the sequence with and without this binding site. For each binding site, the standard deviation of its effects sizes was calculated across the three positions within each sequence (across positions) or across all random sequences for each position (across sequences). Violin plots are as defined in \cref{sfig:accessibility}. Significant differences were determined by two-sided Wilcoxon rank-sum tests: \printsigniflevels. Exact $p$ values are listed in Supplementary Data 4.% TODO: make sure reference is correct
			}%
			\label{sfig:synEnh_TF_sd}%
		\end{sfig}
		
		\begin{sfig}
			\begin{tikzpicture}

	%%% Correlation between TF effects and TF model coefficients %%%
	\coordinate (correlation) at (0, 0);
	
	\def\hits{\textbf{model input:} transcription\\factor binding site counts}
	\def\scores{\textbf{model input:} transcription\\factor motif scores}
	
	\begin{pggfplot}[%
		at = {(correlation)},
		total width = \textwidth,
		xlabel = {$\log_2$(effect size)},
		ylabel = linear model coefficient,
		title is csname,
		col title color from name,
		row title lines = 2,
		xytick = {-1.2, -.9, ..., 1.2},
		ylabel add to width = .6\baselineskip,
	]{lm_TF_cor_effect}
	
		\pggf_annotate(diagonal)
		
		\pggf_trendline()
	
		\pggf_scatter(
			color from column name
		)
	
		\pggf_stats(stats = {n,spearman,rsquare})
	
	\end{pggfplot}

\end{tikzpicture}%
			\caption{%
				\captiontitle{Linear model coefficients are correlated with the effects of transcription factor binding sites in the ACR library}%
				Correlation between the coefficients of the linear models from \cref{efig:lm_TFs} and the transcription factor binding site effects in the ACR library (see \cref{fig:TFBS}\subfigunformatted{A}). The solid and dashed lines represents a $y = x$ line and a linear regression line, respectively. Pearson's $R^2$, Spearman's $\rho$, and the number ($n$) of samples are indicated.
			}%
			\label{sfig:cor_TF_effects_lm}%
		\end{sfig}
		
		\begin{sfig}
			\input{SuppFig12}%
			\caption{%
				\captiontitle{TF-MoDISco patterns match known transcription factor binding motifs}%
				\subfig{A} DNA Patterns identified by TF-MoDISco that positively contribute to plantGREP predictions of enhancer strength in tobacco leaves in the light.\nextentry
				\subfig{B} Transcription factor binding motifs that match the patterns in \parensubfig{A}.
			}%
			\label{sfig:TF-MoDISco_pos_light}%
		\end{sfig}
		
		\begin{sfig}
			\begin{tikzpicture}
	\def\NA{\relax}
	\tikzset{
		TF name/.style = {
			/pgfplots/extra description/.append code = {
				\node[anchor = west] at (1, .5) {\usename{#1}};
			}
		}
	}

	%%% TF-MoDISco patterns %%%
	\coordinate (modisco) at (0, 0);
	
	\planticon[height = .5cm]{light} at ($(modisco) + (.6, -.25)$);
	
	% get number of bases in motif
	\pgfplotstableread[col sep = tab]{rawData/modisco_motifs_neg_light_axes.tsv}\axistable
	
	\pgfplotstablegetelem{0}{xmax}\of\axistable
	\pgfmathsetmacro{\nbases}{int(\pgfplotsretval)}

	\begin{pggfplot}[%
		at = {(modisco)},
		width = \nbases * 2.6mm,
		height = 6.6mm,
		xlabel = {},
		ylabel = {average DeepLIFT contribution score},
		y decimals = 2,
		enlarge x limits = 0,
		enlarge y limits = .1,
		xmajorticks = false,
		row titles from table = false,
		col title = {\raisebox{.5\baselineskip}{TF-MoDISco pattern}},
		col title lines = 2,
		xlabel height = 5pt,
		tight,
	]{modisco_motifs_neg_light}
	
		\pggf_annotate(
			zeroline,
			annotation = {
				\node[anchor = south west, node font = \figsmall, text depth = 0pt] at (current axis.south west) {pattern \pggfrowname};
			},
			skip facet = 13
		)
	
		\pggf_annotate(
			zeroline,
			annotation = {
				\node[anchor = north west, node font = \figsmall, text depth = 0pt] at (current axis.north west) {pattern \pggfrowname};
			},
			only facet = 13
		)
	
		\pggf_logo(letter colors = DNA)
	
	\end{pggfplot}
	
	
	%%% matching TF motifs %%%
	\coordinate[xshift = \columnsep] (TFs) at (modisco -| pggfplot c1r1.east);
	
	% get number of bases in motif
	\pgfplotstableread[col sep = tab]{rawData/modisco_TFs_neg_light_axes.tsv}\axistable
	
	\pgfplotstablegetelem{0}{xmax}\of\axistable
	\pgfmathsetmacro{\nbases}{int(\pgfplotsretval)}
	
	\begin{pggfplot}[%
		at = {(TFs)},
		width = \nbases * 2.6mm,
		height = 6.6mm,
		xlabel = {},
		ylabel = {information content},
		ymin = 0,
		ymax = 2,
		ytick = {0, 1, 2},
		enlarge x limits = 0,
		enlarge y limits = {lower = 0, upper = .15},
		xmajorticks = false,
		row titles from table = false,
		col title = {matched transcription\\factor binding motif},
		col title lines = 2,
		xlabel height = 5pt,
		empty facet warning = false,
		tight,
	]{modisco_TFs_neg_light}
	
		\pggf_logo(letter colors = DNA)
	
	\end{pggfplot}
	
	
	%%% subfigure labels %%%
	\subfiglabel{modisco}
	\subfiglabel{TFs}

\end{tikzpicture}%
			\caption{%
				\captiontitle{TF-MoDISco patterns match known transcription factor binding motifs}%
				\subfig{A} DNA Patterns identified by TF-MoDISco that negatively contribute to plantGREP predictions of enhancer strength in tobacco leaves in the light.\nextentry
				\subfig{B} Transcription factor binding motifs that match the patterns in \parensubfig{A}.
			}%
			\label{sfig:TF-MoDISco_neg_light}%
		\end{sfig}
		
		\begin{sfig}
			\input{SuppFig14}%
			\caption{%
				\captiontitle{TF-MoDISco patterns match known transcription factor binding motifs}%
				\subfig{A} DNA Patterns identified by TF-MoDISco that positively contribute to plantGREP predictions of enhancer strength in tobacco leaves in the dark.\nextentry
				\subfig{B} Transcription factor binding motifs that match the patterns in \parensubfig{A}.
			}%
			\label{sfig:TF-MoDISco_pos_dark}%
		\end{sfig}
		
		\begin{sfig}
			\begin{tikzpicture}
	\def\NA{\relax}
	\tikzset{
		TF name/.style = {
			/pgfplots/extra description/.append code = {
				\node[anchor = west] at (1, .5) {\usename{#1}};
			}
		}
	}

	%%% TF-MoDISco patterns %%%
	\coordinate (modisco) at (0, 0);
	
	\planticon[height = .5cm]{dark} at ($(modisco) + (.6, -.25)$);
	
	% get number of bases in motif
	\pgfplotstableread[col sep = tab]{rawData/modisco_motifs_neg_dark_axes.tsv}\axistable
	
	\pgfplotstablegetelem{0}{xmax}\of\axistable
	\pgfmathsetmacro{\nbases}{int(\pgfplotsretval)}

	\begin{pggfplot}[%
		at = {(modisco)},
		width = \nbases * 2.6mm,
		height = 6.6mm,
		xlabel = {},
		ylabel = {average DeepLIFT contribution score},
		y decimals = 2,
		enlarge x limits = 0,
		enlarge y limits = .1,
		xmajorticks = false,
		row titles from table = false,
		col title = {\raisebox{.5\baselineskip}{TF-MoDISco pattern}},
		col title lines = 2,
		xlabel height = 5pt,
		tight,
	]{modisco_motifs_neg_dark}
	
		\pggf_annotate(
			zeroline,
			annotation = {
				\node[anchor = south west, node font = \figsmall] at (current axis.south west) {pattern \pggfrowname};
			},
			skip facets = {6, 7, 15, 17}
		)
		
		\pggf_annotate(
			zeroline,
			annotation = {
				\node[anchor = north west, node font = \figsmall] at (current axis.north west) {pattern \pggfrowname};
			},
			only facets = {6}
		)
		
		\pggf_annotate(
			zeroline,
			annotation = {
				\node[anchor = north east, node font = \figsmall] at (current axis.north east) {pattern \pggfrowname};
			},
			only facets = {7, 15, 17}
		)
	
		\pggf_logo(letter colors = DNA)
	
	\end{pggfplot}
	
	
	%%% matching TF motifs %%%
	\coordinate[xshift = \columnsep] (TFs) at (modisco -| pggfplot c1r1.east);
	
	% get number of bases in motif
	\pgfplotstableread[col sep = tab]{rawData/modisco_TFs_neg_dark_axes.tsv}\axistable
	
	\pgfplotstablegetelem{0}{xmax}\of\axistable
	\pgfmathsetmacro{\nbases}{int(\pgfplotsretval)}
	
	\begin{pggfplot}[%
		at = {(TFs)},
		width = \nbases * 2.6mm,
		height = 6.6mm,
		xlabel = {},
		ylabel = {information content},
		ymin = 0,
		ymax = 2,
		ytick = {0, 1, 2},
		enlarge x limits = 0,
		enlarge y limits = {lower = 0, upper = .15},
		xmajorticks = false,
		row titles from table = false,
		col title = {matched transcription\\factor binding motif},
		col title lines = 2,
		xlabel height = 5pt,
		empty facet warning = false,
		tight,
	]{modisco_TFs_neg_dark}
	
		\pggf_annotate(
			annotation = {
				\node[align = center, node font = \figtiny, gray] at (current axis) {no matching transcription\\factor binding motif found};
			},
			only facets = {6, 7, 10, 13, 14, 17, 18, 19}
		)
	
		\pggf_logo(letter colors = DNA)
	
	\end{pggfplot}
	
	
	%%% subfigure labels %%%
	\subfiglabel{modisco}
	\subfiglabel{TFs}

\end{tikzpicture}%
			\caption{%
				\captiontitle{TF-MoDISco patterns match known transcription factor binding motifs}%
				\subfig{A} DNA Patterns identified by TF-MoDISco that negatively contribute to plantGREP predictions of enhancer strength in tobacco leaves in the dark.\nextentry
				\subfig{B} Transcription factor binding motifs that match the patterns in \parensubfig{A}.
			}%
			\label{sfig:TF-MoDISco_neg_dark}%
		\end{sfig}

		\begin{sfig}
			\input{SuppFig16}%
			\caption{%
				\captiontitle{TF-MoDISco patterns match known transcription factor binding motifs}%
				\subfig{A} DNA Patterns identified by TF-MoDISco that positively contribute to plantGREP predictions of enhancer strength in tobacco leaves at elevated ambient temperature.\nextentry
				\subfig{B} Transcription factor binding motifs that match the patterns in \parensubfig{A}.
			}%
			\label{sfig:TF-MoDISco_pos_warm}%
		\end{sfig}
		
		\begin{sfig}
			\begin{tikzpicture}
	\def\NA{\relax}
	\tikzset{
		TF name/.style = {
			/pgfplots/extra description/.append code = {
				\node[anchor = west] at (1, .5) {\usename{#1}};
			}
		}
	}

	%%% TF-MoDISco patterns %%%
	\coordinate (modisco) at (0, 0);
	
	\planticon[height = .5cm]{warm} at ($(modisco) + (.6, -.25)$);
	
	% get number of bases in motif
	\pgfplotstableread[col sep = tab]{rawData/modisco_motifs_neg_warm_axes.tsv}\axistable
	
	\pgfplotstablegetelem{0}{xmax}\of\axistable
	\pgfmathsetmacro{\nbases}{int(\pgfplotsretval)}

	\begin{pggfplot}[%
		at = {(modisco)},
		width = \nbases * 2.6mm,
		height = 6.6mm,
		xlabel = {},
		ylabel = {average DeepLIFT contribution score},
		y decimals = 2,
		enlarge x limits = 0,
		enlarge y limits = .1,
		xmajorticks = false,
		row titles from table = false,
		col title = {\raisebox{.5\baselineskip}{TF-MoDISco pattern}},
		col title lines = 2,
		xlabel height = 5pt,
		tight,
	]{modisco_motifs_neg_warm}
	
		\pggf_annotate(
			zeroline,
			annotation = {
				\node[anchor = south west, node font = \figsmall] at (current axis.south west) {pattern \pggfrowname};
			},
			skip facet = 8
		)
	
		\pggf_annotate(
			zeroline,
			annotation = {
				\node[anchor = south east, node font = \figsmall] at (current axis.south east) {pattern \pggfrowname};
			},
			only facet = 8
		)
	
		\pggf_logo(letter colors = DNA)
	
	\end{pggfplot}
	
	
	%%% matching TF motifs %%%
	\coordinate[xshift = \columnsep] (TFs) at (modisco -| pggfplot c1r1.east);
	
	% get number of bases in motif
	\pgfplotstableread[col sep = tab]{rawData/modisco_TFs_neg_warm_axes.tsv}\axistable
	
	\pgfplotstablegetelem{0}{xmax}\of\axistable
	\pgfmathsetmacro{\nbases}{int(\pgfplotsretval)}
	
	\begin{pggfplot}[%
		at = {(TFs)},
		width = \nbases * 2.6mm,
		height = 6.6mm,
		xlabel = {},
		ylabel = {information content},
		ymin = 0,
		ymax = 2,
		ytick = {0, 1, 2},
		enlarge x limits = 0,
		enlarge y limits = {lower = 0, upper = .15},
		xmajorticks = false,
		row titles from table = false,
		col title = {matched transcription\\factor binding motif},
		col title lines = 2,
		xlabel height = 5pt,
		empty facet warning = false,
		tight,
	]{modisco_TFs_neg_warm}
	
		\pggf_annotate(
			annotation = {
				\node[align = center, node font = \figtiny, gray] at (current axis) {no matching transcription\\factor binding motif found};
			},
			only facets = {5, 7}
		)
	
		\pggf_logo(letter colors = DNA)
	
	\end{pggfplot}
	
	
	%%% subfigure labels %%%
	\subfiglabel{modisco}
	\subfiglabel{TFs}

\end{tikzpicture}%
			\caption{%
				\captiontitle{TF-MoDISco patterns match known transcription factor binding motifs}%
				\subfig{A} DNA Patterns identified by TF-MoDISco that negatively contribute to plantGREP predictions of enhancer strength in tobacco leaves at elevated ambient temperature.\nextentry
				\subfig{B} Transcription factor binding motifs that match the patterns in \parensubfig{A}.
			}%
			\label{sfig:TF-MoDISco_neg_warm}%
		\end{sfig}
		
		\begin{sfig}
			\input{SuppFig18}%
			\caption{%
				\captiontitle{TF-MoDISco patterns match known transcription factor binding motifs}%
				\subfig{A} DNA Patterns identified by TF-MoDISco that positively contribute to plantGREP predictions of enhancer strength in tobacco leaves at reduced ambient temperature.\nextentry
				\subfig{B} Transcription factor binding motifs that match the patterns in \parensubfig{A}.
			}%
			\label{sfig:TF-MoDISco_pos_cold}%
		\end{sfig}
		
		\begin{sfig}
			\input{SuppFig19}%
			\caption{%
				\captiontitle{TF-MoDISco patterns match known transcription factor binding motifs}%
				\subfig{A} DNA Patterns identified by TF-MoDISco that negatively contribute to plantGREP predictions of enhancer strength in tobacco leaves at reduced ambient temperature.\nextentry
				\subfig{B} Transcription factor binding motifs that match the patterns in \parensubfig{A}.
			}%
			\label{sfig:TF-MoDISco_neg_cold}%
		\end{sfig}
		
		\begin{sfig}
			\input{SuppFig20}%
			\caption{%
				\captiontitle{TF-MoDISco patterns match known transcription factor binding motifs}%
				\subfig{A} DNA Patterns identified by TF-MoDISco that positively contribute to plantGREP predictions of enhancer strength in maize protoplasts.\nextentry
				\subfig{B} Transcription factor binding motifs that match the patterns in \parensubfig{A}.
			}%
			\label{sfig:TF-MoDISco_pos_maize}%
		\end{sfig}
		
		\begin{sfig}
			\input{SuppFig21}%
			\caption{%
				\captiontitle{TF-MoDISco patterns match known transcription factor binding motifs}%
				\subfig{A} DNA Patterns identified by TF-MoDISco that negatively contribute to plantGREP predictions of enhancer strength in maize protoplasts.\nextentry
				\subfig{B} Transcription factor binding motifs that match the patterns in \parensubfig{A}.
			}%
			\label{sfig:TF-MoDISco_neg_maize}%
		\end{sfig}
		
		\begin{sfig}
			\begin{tikzpicture}

	\pggfset{
		ylabel add to width = 2pt
	}

	%%% plantGREP ensemble accuracy at defined ensemble sizes %%%
	\coordinate (raw) at (0, 0);
	
	\begin{pggfplot}[%
		at = {(raw)},
		total width = \twothirdcolumnwidth,
		xlabel = ensemble size (number of models),
		ylabel = prediction accuracy ($R^2$),
		ytick = {.4, .45, ..., .8},
		y decimals = 2,
		title is csname,
		title color from name,
		enlarge x limits = 0,
		enlarge y limits = {lower = 0.1, upper = 0.1},
	]{plantGREP_ensemble_raw}
	
		\pggf_boxplot(
			color from col name,
			draw = black,
			cld = {anchor = south, text depth = 0pt}{max}
		)
	
	\end{pggfplot}
	
	
	%%% plantGREP ensemble accuracy summary %%%
	\coordinate (summary) at (raw -| \textwidth - \threecolumnwidth, 0);
	
	% get ymin and ymax from last plot
	\pgfplotstableread[col sep = tab]{rawData/plantGREP_ensemble_raw_axes.tsv}\axistable
	
	\pgfplotstablegetelem{0}{ymin}\of\axistable
	\edef\ymin{\expandonce{\pgfplotsretval}}
	
	\pgfplotstablegetelem{0}{ymax}\of\axistable
	\edef\ymax{\expandonce{\pgfplotsretval}}
	
	\begin{pggfplot}[%
		anchor = left of north west,
		at = {(summary |- pggfplot c1r1.north)},
		total width = \threecolumnwidth,
		xlabel = ensemble size (number of models),
		ylabel = prediction accuracy ($R^2$),
		ymin = \ymin,
		ymax = \ymax,
		xtick = {1, 10, 20, ..., 50},
		ytick = {.4, .45, ..., .8},
		y decimals = 2,
		line styles = {light,dark,warm,cold,maize},
		enlarge x limits = {lower = 0},
		enlarge y limits = {lower = 0.1, upper = 0.1},
		legend columns = 5,
		legend style = {anchor = south east, at = {(1, 1)}},
	]{plantGREP_ensemble_summary}
	
		\pggf_line(
			error = ribbon,
			error ribbon style = {fill opacity = .5},
			group column = condition,
			thick
		)
	
		\pggf_annotate(line legend)
	
	\end{pggfplot}
	
	
	%%% subfigure labels %%%
	\subfiglabel{raw}
	\subfiglabel{summary}

\end{tikzpicture}%
			\caption{%
				\captiontitle{A model ensemble increases the prediction accuracy of plantGREP}%
				\subfigtwo{A}{B} To build a plantGREP model ensemble, 50 models were trained with different, random initializations of the starting model weights. Box plots (as defined in \cref{sfig:TFBS_shuffle_strength}) in \parensubfig{A} show the prediction accuracy of a plantGREP model ensemble consisting of 1, 5, 10, 20, or 50 randomly sampled models. Lines in \parensubfig{B} indicate the mean prediction accuracy and the shaded area indicates the mean $\pm$ standard deviation. For each ensemble size, models were sampled independently for a total of 50 times.\nextentry
				Letters above the plots in \parensubfig{A} indicate significance groups determined by post-hoc Tukey tests performed separately for each condition and species. Exact $p$ values are listed in Supplementary Data 4.% TODO: make sure reference is correct
			}%
			\label{sfig:plantGREP_ensemble}%
		\end{sfig}

		\begin{sfig}
			\begin{tikzpicture}

	%%% in silico evolution with multiple mutations per round %%%
	\coordinate (muts per round) at (0, 0);
	
	% group styles
	\pgfplotsset{
		mpr1/.code = {
			\def\tmpA{dark}
			\ifx\tmpA\pggfcolname\relax
				\pgfkeysalso{fill = \pggfcolname!33!white}
			\else
				\pgfkeysalso{fill = \pggfcolname!50!white}
			\fi
		},
		mpr2/.code = {
			\def\tmpA{dark}
			\ifx\tmpA\pggfcolname\relax
				\pgfkeysalso{fill = \pggfcolname!67!white}
			\else
				\pgfkeysalso{fill = \pggfcolname}
			\fi
		},
		mpr3/.code = {
			\def\tmpA{dark}
			\ifx\tmpA\pggfcolname\relax
				\pgfkeysalso{fill = \pggfcolname}
			\else
				\pgfkeysalso{fill = \pggfcolname!50!black}
			\fi
		}
	}
	
	\begin{pggfplot}[%
		at = {(muts per round)},
		yshift = -\baselineskip,
		total width = \textwidth,
		xlabel = total number of mutations introduced during \textit{in silico} evolution,
		ylabel = \enrichment,
		xtick = {0, 0.76667, 1.86667},
		xticklabels = {0, 6, 12},
		ytick = {-12, -10, ..., 12},
		enlarge x limits = {0, abs lower = 1/6, abs upper = .5},
		enlarge y limits = {lower = 0.1, upper = 0.25},
		title is csname,
		title color from name,
		legend style = {anchor = south east, at = {(1, 1)}, yshift = 2\baselineskip},
		legend columns = 4,
		legend image width = .75\baselineskip,
		every legend image post/.append style = {fill opacity = 0.5},
		tight y
	]{evolution_muts_per_round}
	
		\pggf_annotate(zeroline)
		
		\pggf_annotate(
			manual legend = {
				{\textbf{mutations per round:}\hspace*{-\groupplotsep}} = {empty legend},
				1 = {violin legend = y, fill = gray!25!white},
				2 = {violin legend = y, fill = gray},
				3 = {violin legend = y, fill = gray!25!black}
			},
			only facet = 5
		)
	
		\pggf_violin(
			style from column = {muts_per_round},
			forget plot
		)
		
		\pggf_signif()
	
	\end{pggfplot}

\end{tikzpicture}%
			\caption{%
				\captiontitle{Increasing the number of mutations per round during \textit{in silico} evolution did not increase performance}%
				\textit{In silico} evolution was performed as in \cref{fig:synEnh}\subfigunformatted{A} (\textit{i.e.}, with one mutation per round) or by inserting two or three mutations within an 8-bp window in each round. Sequences before (0 mutations) or after inserting six or twelve mutations during \textit{in silico} evolution were subjected to Plant STARR-seq in the indicated conditions and species. Violin plots are as defined in \cref{sfig:accessibility}. Significant differences between sequences generated with one, two, or three mutations per round were determined by two-sided Wilcoxon rank-sum tests. The obtained $p$ values were Bonferroni-adjusted separately for each condition and species and are indicated: \printsigniflevels. Exact $p$ values are listed in Supplementary Data 4.% TODO: make sure reference is correct
			}%
			\label{sfig:evolution_muts_per_round}%
		\end{sfig}
		
		\begin{sfig}
			\begin{tikzpicture}

	%%% in silico evolution: species-/condition-specificity %%%
	\coordinate (evo specificity) at (0, 0);
	
	\colorlet{high}{OkabeIto4}
	\colorlet{any}{gray}
	\colorlet{low}{OkabeIto2}
	
	\begin{pggfplot}[%
		at = {(evo specificity)},
		yshift = -\baselineskip,
		total width = \textwidth,
		xlabel = \textit{in silico} evolution objective,
		ylabel = \enrichment,
		enlarge x limits = 0,
		enlarge y limits = {lower = 0.1, upper = 0.1},
		title is csname,
		title color from name,
		xytick = {-12, -10, ..., 12},
		xticklabel style = {rotate = 45, anchor = north east, inner sep = 0.236em},
		legend style = {anchor = south east, at = {(1, 1)}, yshift = 2\baselineskip},
		legend columns = 4,
		legend image width = .75\baselineskip,
		every legend image post/.append style = {fill opacity = 0.5},
		tight y
	]{evolution_objectives}
	
		\pggf_annotate(zeroline)
		
		\pggf_annotate(
			manual legend = {
				{\textbf{expected activity:}\hspace*{-\groupplotsep}} = {empty legend},
				high = {violin legend = y, fill = high},
				any = {violin legend = y, fill = any},
				low = {violin legend = y, fill = low}
			},
			only facet = 5
		)
	
		\pggf_violin(
			style from column = {fill = target},
			cld = {anchor = south, text depth = 0pt}{max},
			forget plot
		)
	
	\end{pggfplot}

\end{tikzpicture}%
			\caption{
				\captiontitle{\textit{In silico} evolved enhancers show species- and condition-specific activity}%
				Violin plots (as defined in \cref{sfig:accessibility} and colored according to the expected activity in the indicated condition) of the strength of enhancers before (no evolution) or after twelve rounds of \textit{in silico} evolution using plantGREP with different objectives to generate constitutive, species-, or condition-specific enhancers. Letters above the plots indicate significance groups determined by post-hoc Tukey tests performed separately for each condition and species. Exact $p$ values are listed in Supplementary Data 4.% TODO: make sure reference is correct
			}%
			\label{sfig:evolution_strength}
		\end{sfig}
	\fi
	%%% Supplementary figures end
	
	%%% Supplemenary tables start
	\ifstab
	
		\begin{stab}
			\caption{\captiontitle{Composition of the ACR library}}%
			\label{stab:tamsACR_composition}%
			\rowcolors{2}{white}{gray!10}

\pgfset{/pgf/number format/fixed}

\begin{tabular}{rrrrrr}
	\toprule
	\textbf{category} & \textbf{\usename{Arabidopsis}} & \textbf{\usename{tomato}} & \textbf{\usename{maize}} & \textbf{\usename{sorghum}} & \textbf{total} \\
	\midrule
	ACR sequences & \pgfmathprintnumber{58068} & \pgfmathprintnumber{24166} & \pgfmathprintnumber{52282} & \pgfmathprintnumber{35314} & \pgfmathprintnumber{169830} \\
	non-ACR sequences & \pgfmathprintnumber{1673} & \pgfmathprintnumber{2164} & \pgfmathprintnumber{1561} & \pgfmathprintnumber{1066} & \pgfmathprintnumber{6464} \\
	genes of interest &  & \pgfmathprintnumber{15670} &  &  & \pgfmathprintnumber{15670} \\
	known regulatory elements &  &  &  &  & \pgfmathprintnumber{36} \\
	 &  &  &  &  & \textbf{\pgfmathprintnumber[assume math mode]{192000}} \\
	\bottomrule
\end{tabular}%
		\end{stab}
		
		\begin{longstab}
			\rowcolors{2}{gray!10}{white}

\pgfplotstabletypeset[
	col sep = tab,
	string type,
	begin table = {\begin{longtable}[l]},
	end table = {\end{longtable}},
	every head row/.append style = {
		output empty row,
		before row = {%
			\hiderowcolors
			\caption{\captiontitle{Tomato genes of interest with sequences tiled across their upstream regions}}\label{stab:tomato_ROIs}\\%
			\showrowcolors
			\toprule
			\rowcolor{white}\textbf{gene ID} & \textbf{description} & \textbf{\# sequences} \\
			\midrule
			\endfirsthead
			\toprule
			\rowcolor{white}\textbf{gene ID} & \textbf{description} & \textbf{\# sequences} \\
			\midrule
			\endhead
			\bottomrule
			\endfoot
		},
	},
	columns/gene ID/.style = {
		column type = r
	},
	columns/description/.style = {
		column type = l
	},
	columns/tiles/.style = {
		column type = r
	}
]{../data/extra_files/Tomato_ROIs.tsv}%
		\end{longstab}
		
		\begin{stab}
			\caption{\captiontitle{Regulatory activity of test sequences}}%
			\label{stab:tamsACR_categories}%
			\rowcolors{2}{white}{gray!10}

\begin{tabular}{rrrrrrrr}
	\toprule
	\textbf{category}\textsuperscript{a} & \textbf{\usename{light}} & \textbf{\usename{dark}} & \textbf{\usename{warm}} & \textbf{\usename{cold}} & \textbf{\usename{maize}} & \textbf{at least one condition or assay system} & \textbf{all conditions and assay systems} \\
	\midrule
	enhancing sequences & 13.6\% & 19.3\% & 22.1\% & 11.1\% & 4.7\% & 33.0\% & 0.9\% \\
	inactive sequences & 75.5\% & 69.2\% & 76.1\% & 77.9\% & 84.3\% & 98.7\% & 42.9\% \\
	repressive sequences & 10.9\% & 11.5\% & 1.8\% & 11.0\% & 11\% & 26.7\% & 0.3\% \\
	\bottomrule
\end{tabular}

\smallskip

{\figsmall\textsuperscript{a} Categories: enhancing sequences, $\log_2$(enhancer strength) > 1; inactive sequences, $\log_2$(enhancer strength) between \textminus2 and 1; repressive sequences, $\log_2$(enhancer strength) < \textminus2}%
		\end{stab}
		
		\begin{longstab}
			% externalize TF motif logo plots
\setbox1=\hbox{
	\pgfplotstableforeachcolumnelement{id}\of\TFtable\as\TFid{%
		\tikzsetnextfilename{TF_logos/\TFid}%
		\begin{tikzpicture}
			\begin{pggfplot}[%
				width = 15cm,
				height = 1cm,
				enlargelimits = 0,
				enlarge y limits = {upper = 0.25},
				axis lines = none,
				axis background/.style = {fill = none},
				tight
			]{TF_logos/\TFid}
				\pggf_logo(letter colors = DNA)
			\end{pggfplot}
		\end{tikzpicture}%
		\tikzsetnextfilename{TF_logos/\TFid_RC}%
		\begin{tikzpicture}
			\begin{pggfplot}[%
				width = 15cm,
				height = 1cm,
				enlargelimits = 0,
				enlarge y limits = {upper = 0.25},
				axis lines = none,
				axis background/.style = {fill = none},
				tight
			]{TF_logos/\TFid_RC}
				\pggf_logo(letter colors = DNA)
			\end{pggfplot}
		\end{tikzpicture}%
	}%
}%

% set up macros to parse table content
\makeatletter
	\def\motifid TF_#1{#1}
	\def\TFfamilies #1 (#2){\def\slash{, }#1}
	\def\TFlogo#1{\tikzifexternalizing{\relax}{\pgfimage[width = \TX@col@width, height = 1.75\baselineskip]{extFigures/TF_logos/#1}}}
\makeatother

% typeset table
\rowcolors{2}{gray!10}{white}

\pgfplotstabletypeset[
	col sep = tab,
	string type,
	begin table = {\begin{xltabular}{\textwidth}},
	end table = {\end{xltabular}},
	every head row/.append style = {
		output empty row,
		before row = {%
			\hiderowcolors
			\caption{\captiontitle{Consensus binding motifs for transcription factors of the indicated families} Consensus motifs were obtained from ref. \textsuperscript{39}}\label{stab:TF_logos}\\% TODO: make sure reference is correct
			\showrowcolors
			\toprule
			\rowcolor{white}\textbf{motif ID} & \textbf{motif logo (forward)} & \textbf{motif logo (reverse-complement)} & \textbf{transcription factor families} \\
			\midrule
			\endfirsthead
			\toprule
			\rowcolor{white}\textbf{motif ID} & \textbf{motif logo (forward)} & \textbf{motif logo (reverse-complement)} & \textbf{transcription factor families} \\
			\midrule
			\endhead
			\bottomrule
			\endfoot
		},
	},
	create on use/logo/.style = {create col/copy = {id}},
	create on use/RClogo/.style = {create col/copy = {id}},
	columns = {id, logo, RClogo, family},
	columns/id/.style = {
		column type = r,
  	postproc cell content/.style = {
  		@cell content = \motifid##1
  	}
	},
	columns/logo/.style={
		column type = X,
  	postproc cell content/.style = {
  		@cell content = \TFlogo{##1}
  	}
	},
	columns/RClogo/.style={
		column type = X,
  	postproc cell content/.style = {
  		@cell content = \TFlogo{##1_RC}
  	}
	},
	columns/family/.style = {
		column type = l,
  	postproc cell content/.style = {
  		@cell content = \TFfamilies##1
  	}
	},
]{\TFtable}% TODO: make sure reference is correct
		\end{longstab}
	
		\begin{longstab}
			\rowcolors{2}{gray!10}{white}

\pgfplotstabletypeset[
	col sep = tab,
	string type,
	begin table = {\begin{xltabular}{\textwidth}},
	end table = {\end{xltabular}},
	every head row/.append style = {
		output empty row,
		before row = {%
			\hiderowcolors
			\caption{\captiontitle{Oligonucleotides used in this study}}\label{stab:primer}\\%
			\showrowcolors
			\toprule
			\rowcolor{white}\textbf{sequence} & \textbf{purpose} \\
			\midrule
			\endfirsthead
			\toprule
			\rowcolor{white}\textbf{sequence} & \textbf{purpose} \\
			\midrule
			\endhead
			\bottomrule
			\endfoot
		},
	},
	columns/sequence/.style = {
		column type = X,
		postproc cell content/.style = {
			@cell content = \courier\seqsplit{##1}
		}
	},
	columns/purpose/.style = {
		column type = {>{\raggedright\arraybackslash}m{.3\textwidth}}
	}
]{../data/extra_files/primer.tsv}%
		\end{longstab}
		
	\fi
	%%% Supplemenary tables end

\end{document}